% Options for packages loaded elsewhere
% Options for packages loaded elsewhere
\PassOptionsToPackage{unicode}{hyperref}
\PassOptionsToPackage{hyphens}{url}
%
\documentclass[
  10pt,
  ignorenonframetext,
  aspectratio=169,
  spanish,
]{beamer}
\newif\ifbibliography
\usepackage{pgfpages}
\setbeamertemplate{caption}[numbered]
\setbeamertemplate{caption label separator}{: }
\setbeamercolor{caption name}{fg=normal text.fg}
\beamertemplatenavigationsymbolsempty
% remove section numbering
\setbeamertemplate{part page}{
  \centering
  \begin{beamercolorbox}[sep=16pt,center]{part title}
    \usebeamerfont{part title}\insertpart\par
  \end{beamercolorbox}
}
\setbeamertemplate{section page}{
  \centering
  \begin{beamercolorbox}[sep=12pt,center]{section title}
    \usebeamerfont{section title}\insertsection\par
  \end{beamercolorbox}
}
\setbeamertemplate{subsection page}{
  \centering
  \begin{beamercolorbox}[sep=8pt,center]{subsection title}
    \usebeamerfont{subsection title}\insertsubsection\par
  \end{beamercolorbox}
}
% Prevent slide breaks in the middle of a paragraph
\widowpenalties 1 10000
\raggedbottom
\AtBeginPart{
  \frame{\partpage}
}
\AtBeginSection{
  \ifbibliography
  \else
    \frame{\sectionpage}
  \fi
}
\AtBeginSubsection{
  \frame{\subsectionpage}
}
\usepackage{iftex}
\ifPDFTeX
  \usepackage[T1]{fontenc}
  \usepackage[utf8]{inputenc}
  \usepackage{textcomp} % provide euro and other symbols
\else % if luatex or xetex
  \usepackage{unicode-math} % this also loads fontspec
  \defaultfontfeatures{Scale=MatchLowercase}
  \defaultfontfeatures[\rmfamily]{Ligatures=TeX,Scale=1}
\fi
\usepackage{lmodern}

\usetheme[]{Madrid}
\usecolortheme[]{dolphin}
\usefonttheme[]{professionalfonts}
\usefonttheme{serif} % use mainfont rather than sansfont for slide text
\ifPDFTeX\else
  % xetex/luatex font selection
  \setmainfont[]{DejaVu Sans}
  \setsansfont[]{DejaVu Sans}
  \setmonofont[]{DejaVu Sans Mono}
\fi
% Use upquote if available, for straight quotes in verbatim environments
\IfFileExists{upquote.sty}{\usepackage{upquote}}{}
\IfFileExists{microtype.sty}{% use microtype if available
  \usepackage[]{microtype}
  \UseMicrotypeSet[protrusion]{basicmath} % disable protrusion for tt fonts
}{}
\makeatletter
\@ifundefined{KOMAClassName}{% if non-KOMA class
  \IfFileExists{parskip.sty}{%
    \usepackage{parskip}
  }{% else
    \setlength{\parindent}{0pt}
    \setlength{\parskip}{6pt plus 2pt minus 1pt}}
}{% if KOMA class
  \KOMAoptions{parskip=half}}
\makeatother

\usepackage{color}
\usepackage{fancyvrb}
\newcommand{\VerbBar}{|}
\newcommand{\VERB}{\Verb[commandchars=\\\{\}]}
\DefineVerbatimEnvironment{Highlighting}{Verbatim}{commandchars=\\\{\}}
% Add ',fontsize=\small' for more characters per line
\usepackage{framed}
\definecolor{shadecolor}{RGB}{241,243,245}
\newenvironment{Shaded}{\begin{snugshade}}{\end{snugshade}}
\newcommand{\AlertTok}[1]{\textcolor[rgb]{0.68,0.00,0.00}{#1}}
\newcommand{\AnnotationTok}[1]{\textcolor[rgb]{0.37,0.37,0.37}{#1}}
\newcommand{\AttributeTok}[1]{\textcolor[rgb]{0.40,0.45,0.13}{#1}}
\newcommand{\BaseNTok}[1]{\textcolor[rgb]{0.68,0.00,0.00}{#1}}
\newcommand{\BuiltInTok}[1]{\textcolor[rgb]{0.00,0.23,0.31}{#1}}
\newcommand{\CharTok}[1]{\textcolor[rgb]{0.13,0.47,0.30}{#1}}
\newcommand{\CommentTok}[1]{\textcolor[rgb]{0.37,0.37,0.37}{#1}}
\newcommand{\CommentVarTok}[1]{\textcolor[rgb]{0.37,0.37,0.37}{\textit{#1}}}
\newcommand{\ConstantTok}[1]{\textcolor[rgb]{0.56,0.35,0.01}{#1}}
\newcommand{\ControlFlowTok}[1]{\textcolor[rgb]{0.00,0.23,0.31}{\textbf{#1}}}
\newcommand{\DataTypeTok}[1]{\textcolor[rgb]{0.68,0.00,0.00}{#1}}
\newcommand{\DecValTok}[1]{\textcolor[rgb]{0.68,0.00,0.00}{#1}}
\newcommand{\DocumentationTok}[1]{\textcolor[rgb]{0.37,0.37,0.37}{\textit{#1}}}
\newcommand{\ErrorTok}[1]{\textcolor[rgb]{0.68,0.00,0.00}{#1}}
\newcommand{\ExtensionTok}[1]{\textcolor[rgb]{0.00,0.23,0.31}{#1}}
\newcommand{\FloatTok}[1]{\textcolor[rgb]{0.68,0.00,0.00}{#1}}
\newcommand{\FunctionTok}[1]{\textcolor[rgb]{0.28,0.35,0.67}{#1}}
\newcommand{\ImportTok}[1]{\textcolor[rgb]{0.00,0.46,0.62}{#1}}
\newcommand{\InformationTok}[1]{\textcolor[rgb]{0.37,0.37,0.37}{#1}}
\newcommand{\KeywordTok}[1]{\textcolor[rgb]{0.00,0.23,0.31}{\textbf{#1}}}
\newcommand{\NormalTok}[1]{\textcolor[rgb]{0.00,0.23,0.31}{#1}}
\newcommand{\OperatorTok}[1]{\textcolor[rgb]{0.37,0.37,0.37}{#1}}
\newcommand{\OtherTok}[1]{\textcolor[rgb]{0.00,0.23,0.31}{#1}}
\newcommand{\PreprocessorTok}[1]{\textcolor[rgb]{0.68,0.00,0.00}{#1}}
\newcommand{\RegionMarkerTok}[1]{\textcolor[rgb]{0.00,0.23,0.31}{#1}}
\newcommand{\SpecialCharTok}[1]{\textcolor[rgb]{0.37,0.37,0.37}{#1}}
\newcommand{\SpecialStringTok}[1]{\textcolor[rgb]{0.13,0.47,0.30}{#1}}
\newcommand{\StringTok}[1]{\textcolor[rgb]{0.13,0.47,0.30}{#1}}
\newcommand{\VariableTok}[1]{\textcolor[rgb]{0.07,0.07,0.07}{#1}}
\newcommand{\VerbatimStringTok}[1]{\textcolor[rgb]{0.13,0.47,0.30}{#1}}
\newcommand{\WarningTok}[1]{\textcolor[rgb]{0.37,0.37,0.37}{\textit{#1}}}

\usepackage{longtable,booktabs,array}
\usepackage{calc} % for calculating minipage widths
\usepackage{caption}
% Make caption package work with longtable
\makeatletter
\def\fnum@table{\tablename~\thetable}
\makeatother
\usepackage{graphicx}
\makeatletter
\newsavebox\pandoc@box
\newcommand*\pandocbounded[1]{% scales image to fit in text height/width
  \sbox\pandoc@box{#1}%
  \Gscale@div\@tempa{\textheight}{\dimexpr\ht\pandoc@box+\dp\pandoc@box\relax}%
  \Gscale@div\@tempb{\linewidth}{\wd\pandoc@box}%
  \ifdim\@tempb\p@<\@tempa\p@\let\@tempa\@tempb\fi% select the smaller of both
  \ifdim\@tempa\p@<\p@\scalebox{\@tempa}{\usebox\pandoc@box}%
  \else\usebox{\pandoc@box}%
  \fi%
}
% Set default figure placement to htbp
\def\fps@figure{htbp}
\makeatother


% definitions for citeproc citations
\NewDocumentCommand\citeproctext{}{}
\NewDocumentCommand\citeproc{mm}{%
  \begingroup\def\citeproctext{#2}\cite{#1}\endgroup}
\makeatletter
 % allow citations to break across lines
 \let\@cite@ofmt\@firstofone
 % avoid brackets around text for \cite:
 \def\@biblabel#1{}
 \def\@cite#1#2{{#1\if@tempswa , #2\fi}}
\makeatother
\newlength{\cslhangindent}
\setlength{\cslhangindent}{1.5em}
\newlength{\csllabelwidth}
\setlength{\csllabelwidth}{3em}
\newenvironment{CSLReferences}[2] % #1 hanging-indent, #2 entry-spacing
 {\begin{list}{}{%
  \setlength{\itemindent}{0pt}
  \setlength{\leftmargin}{0pt}
  \setlength{\parsep}{0pt}
  % turn on hanging indent if param 1 is 1
  \ifodd #1
   \setlength{\leftmargin}{\cslhangindent}
   \setlength{\itemindent}{-1\cslhangindent}
  \fi
  % set entry spacing
  \setlength{\itemsep}{#2\baselineskip}}}
 {\end{list}}
\usepackage{calc}
\newcommand{\CSLBlock}[1]{\hfill\break\parbox[t]{\linewidth}{\strut\ignorespaces#1\strut}}
\newcommand{\CSLLeftMargin}[1]{\parbox[t]{\csllabelwidth}{\strut#1\strut}}
\newcommand{\CSLRightInline}[1]{\parbox[t]{\linewidth - \csllabelwidth}{\strut#1\strut}}
\newcommand{\CSLIndent}[1]{\hspace{\cslhangindent}#1}

\ifLuaTeX
\usepackage[bidi=basic]{babel}
\else
\usepackage[bidi=default]{babel}
\fi
\ifPDFTeX
\else
\babelfont{rm}[]{DejaVu Sans}
\fi
% get rid of language-specific shorthands (see #6817):
\let\LanguageShortHands\languageshorthands
\def\languageshorthands#1{}


\setlength{\emergencystretch}{3em} % prevent overfull lines

\providecommand{\tightlist}{%
  \setlength{\itemsep}{0pt}\setlength{\parskip}{0pt}}



 


% Estilo separado del contenido. Compilación comprobada con XeLaTeX.
\definecolor{SYSBAzul}{HTML}{0B3954}
\definecolor{SYSBVerde}{HTML}{168C91}
\definecolor{SYSBTexto}{HTML}{203E50}
\setbeamercolor{structure}{fg=SYSBAzul}
\setbeamercolor{normal text}{fg=SYSBTexto,bg=white}
\setbeamercolor{frametitle}{fg=white,bg=SYSBAzul}
\setbeamercolor{block title}{fg=white,bg=SYSBAzul}
\setbeamercolor{block body}{fg=SYSBTexto,bg=SYSBAzul!4}
\setbeamerfont{frametitle}{size=\large,series=\bfseries}
\setbeamerfont{footline}{size=\tiny}
\setbeamersize{text margin left=6mm,text margin right=6mm}
\setbeamertemplate{navigation symbols}{}
\setbeamertemplate{footline}{\hfill\textcolor{SYSBAzul}{SYSB\quad\insertframenumber/\inserttotalframenumber}\hspace{6mm}\vspace{2mm}}
\setbeameroption{hide notes}
\setlength{\parskip}{2pt}
\setlength{\abovedisplayskip}{5pt}
\setlength{\belowdisplayskip}{5pt}
\hypersetup{pdftitle={Sistemas y Señales Biomédicos},pdfauthor={Pablo Eduardo Caicedo-Rodríguez}}

\makeatletter
\@ifpackageloaded{tcolorbox}{}{\usepackage[skins,breakable]{tcolorbox}}
\@ifpackageloaded{fontawesome5}{}{\usepackage{fontawesome5}}
\definecolor{quarto-callout-color}{HTML}{909090}
\definecolor{quarto-callout-note-color}{HTML}{0758E5}
\definecolor{quarto-callout-important-color}{HTML}{CC1914}
\definecolor{quarto-callout-warning-color}{HTML}{EB9113}
\definecolor{quarto-callout-tip-color}{HTML}{00A047}
\definecolor{quarto-callout-caution-color}{HTML}{FC5300}
\definecolor{quarto-callout-color-frame}{HTML}{acacac}
\definecolor{quarto-callout-note-color-frame}{HTML}{4582ec}
\definecolor{quarto-callout-important-color-frame}{HTML}{d9534f}
\definecolor{quarto-callout-warning-color-frame}{HTML}{f0ad4e}
\definecolor{quarto-callout-tip-color-frame}{HTML}{02b875}
\definecolor{quarto-callout-caution-color-frame}{HTML}{fd7e14}
\makeatother
\makeatletter
\@ifpackageloaded{caption}{}{\usepackage{caption}}
\AtBeginDocument{%
\ifdefined\contentsname
  \renewcommand*\contentsname{Tabla de contenidos}
\else
  \newcommand\contentsname{Tabla de contenidos}
\fi
\ifdefined\listfigurename
  \renewcommand*\listfigurename{Listado de Figuras}
\else
  \newcommand\listfigurename{Listado de Figuras}
\fi
\ifdefined\listtablename
  \renewcommand*\listtablename{Listado de Tablas}
\else
  \newcommand\listtablename{Listado de Tablas}
\fi
\ifdefined\figurename
  \renewcommand*\figurename{Figura}
\else
  \newcommand\figurename{Figura}
\fi
\ifdefined\tablename
  \renewcommand*\tablename{Tabla}
\else
  \newcommand\tablename{Tabla}
\fi
}
\@ifpackageloaded{float}{}{\usepackage{float}}
\floatstyle{ruled}
\@ifundefined{c@chapter}{\newfloat{codelisting}{h}{lop}}{\newfloat{codelisting}{h}{lop}[chapter]}
\floatname{codelisting}{Listado}
\newcommand*\listoflistings{\listof{codelisting}{Listado de Listados}}
\makeatother
\makeatletter
\makeatother
\makeatletter
\@ifpackageloaded{caption}{}{\usepackage{caption}}
\@ifpackageloaded{subcaption}{}{\usepackage{subcaption}}
\makeatother

\usepackage{bookmark}
\IfFileExists{xurl.sty}{\usepackage{xurl}}{} % add URL line breaks if available
\urlstyle{same}
\hypersetup{
  pdftitle={Sistemas y Señales Biomédicos},
  pdfauthor={Pablo Eduardo Caicedo-Rodríguez},
  pdflang={es},
  hidelinks,
  pdfcreator={LaTeX via pandoc}}


\author{}
\date{}

\begin{document}


\begin{frame}{Sistemas y Señales Biomédicos}
\phantomsection\label{d2-001}
\textbf{Semanas 6--10 · Espectro y sistemas LTI}

Ph. D. Pablo Eduardo Caicedo-Rodríguez

Material de clase y estudio · Ingeniería Biomédica

\note{\textbf{Libro guía:} John Semmlow, \emph{Circuits, Signals, and
Systems for Bioengineers}, 3.ª ed., 2018.

§3.4, \textbf{Time--Frequency Transformation In The Discrete Domain}
(página 142 del visor PDF); §4.2, \textbf{Data Acquisition And Storage}
(página 175 del visor PDF); §4.6, \textbf{Time--Frequency Analysis}
(página 205 del visor PDF); §5.4, \textbf{Using The Impulse Response To
Find A Systems Output To Any Input---Convolution} (página 219 del visor
PDF).

\textbf{Correspondencia:} Síntesis o encuadre que reúne varias
secciones; no corresponde a un único apartado del libro. Lectura
principal: Semmlow (2018), tercera edición. La correspondencia usa la
edición del archivo aportado, que difiere de la segunda edición de 2012
citada en el microcurrículo.

\textbf{Lectura:} use estas secciones como ruta de preparación y
consulta. La bibliografía aparece al final del bloque.}
\end{frame}

\begin{frame}{Cómo estudiar las semanas 6--10}
\phantomsection\label{d2-002}
\note{\textbf{Libro guía:} John Semmlow, \emph{Circuits, Signals, and
Systems for Bioengineers}, 3.ª ed., 2018.

§3.4, \textbf{Time--Frequency Transformation In The Discrete Domain}
(página 142 del visor PDF); §4.2, \textbf{Data Acquisition And Storage}
(página 175 del visor PDF); §4.6, \textbf{Time--Frequency Analysis}
(página 205 del visor PDF); §5.4, \textbf{Using The Impulse Response To
Find A Systems Output To Any Input---Convolution} (página 219 del visor
PDF).

\textbf{Correspondencia:} Síntesis o encuadre que reúne varias
secciones; no corresponde a un único apartado del libro.

\textbf{Lectura:} use estas secciones como ruta de preparación y
consulta. La bibliografía aparece al final del bloque.}
\end{frame}

\begin{frame}{Del espectro ideal a decisiones reproducibles}
\phantomsection\label{d2-003}
Al terminar la semana 10, el estudiante debe poder:

\begin{enumerate}
\tightlist
\item
  predecir el efecto de operaciones temporales sobre la transformada de
  Fourier;
\item
  calcular e interpretar DFT y FFT con escalamiento explícito;
\item
  explicar resolución frecuencial, fuga y efecto de las ventanas;
\item
  relacionar adquisición, muestreo y almacenamiento con el espectro
  observado;
\item
  estimar PSD mediante periodograma y Welch;
\item
  analizar señales no estacionarias con STFT;
\item
  verificar linealidad, invariancia temporal y respuesta al impulso.
\end{enumerate}

\begin{block}{Idea conductora}

El espectro calculado depende tanto de la señal fisiológica como de la
adquisición, la duración observada y el estimador seleccionado.

\end{block}

\note{\textbf{Libro guía:} John Semmlow, \emph{Circuits, Signals, and
Systems for Bioengineers}, 3.ª ed., 2018.

§3.4, \textbf{Time--Frequency Transformation In The Discrete Domain}
(página 142 del visor PDF); §4.2, \textbf{Data Acquisition And Storage}
(página 175 del visor PDF); §4.6, \textbf{Time--Frequency Analysis}
(página 205 del visor PDF); §5.4, \textbf{Using The Impulse Response To
Find A Systems Output To Any Input---Convolution} (página 219 del visor
PDF).

\textbf{Correspondencia:} Síntesis o encuadre que reúne varias
secciones; no corresponde a un único apartado del libro.

\textbf{Lectura:} use estas secciones como ruta de preparación y
consulta. La bibliografía aparece al final del bloque.}
\end{frame}

\begin{frame}{Mapa de las diez sesiones}
\phantomsection\label{d2-004}
\begin{longtable}[]{@{}
  >{\raggedright\arraybackslash}p{(\linewidth - 4\tabcolsep) * \real{0.3333}}
  >{\raggedright\arraybackslash}p{(\linewidth - 4\tabcolsep) * \real{0.3333}}
  >{\raggedright\arraybackslash}p{(\linewidth - 4\tabcolsep) * \real{0.3333}}@{}}
\toprule\noalign{}
\begin{minipage}[b]{\linewidth}\raggedright
Semana
\end{minipage} & \begin{minipage}[b]{\linewidth}\raggedright
Sesión 1
\end{minipage} & \begin{minipage}[b]{\linewidth}\raggedright
Sesión 2
\end{minipage} \\
\midrule\noalign{}
\endhead
6 & Propiedades de Fourier & DFT y FFT \\
7 & Resolución, fuga y ventanas & Cadena de adquisición \\
8 & Muestreo, Nyquist y aliasing & Potencia, PSD y periodograma \\
9 & Promediado y Welch & No estacionariedad y STFT \\
10 & Ventanas tiempo-frecuencia & Sistemas LTI y respuesta al impulso \\
\bottomrule\noalign{}
\end{longtable}

Cada sesión integra fundamento matemático, interpretación biomédica,
ejemplo computacional y verificación conceptual.

\note{\textbf{Libro guía:} John Semmlow, \emph{Circuits, Signals, and
Systems for Bioengineers}, 3.ª ed., 2018.

§3.4, \textbf{Time--Frequency Transformation In The Discrete Domain}
(página 142 del visor PDF); §4.2, \textbf{Data Acquisition And Storage}
(página 175 del visor PDF); §4.6, \textbf{Time--Frequency Analysis}
(página 205 del visor PDF); §5.4, \textbf{Using The Impulse Response To
Find A Systems Output To Any Input---Convolution} (página 219 del visor
PDF).

\textbf{Correspondencia:} Síntesis o encuadre que reúne varias
secciones; no corresponde a un único apartado del libro.

\textbf{Lectura:} use estas secciones como ruta de preparación y
consulta. La bibliografía aparece al final del bloque.}
\end{frame}

\begin{frame}{Organización de cada sesión y trabajo independiente}
\phantomsection\label{d2-005}
\begin{longtable}[]{@{}llr@{}}
\toprule\noalign{}
Momento & Actividad & Tiempo \\
\midrule\noalign{}
\endhead
Apertura & Lectura previa y predicción & 10 min \\
Fundamento & Ecuación, supuestos y ejemplo & 30 min \\
Aplicación & Cálculo o código en parejas & 30 min \\
Cierre & Comprobación individual y límites & 20 min \\
\bottomrule\noalign{}
\end{longtable}

\begin{block}{Nota}

Prepare la lectura en inglés indicada en las notas. Use las ampliaciones
como consulta y trabajo independiente según los conocimientos previos
del grupo.

\end{block}

\note{\textbf{Libro guía:} John Semmlow, \emph{Circuits, Signals, and
Systems for Bioengineers}, 3.ª ed., 2018.

§1.1.1, \textbf{Goals Of This Book} (página 10 del visor PDF).

\textbf{Correspondencia:} Organización docente basada en las actividades
y RAE del microcurrículo. No corresponde a una sección temática del
libro. Microcurrículo: actividades de aprendizaje y RAE, PDF 2--3.
Distribución docente propuesta para las sesiones de 90 minutos.

\textbf{Lectura:} use estas secciones como ruta de preparación y
consulta. La bibliografía aparece al final del bloque.}
\end{frame}

\begin{frame}{Evidencias de aprendizaje y criterios de revisión}
\phantomsection\label{d2-006}
\begin{longtable}[]{@{}
  >{\raggedright\arraybackslash}p{(\linewidth - 2\tabcolsep) * \real{0.5000}}
  >{\raggedright\arraybackslash}p{(\linewidth - 2\tabcolsep) * \real{0.5000}}@{}}
\toprule\noalign{}
\begin{minipage}[b]{\linewidth}\raggedright
Evidencia
\end{minipage} & \begin{minipage}[b]{\linewidth}\raggedright
Criterio observable
\end{minipage} \\
\midrule\noalign{}
\endhead
Lectura y explicación oral & Identifica el supuesto del texto y lo
explica con sus palabras \\
Cálculo individual & Procedimiento correcto, unidades y caso límite \\
Cuaderno en parejas & Código reproducible, datos trazables y
comprobación numérica \\
Informe breve & Interpreta el resultado, cita las fuentes y declara
limitaciones \\
\bottomrule\noalign{}
\end{longtable}

\begin{block}{Concepto clave}

Cada entrega debe permitir distinguir aportes individuales y trabajo
conjunto. La corrección matemática y la justificación biomédica se
revisan por separado. Los porcentajes de evaluación se rigen por el
programa que confirme el profesor.

\end{block}

\note{\textbf{Libro guía:} John Semmlow, \emph{Circuits, Signals, and
Systems for Bioengineers}, 3.ª ed., 2018.

§1.1.1, \textbf{Goals Of This Book} (página 10 del visor PDF).

\textbf{Correspondencia:} Organización docente basada en las actividades
y RAE del microcurrículo. No corresponde a una sección temática del
libro. Microcurrículo: RAE y actividades de evaluación, PDF 2--3. No se
resuelven aquí las inconsistencias administrativas del documento.

\textbf{Lectura:} use estas secciones como ruta de preparación y
consulta. La bibliografía aparece al final del bloque.}
\end{frame}

\begin{frame}[fragile]{Convenciones de los ejemplos en Python}
\phantomsection\label{d2-007}
Los fragmentos que usan \texttt{x} y \texttt{fs} reciben una señal real,
unidimensional, finita y uniformemente muestreada. En la actividad,
documente unidades y metadatos antes de sustituirla.

\begin{Shaded}
\begin{Highlighting}[]
\ImportTok{import}\NormalTok{ numpy }\ImportTok{as}\NormalTok{ np}
\NormalTok{fs }\OperatorTok{=} \FloatTok{256.0}
\NormalTok{t }\OperatorTok{=}\NormalTok{ np.arange(}\DecValTok{2048}\NormalTok{)}\OperatorTok{/}\NormalTok{fs}
\NormalTok{x }\OperatorTok{=}\NormalTok{ np.cos(}\DecValTok{2}\OperatorTok{*}\NormalTok{np.pi}\OperatorTok{*}\DecValTok{10}\OperatorTok{*}\NormalTok{t) }\OperatorTok{+} \FloatTok{.3}\OperatorTok{*}\NormalTok{np.sin(}\DecValTok{2}\OperatorTok{*}\NormalTok{np.pi}\OperatorTok{*}\DecValTok{23}\OperatorTok{*}\NormalTok{t)}
\NormalTok{L, H, Nfft }\OperatorTok{=} \DecValTok{256}\NormalTok{, }\DecValTok{128}\NormalTok{, }\DecValTok{512}
\end{Highlighting}
\end{Shaded}

\begin{block}{Nota}

Este es un ejemplo sintético en unidades arbitrarias. Las variables
específicas de una actividad (etiquetas, sujetos, umbrales, referencias)
deben definirse para ese problema. Los casos reproducibles están en
\texttt{codigo/ejemplos\_verificados.py}.

\end{block}

\note{\textbf{Libro guía:} John Semmlow, \emph{Circuits, Signals, and
Systems for Bioengineers}, 3.ª ed., 2018.

§1.1.1, \textbf{Goals Of This Book} (página 10 del visor PDF).

\textbf{Correspondencia:} Ejemplo, figura o implementación de
elaboración docente a partir de estos fundamentos. No es una
transcripción del código MATLAB del libro. El lenguaje del
microcurrículo suministrado es Python. El libro guía utiliza MATLAB; se
conservan conceptos y se explicitan diferencias de convenciones.

\textbf{Lectura:} use estas secciones como ruta de preparación y
consulta. La bibliografía aparece al final del bloque.}
\end{frame}

\begin{frame}{Convenciones del bloque}
\phantomsection\label{d2-008}
Transformada continua en frecuencia ordinaria:

\[X(f)=\int_{-\infty}^{\infty}x(t)e^{-j2\pi ft}\,dt\]

DFT de \(N\) muestras:

\[X[k]=\sum_{n=0}^{N-1}x[n]e^{-j2\pi kn/N}\]

\begin{alertblock}{Atención}

Las constantes de normalización cambian entre textos y bibliotecas. Toda
interpretación cuantitativa debe declarar la convención utilizada.

\end{alertblock}

\note{\textbf{Libro guía:} John Semmlow, \emph{Circuits, Signals, and
Systems for Bioengineers}, 3.ª ed., 2018.

§3.4, \textbf{Time--Frequency Transformation In The Discrete Domain}
(página 142 del visor PDF); §4.2, \textbf{Data Acquisition And Storage}
(página 175 del visor PDF); §4.6, \textbf{Time--Frequency Analysis}
(página 205 del visor PDF); §5.4, \textbf{Using The Impulse Response To
Find A Systems Output To Any Input---Convolution} (página 219 del visor
PDF).

\textbf{Correspondencia:} Síntesis o encuadre que reúne varias
secciones; no corresponde a un único apartado del libro.

\textbf{Lectura:} use estas secciones como ruta de preparación y
consulta. La bibliografía aparece al final del bloque.}
\end{frame}

\begin{frame}{Semana 6 · Propiedades de Fourier y DFT}
\phantomsection\label{d2-009}
\note{\textbf{Libro guía:} John Semmlow, \emph{Circuits, Signals, and
Systems for Bioengineers}, 3.ª ed., 2018.

§3.4, \textbf{Time--Frequency Transformation In The Discrete Domain}
(página 142 del visor PDF); §4.2, \textbf{Data Acquisition And Storage}
(página 175 del visor PDF); §4.6, \textbf{Time--Frequency Analysis}
(página 205 del visor PDF); §5.4, \textbf{Using The Impulse Response To
Find A Systems Output To Any Input---Convolution} (página 219 del visor
PDF).

\textbf{Correspondencia:} Síntesis o encuadre que reúne varias
secciones; no corresponde a un único apartado del libro.

\textbf{Lectura:} use estas secciones como ruta de preparación y
consulta. La bibliografía aparece al final del bloque.}
\end{frame}

\begin{frame}{Sesión 11 · Las propiedades permiten razonar sin
reintegrar}
\phantomsection\label{d2-010}
Si conocemos el par

\[x(t)\longleftrightarrow X(f),\]

podemos anticipar la transformada de una señal modificada mediante
propiedades algebraicas.

\begin{block}{Propósito}

Las propiedades no reemplazan la interpretación. Permiten conectar una
operación en tiempo con su consecuencia en magnitud y fase.

\end{block}

\note{\textbf{Libro guía:} John Semmlow, \emph{Circuits, Signals, and
Systems for Bioengineers}, 3.ª ed., 2018.

§3.3.7, \textbf{The Continuous Fourier Transform} (página 138 del visor
PDF); §5.6, \textbf{Convolution In The Frequency Domain} (página 237 del
visor PDF).

\textbf{Correspondencia:} Las secciones aportan el fundamento. La
diapositiva amplía la formulación o precisa condiciones que no se
presentan con idéntica notación en el libro. Desarrollo de propiedades
con convención en Hz; las secciones son el fundamento y no una tabla
idéntica a la diapositiva.

\textbf{Microcurrículo:} semana 6, sesión 1 (sesión global 11). Los
numerales del cronograma son identificadores curriculares; no son los
apartados del libro citados arriba.}
\end{frame}

\begin{frame}{Linealidad conserva la combinación}
\phantomsection\label{d2-011}
Si

\[x_1(t)\longleftrightarrow X_1(f),\qquad x_2(t)\longleftrightarrow X_2(f),\]

entonces

\[a x_1(t)+b x_2(t)\longleftrightarrow aX_1(f)+bX_2(f).\]

\begin{block}{Nota}

En un ECG con deriva e interferencia de red, el espectro observado
combina las transformadas de cada componente solo si el modelo aditivo
resulta razonable.

\end{block}

\note{\textbf{Libro guía:} John Semmlow, \emph{Circuits, Signals, and
Systems for Bioengineers}, 3.ª ed., 2018.

§3.3.7, \textbf{The Continuous Fourier Transform} (página 138 del visor
PDF); §5.6, \textbf{Convolution In The Frequency Domain} (página 237 del
visor PDF).

\textbf{Correspondencia:} Las secciones aportan el fundamento. La
diapositiva amplía la formulación o precisa condiciones que no se
presentan con idéntica notación en el libro. Desarrollo de propiedades
con convención en Hz; las secciones son el fundamento y no una tabla
idéntica a la diapositiva.

\textbf{Microcurrículo:} semana 6, sesión 1 (sesión global 11). Los
numerales del cronograma son identificadores curriculares; no son los
apartados del libro citados arriba.}
\end{frame}

\begin{frame}{Linealidad no separa fuentes automáticamente}
\phantomsection\label{d2-012}
Modelo aditivo:

\[x_{obs}(t)=s(t)+i(t)+a(t).\]

\begin{itemize}
\tightlist
\item
  \(s(t)\): señal de interés.
\item
  \(i(t)\): interferencia estructurada.
\item
  \(a(t)\): artefacto.
\end{itemize}

Entonces \(X_{obs}(f)=S(f)+I(f)+A(f)\).

\begin{alertblock}{Atención}

La suma espectral puede producir cancelación por fase. La magnitud de la
suma no equivale a sumar magnitudes.

\end{alertblock}

\note{\textbf{Libro guía:} John Semmlow, \emph{Circuits, Signals, and
Systems for Bioengineers}, 3.ª ed., 2018.

§3.3.7, \textbf{The Continuous Fourier Transform} (página 138 del visor
PDF); §5.6, \textbf{Convolution In The Frequency Domain} (página 237 del
visor PDF).

\textbf{Correspondencia:} Las secciones aportan el fundamento. La
diapositiva amplía la formulación o precisa condiciones que no se
presentan con idéntica notación en el libro. Desarrollo de propiedades
con convención en Hz; las secciones son el fundamento y no una tabla
idéntica a la diapositiva.

\textbf{Microcurrículo:} semana 6, sesión 1 (sesión global 11). Los
numerales del cronograma son identificadores curriculares; no son los
apartados del libro citados arriba.}
\end{frame}

\begin{frame}{Un retraso añade fase lineal}
\phantomsection\label{d2-013}
\[x(t-t_0)\longleftrightarrow X(f)e^{-j2\pi ft_0}\]

Por tanto:

\[|X_{ret}(f)|=|X(f)|,\]

\[\angle X_{ret}(f)=\angle X(f)-2\pi ft_0.\]

\begin{block}{Concepto clave}

Un retraso puro conserva la magnitud. La información temporal aparece en
la pendiente de fase.

\end{block}

\note{\textbf{Libro guía:} John Semmlow, \emph{Circuits, Signals, and
Systems for Bioengineers}, 3.ª ed., 2018.

§3.3.7, \textbf{The Continuous Fourier Transform} (página 138 del visor
PDF); §5.6, \textbf{Convolution In The Frequency Domain} (página 237 del
visor PDF).

\textbf{Correspondencia:} Las secciones aportan el fundamento. La
diapositiva amplía la formulación o precisa condiciones que no se
presentan con idéntica notación en el libro. Desarrollo de propiedades
con convención en Hz; las secciones son el fundamento y no una tabla
idéntica a la diapositiva.

\textbf{Microcurrículo:} semana 6, sesión 1 (sesión global 11). Los
numerales del cronograma son identificadores curriculares; no son los
apartados del libro citados arriba.}
\end{frame}

\begin{frame}{La fase lineal representa un retardo constante}
\phantomsection\label{d2-014}
El retardo de grupo se define como

\[\tau_g(f)=-\frac{1}{2\pi}\frac{d\phi(f)}{df}.\]

Si \(\phi(f)=-2\pi ft_0\), entonces \(\tau_g(f)=t_0\).

\begin{block}{Nota}

Con retardo de grupo constante y ganancia constante sobre la banda de la
señal, la forma se conserva salvo escala y desplazamiento. La fase
lineal por sí sola no evita distorsión de amplitud.

\end{block}

\note{\textbf{Libro guía:} John Semmlow, \emph{Circuits, Signals, and
Systems for Bioengineers}, 3.ª ed., 2018.

§6.4.1, \textbf{The Spectrum Of A Transfer Function} (página 259 del
visor PDF); §8.4.1, \textbf{Filter Properties} (página 360 del visor
PDF).

\textbf{Correspondencia:} Las secciones aportan el fundamento. La
diapositiva amplía la formulación o precisa condiciones que no se
presentan con idéntica notación en el libro. Formalización del retardo
de grupo a partir de la fase. El criterio de no distorsión también exige
ganancia constante.

\textbf{Microcurrículo:} semana 6, sesión 1 (sesión global 11). Los
numerales del cronograma son identificadores curriculares; no son los
apartados del libro citados arriba.}
\end{frame}

\begin{frame}{Desplazar en frecuencia equivale a modular}
\phantomsection\label{d2-015}
\[x(t)e^{j2\pi f_0t}\longleftrightarrow X(f-f_0).\]

Para modulación real:

\[x(t)\cos(2\pi f_0t)
\longleftrightarrow \frac{1}{2}X(f-f_0)+\frac{1}{2}X(f+f_0).\]

\begin{block}{Comprobación}

La multiplicación por una portadora crea copias desplazadas del
espectro. Este principio aparece en instrumentación y comunicaciones
biomédicas.

\end{block}

\note{\textbf{Libro guía:} John Semmlow, \emph{Circuits, Signals, and
Systems for Bioengineers}, 3.ª ed., 2018.

§3.3.7, \textbf{The Continuous Fourier Transform} (página 138 del visor
PDF); §5.6, \textbf{Convolution In The Frequency Domain} (página 237 del
visor PDF).

\textbf{Correspondencia:} Las secciones aportan el fundamento. La
diapositiva amplía la formulación o precisa condiciones que no se
presentan con idéntica notación en el libro. Desarrollo de propiedades
con convención en Hz; las secciones son el fundamento y no una tabla
idéntica a la diapositiva.

\textbf{Microcurrículo:} semana 6, sesión 1 (sesión global 11). Los
numerales del cronograma son identificadores curriculares; no son los
apartados del libro citados arriba.}
\end{frame}

\begin{frame}{El escalamiento temporal modifica ancho y amplitud}
\phantomsection\label{d2-016}
\[x(at)\longleftrightarrow \frac{1}{|a|}X\left(\frac{f}{a}\right).\]

\begin{itemize}
\tightlist
\item
  \(|a|>1\): la señal se comprime en tiempo y su espectro se expande.
\item
  \(0<|a|<1\): la señal se expande en tiempo y su espectro se comprime.
\item
  \(a<0\): también ocurre inversión temporal.
\end{itemize}

\begin{alertblock}{Atención}

El factor \(1/|a|\) forma parte de la propiedad. Omitirlo cambia la
escala de la transformada.

\end{alertblock}

\note{\textbf{Libro guía:} John Semmlow, \emph{Circuits, Signals, and
Systems for Bioengineers}, 3.ª ed., 2018.

§3.3.7, \textbf{The Continuous Fourier Transform} (página 138 del visor
PDF); §5.6, \textbf{Convolution In The Frequency Domain} (página 237 del
visor PDF).

\textbf{Correspondencia:} Las secciones aportan el fundamento. La
diapositiva amplía la formulación o precisa condiciones que no se
presentan con idéntica notación en el libro. Desarrollo de propiedades
con convención en Hz; las secciones son el fundamento y no una tabla
idéntica a la diapositiva.

\textbf{Microcurrículo:} semana 6, sesión 1 (sesión global 11). Los
numerales del cronograma son identificadores curriculares; no son los
apartados del libro citados arriba.}
\end{frame}

\begin{frame}{Derivar resalta componentes rápidas}
\phantomsection\label{d2-017}
\[\frac{dx(t)}{dt}\longleftrightarrow j2\pi fX(f).\]

Consecuencias:

\begin{itemize}
\tightlist
\item
  frecuencia cero queda anulada;
\item
  las componentes altas reciben mayor ponderación;
\item
  el ruido de alta frecuencia puede amplificarse.
\end{itemize}

\begin{alertblock}{Atención}

Un diferenciador ayuda a localizar cambios rápidos, pero puede
deteriorar una bioseñal si no se controla el ruido.

\end{alertblock}

\note{\textbf{Libro guía:} John Semmlow, \emph{Circuits, Signals, and
Systems for Bioengineers}, 3.ª ed., 2018.

§6.5.2, \textbf{Derivative Element} (página 264 del visor PDF).

\textbf{Correspondencia:} Las secciones aportan el fundamento. La
diapositiva amplía la formulación o precisa condiciones que no se
presentan con idéntica notación en el libro. Desarrollo de propiedades
con convención en Hz; las secciones son el fundamento y no una tabla
idéntica a la diapositiva.

\textbf{Microcurrículo:} semana 6, sesión 1 (sesión global 11). Los
numerales del cronograma son identificadores curriculares; no son los
apartados del libro citados arriba.}
\end{frame}

\begin{frame}{Integrar favorece variaciones lentas}
\phantomsection\label{d2-018}
Para \(y(t)=\int_{-\infty}^{t}x(\tau)d\tau\) y condiciones de existencia
apropiadas:

\[Y(f)=\operatorname{PV}\!\left[\frac{X(f)}{j2\pi f}\right]
+\frac{X(0)}{2}\delta(f).\]

\begin{block}{Nota}

\(\operatorname{PV}\) denota el valor principal distribucional. Si la
integral acumulada no retorna a cero, su transformada puede requerir
distribuciones. No divida numéricamente el bin DC por cero.

\end{block}

La interpretación y los límites de la integración se retoman con Laplace
en la sesión 24.

\note{\textbf{Libro guía:} John Semmlow, \emph{Circuits, Signals, and
Systems for Bioengineers}, 3.ª ed., 2018.

§6.5.3, \textbf{Integrator Element} (página 266 del visor PDF); §7.2.2,
\textbf{Calculus Operations In The Laplace Domain} (página 301 del visor
PDF).

\textbf{Correspondencia:} Las secciones aportan el fundamento. La
diapositiva amplía la formulación o precisa condiciones que no se
presentan con idéntica notación en el libro. Desarrollo de propiedades
con convención en Hz; las secciones son el fundamento y no una tabla
idéntica a la diapositiva.

\textbf{Microcurrículo:} semana 6, sesión 1 (sesión global 11). Los
numerales del cronograma son identificadores curriculares; no son los
apartados del libro citados arriba.}
\end{frame}

\begin{frame}{La convolución se convierte en multiplicación}
\phantomsection\label{d2-019}
\[y(t)=x(t)*h(t)\]

\[Y(f)=X(f)H(f).\]

Interpretación:

\begin{itemize}
\tightlist
\item
  \(X(f)\) describe el contenido de la entrada;
\item
  \(H(f)\) describe la acción frecuencial del sistema;
\item
  \(Y(f)\) conserva, atenúa o cambia la fase de cada componente.
\end{itemize}

\begin{block}{Concepto clave}

El teorema de convolución conecta el análisis de señales con el análisis
de sistemas LTI.

\end{block}

\note{\textbf{Libro guía:} John Semmlow, \emph{Circuits, Signals, and
Systems for Bioengineers}, 3.ª ed., 2018.

§5.6, \textbf{Convolution In The Frequency Domain} (página 237 del visor
PDF).

\textbf{Correspondencia:} Las secciones aportan el fundamento. La
diapositiva amplía la formulación o precisa condiciones que no se
presentan con idéntica notación en el libro. Desarrollo de propiedades
con convención en Hz; las secciones son el fundamento y no una tabla
idéntica a la diapositiva.

\textbf{Microcurrículo:} semana 6, sesión 1 (sesión global 11). Los
numerales del cronograma son identificadores curriculares; no son los
apartados del libro citados arriba.}
\end{frame}

\begin{frame}{Multiplicar en tiempo produce convolución espectral}
\phantomsection\label{d2-020}
\[y(t)=x(t)w(t)\]

\[Y(f)=X(f)*W(f).\]

\begin{alertblock}{Atención}

Seleccionar un segmento mediante una ventana mezcla el espectro de la
señal con el de la ventana. Esta es la base matemática de la fuga
espectral.

\end{alertblock}

\note{\textbf{Libro guía:} John Semmlow, \emph{Circuits, Signals, and
Systems for Bioengineers}, 3.ª ed., 2018.

§4.2.3, \textbf{Data Truncation} (página 184 del visor PDF); §4.2.4,
\textbf{Data Truncation---Window Functions} (página 188 del visor PDF).

\textbf{Correspondencia:} Las secciones aportan el fundamento. La
diapositiva amplía la formulación o precisa condiciones que no se
presentan con idéntica notación en el libro. Desarrollo de propiedades
con convención en Hz; las secciones son el fundamento y no una tabla
idéntica a la diapositiva.

\textbf{Microcurrículo:} semana 6, sesión 1 (sesión global 11). Los
numerales del cronograma son identificadores curriculares; no son los
apartados del libro citados arriba.}
\end{frame}

\begin{frame}{Dualidad: tiempo y frecuencia intercambian estructuras}
\phantomsection\label{d2-021}
Bajo la convención empleada:

\[x(t)\longleftrightarrow X(f)\]

implica

\[X(t)\longleftrightarrow x(-f).\]

\begin{block}{Nota}

La dualidad ayuda a generar pares de transformada, pero requiere
atención a signos y factores según la convención.

\end{block}

\note{\textbf{Libro guía:} John Semmlow, \emph{Circuits, Signals, and
Systems for Bioengineers}, 3.ª ed., 2018.

§3.3.7, \textbf{The Continuous Fourier Transform} (página 138 del visor
PDF); §5.6, \textbf{Convolution In The Frequency Domain} (página 237 del
visor PDF).

\textbf{Correspondencia:} Las secciones aportan el fundamento. La
diapositiva amplía la formulación o precisa condiciones que no se
presentan con idéntica notación en el libro. Desarrollo de propiedades
con convención en Hz; las secciones son el fundamento y no una tabla
idéntica a la diapositiva.

\textbf{Microcurrículo:} semana 6, sesión 1 (sesión global 11). Los
numerales del cronograma son identificadores curriculares; no son los
apartados del libro citados arriba.}
\end{frame}

\begin{frame}{Tabla de predicción rápida}
\phantomsection\label{d2-022}
\begin{longtable}[]{@{}ll@{}}
\toprule\noalign{}
Operación temporal & Efecto principal en frecuencia \\
\midrule\noalign{}
\endhead
suma ponderada & suma compleja de espectros \\
retraso & fase lineal \\
compresión temporal & expansión espectral \\
derivada & ponderación por \(j2\pi f\) \\
convolución & producto de espectros \\
ventaneo & convolución con \(W(f)\) \\
\bottomrule\noalign{}
\end{longtable}

\begin{block}{Comprobación}

Antes de calcular, prediga qué aspectos deberían cambiar y cuáles
deberían conservarse.

\end{block}

\note{\textbf{Libro guía:} John Semmlow, \emph{Circuits, Signals, and
Systems for Bioengineers}, 3.ª ed., 2018.

§3.3.7, \textbf{The Continuous Fourier Transform} (página 138 del visor
PDF); §5.6, \textbf{Convolution In The Frequency Domain} (página 237 del
visor PDF).

\textbf{Correspondencia:} Las secciones aportan el fundamento. La
diapositiva amplía la formulación o precisa condiciones que no se
presentan con idéntica notación en el libro. Desarrollo de propiedades
con convención en Hz; las secciones son el fundamento y no una tabla
idéntica a la diapositiva.

\textbf{Microcurrículo:} semana 6, sesión 1 (sesión global 11). Los
numerales del cronograma son identificadores curriculares; no son los
apartados del libro citados arriba.}
\end{frame}

\begin{frame}{Caso biomédico: tres modificaciones del ECG}
\phantomsection\label{d2-023}
Considere un latido ideal \(x(t)\).

\begin{enumerate}
\tightlist
\item
  \(x(t-0.2)\) conserva magnitud y añade fase lineal.
\item
  \(x(2t)\) dura la mitad y extiende el contenido frecuencial.
\item
  \(dx/dt\) enfatiza las pendientes rápidas del QRS.
\end{enumerate}

\begin{alertblock}{Atención}

Ninguna operación garantiza mejorar detección o diagnóstico. La utilidad
depende del objetivo y de la contaminación presente.

\end{alertblock}

\note{\textbf{Libro guía:} John Semmlow, \emph{Circuits, Signals, and
Systems for Bioengineers}, 3.ª ed., 2018.

§3.3.7, \textbf{The Continuous Fourier Transform} (página 138 del visor
PDF); §5.6, \textbf{Convolution In The Frequency Domain} (página 237 del
visor PDF).

\textbf{Correspondencia:} Las secciones aportan el fundamento. La
diapositiva amplía la formulación o precisa condiciones que no se
presentan con idéntica notación en el libro. Desarrollo de propiedades
con convención en Hz; las secciones son el fundamento y no una tabla
idéntica a la diapositiva.

\textbf{Microcurrículo:} semana 6, sesión 1 (sesión global 11). Los
numerales del cronograma son identificadores curriculares; no son los
apartados del libro citados arriba.}
\end{frame}

\begin{frame}{Verificación de la sesión 11}
\phantomsection\label{d2-024}
Para \(y(t)=3x(2t-4)\):

\begin{enumerate}
\tightlist
\item
  reescriba el argumento para identificar escalamiento y retraso;
\item
  obtenga \(Y(f)\) a partir de \(X(f)\);
\item
  indique qué ocurre con el ancho espectral;
\item
  determine si la fase cambia por el desplazamiento.
\end{enumerate}

\begin{block}{Criterio de logro}

La solución debe incluir el factor de amplitud, el escalamiento
frecuencial y la fase asociada al retardo efectivo.

\end{block}

\note{\textbf{Libro guía:} John Semmlow, \emph{Circuits, Signals, and
Systems for Bioengineers}, 3.ª ed., 2018.

§3.3.7, \textbf{The Continuous Fourier Transform} (página 138 del visor
PDF); §5.6, \textbf{Convolution In The Frequency Domain} (página 237 del
visor PDF).

\textbf{Correspondencia:} Las secciones aportan el fundamento. La
diapositiva amplía la formulación o precisa condiciones que no se
presentan con idéntica notación en el libro. Desarrollo de propiedades
con convención en Hz; las secciones son el fundamento y no una tabla
idéntica a la diapositiva.

\textbf{Microcurrículo:} semana 6, sesión 1 (sesión global 11). Los
numerales del cronograma son identificadores curriculares; no son los
apartados del libro citados arriba.}
\end{frame}

\begin{frame}{Sesión 12 · La DFT analiza un bloque finito}
\phantomsection\label{d2-025}
La DFT opera sobre \(N\) muestras:

\[X[k]=\sum_{n=0}^{N-1}x[n]e^{-j2\pi kn/N},\qquad k=0,\ldots,N-1.\]

\begin{block}{Concepto clave}

La DFT produce coeficientes de una representación discreta y periódica.
No entrega un continuo de frecuencias.

\end{block}

\note{\textbf{Libro guía:} John Semmlow, \emph{Circuits, Signals, and
Systems for Bioengineers}, 3.ª ed., 2018.

§3.4.1, \textbf{The Discrete Fourier Transform} (página 142 del visor
PDF).

\textbf{Correspondencia:} El tema se desarrolla en las secciones
indicadas. La redacción es una síntesis docente.

\textbf{Microcurrículo:} semana 6, sesión 2 (sesión global 12). Los
numerales del cronograma son identificadores curriculares; no son los
apartados del libro citados arriba.}
\end{frame}

\begin{frame}{La transformada inversa reconstruye el bloque}
\phantomsection\label{d2-026}
\[x[n]=\frac{1}{N}\sum_{k=0}^{N-1}X[k]e^{j2\pi kn/N}.\]

La DFT y la IDFT forman un par exacto salvo error numérico.

\begin{alertblock}{Atención}

Esta convención asigna el factor \(1/N\) a la inversa. Otras
convenciones reparten la normalización entre ambas transformadas.

\end{alertblock}

\note{\textbf{Libro guía:} John Semmlow, \emph{Circuits, Signals, and
Systems for Bioengineers}, 3.ª ed., 2018.

§3.4.4, \textbf{The Inverse Discrete Fourier Transform} (página 163 del
visor PDF).

\textbf{Correspondencia:} El tema se desarrolla en las secciones
indicadas. La redacción es una síntesis docente.

\textbf{Microcurrículo:} semana 6, sesión 2 (sesión global 12). Los
numerales del cronograma son identificadores curriculares; no son los
apartados del libro citados arriba.}
\end{frame}

\begin{frame}{Cada índice corresponde a una frecuencia}
\phantomsection\label{d2-027}
Para muestreo uniforme a \(f_s\):

\[f_k=\frac{k}{N}f_s,\qquad \Delta f=\frac{f_s}{N}=\frac{1}{T}.\]

donde \(T=N/f_s\) es la duración.

\begin{block}{Concepto clave}

La separación entre bins depende de la duración observada. Aumentar
\(f_s\) sin aumentar \(N\) no mejora la resolución frecuencial.

\end{block}

\note{\textbf{Libro guía:} John Semmlow, \emph{Circuits, Signals, and
Systems for Bioengineers}, 3.ª ed., 2018.

§3.4.1, \textbf{The Discrete Fourier Transform} (página 142 del visor
PDF); §3.4.3, \textbf{Details Of The Dft Spectrum} (página 153 del visor
PDF).

\textbf{Correspondencia:} El tema se desarrolla en las secciones
indicadas. La redacción es una síntesis docente.

\textbf{Microcurrículo:} semana 6, sesión 2 (sesión global 12). Los
numerales del cronograma son identificadores curriculares; no son los
apartados del libro citados arriba.}
\end{frame}

\begin{frame}[fragile]{Los índices superiores representan frecuencias
negativas}
\phantomsection\label{d2-028}
La convención de \texttt{np.fft.fftfreq} es:

\[f_k=\begin{cases}
kf_s/N,&0\le k\le\lfloor(N-1)/2\rfloor,\\
(k-N)f_s/N,&\lfloor(N-1)/2\rfloor<k<N.
\end{cases}\]

\begin{block}{Nota}

Para \(N\) par, \texttt{fftfreq} etiqueta el bin de Nyquist como
\(-f_s/2\) y \texttt{rfftfreq} como \(+f_s/2\). Son la misma frecuencia
discreta. Para \(N\) impar no hay bin exactamente en Nyquist.

\end{block}

\note{\textbf{Libro guía:} John Semmlow, \emph{Circuits, Signals, and
Systems for Bioengineers}, 3.ª ed., 2018.

§3.4.1, \textbf{The Discrete Fourier Transform} (página 142 del visor
PDF); §3.4.3, \textbf{Details Of The Dft Spectrum} (página 153 del visor
PDF).

\textbf{Correspondencia:} El tema se desarrolla en las secciones
indicadas. La redacción es una síntesis docente.

\textbf{Microcurrículo:} semana 6, sesión 2 (sesión global 12). Los
numerales del cronograma son identificadores curriculares; no son los
apartados del libro citados arriba.}
\end{frame}

\begin{frame}[fragile]{Señales reales poseen simetría conjugada}
\phantomsection\label{d2-029}
Si \(x[n]\in\mathbb{R}\):

\[X[(-k)\bmod N]=X^*[k].\]

Por tanto, la mitad no negativa contiene información suficiente, con
cuidado especial en DC y Nyquist.

\begin{block}{Comprobación}

\texttt{rfft} calcula solo las frecuencias no negativas de una entrada
real y reduce costo y almacenamiento.

\end{block}

\note{\textbf{Libro guía:} John Semmlow, \emph{Circuits, Signals, and
Systems for Bioengineers}, 3.ª ed., 2018.

§3.4.1, \textbf{The Discrete Fourier Transform} (página 142 del visor
PDF); §3.4.3, \textbf{Details Of The Dft Spectrum} (página 153 del visor
PDF).

\textbf{Correspondencia:} El tema se desarrolla en las secciones
indicadas. La redacción es una síntesis docente.

\textbf{Microcurrículo:} semana 6, sesión 2 (sesión global 12). Los
numerales del cronograma son identificadores curriculares; no son los
apartados del libro citados arriba.}
\end{frame}

\begin{frame}{La FFT es un algoritmo}
\phantomsection\label{d2-030}
\begin{itemize}
\tightlist
\item
  DFT: transformación matemática.
\item
  FFT: familia de algoritmos que calcula la DFT eficientemente.
\item
  La FFT no cambia bins, resolución, fuga ni significado físico.
\end{itemize}

Complejidad aproximada:

\[\mathcal{O}(N^2)\quad\text{frente a}\quad\mathcal{O}(N\log N).\]

\begin{alertblock}{Atención}

Una FFT más rápida no corrige un segmento mal elegido ni una frecuencia
de muestreo inadecuada.

\end{alertblock}

\note{\textbf{Libro guía:} John Semmlow, \emph{Circuits, Signals, and
Systems for Bioengineers}, 3.ª ed., 2018.

§3.4.2, \textbf{Matlab Implementation Of The Discrete Fourier Transform}
(página 148 del visor PDF).

\textbf{Correspondencia:} El tema se desarrolla en las secciones
indicadas. La redacción es una síntesis docente.

\textbf{Microcurrículo:} semana 6, sesión 2 (sesión global 12). Los
numerales del cronograma son identificadores curriculares; no son los
apartados del libro citados arriba.}
\end{frame}

\begin{frame}{La DFT supone extensión periódica del bloque}
\phantomsection\label{d2-031}
El bloque \(x[0],\ldots,x[N-1]\) se interpreta como un periodo de una
secuencia periódica.

Si los extremos no enlazan suavemente, la repetición introduce
discontinuidades.

\begin{block}{Concepto clave}

La fuga espectral surge porque el bloque finito y su extensión periódica
rara vez contienen un número entero de ciclos.

\end{block}

\note{\textbf{Libro guía:} John Semmlow, \emph{Circuits, Signals, and
Systems for Bioengineers}, 3.ª ed., 2018.

§4.2.3, \textbf{Data Truncation} (página 184 del visor PDF).

\textbf{Correspondencia:} El tema se desarrolla en las secciones
indicadas. La redacción es una síntesis docente.

\textbf{Microcurrículo:} semana 6, sesión 2 (sesión global 12). Los
numerales del cronograma son identificadores curriculares; no son los
apartados del libro citados arriba.}
\end{frame}

\begin{frame}{Magnitud bilateral y unilateral}
\phantomsection\label{d2-032}
Con \(X[k]=\mathrm{DFT}\{x[n]\}\):

\[A_{2}[k]=\frac{|X[k]|}{N}.\]

Para espectro unilateral de una señal real, se duplican los bins
interiores:

\[A_{1}[k]=2A_{2}[k],\]

excepto DC y, si existe, Nyquist.

\begin{alertblock}{Atención}

La amplitud pico se recupera así para tonos aislados alineados con bins
y ventana rectangular. Con otra ventana se corrige la ganancia
coherente; los tonos no alineados sufren fuga y pérdida de pico.

\end{alertblock}

\note{\textbf{Libro guía:} John Semmlow, \emph{Circuits, Signals, and
Systems for Bioengineers}, 3.ª ed., 2018.

§3.4.3, \textbf{Details Of The Dft Spectrum} (página 153 del visor PDF).

\textbf{Correspondencia:} El tema se desarrolla en las secciones
indicadas. La redacción es una síntesis docente.

\textbf{Microcurrículo:} semana 6, sesión 2 (sesión global 12). Los
numerales del cronograma son identificadores curriculares; no son los
apartados del libro citados arriba.}
\end{frame}

\begin{frame}{La fase necesita una referencia de amplitud}
\phantomsection\label{d2-033}
\[\phi[k]=\operatorname{angle}(X[k]).\]

En bins donde \(|X[k]|\) es casi cero, la fase resulta numéricamente
inestable.

\begin{block}{Nota}

La fase debe interpretarse solo donde la magnitud o coherencia aporten
evidencia suficiente.

\end{block}

\note{\textbf{Libro guía:} John Semmlow, \emph{Circuits, Signals, and
Systems for Bioengineers}, 3.ª ed., 2018.

§3.4.3, \textbf{Details Of The Dft Spectrum} (página 153 del visor PDF).

\textbf{Correspondencia:} El tema se desarrolla en las secciones
indicadas. La redacción es una síntesis docente.

\textbf{Microcurrículo:} semana 6, sesión 2 (sesión global 12). Los
numerales del cronograma son identificadores curriculares; no son los
apartados del libro citados arriba.}
\end{frame}

\begin{frame}{Convolución circular y periodicidad}
\phantomsection\label{d2-034}
La multiplicación de DFT corresponde a convolución circular:

\[\mathrm{IDFT}\{X[k]H[k]\}=x[n]\circledast h[n].\]

Para obtener convolución lineal sin solapamiento circular, se usa
relleno con ceros hasta al menos \(N_x+N_h-1\).

\begin{block}{Concepto clave}

El relleno evita plegamiento circular en convolución, pero no añade
información espectral al registro.

\end{block}

\note{\textbf{Libro guía:} John Semmlow, \emph{Circuits, Signals, and
Systems for Bioengineers}, 3.ª ed., 2018.

§5.6, \textbf{Convolution In The Frequency Domain} (página 237 del visor
PDF).

\textbf{Correspondencia:} Las secciones aportan el fundamento. La
diapositiva amplía la formulación o precisa condiciones que no se
presentan con idéntica notación en el libro. Convolución circular
mediante DFT: extensión del teorema para bloques finitos. El relleno
especificado evita aliasing circular.

\textbf{Microcurrículo:} semana 6, sesión 2 (sesión global 12). Los
numerales del cronograma son identificadores curriculares; no son los
apartados del libro citados arriba.}
\end{frame}

\begin{frame}[fragile]{Código: eje y espectro unilateral}
\phantomsection\label{d2-035}
\begin{Shaded}
\begin{Highlighting}[]
\ImportTok{import}\NormalTok{ numpy }\ImportTok{as}\NormalTok{ np}

\NormalTok{x0 }\OperatorTok{=}\NormalTok{ x }\OperatorTok{{-}}\NormalTok{ np.mean(x)}
\NormalTok{N }\OperatorTok{=} \BuiltInTok{len}\NormalTok{(x0)}
\NormalTok{X }\OperatorTok{=}\NormalTok{ np.fft.rfft(x0)}
\NormalTok{f }\OperatorTok{=}\NormalTok{ np.fft.rfftfreq(N, d}\OperatorTok{=}\DecValTok{1}\OperatorTok{/}\NormalTok{fs)}
\NormalTok{A }\OperatorTok{=}\NormalTok{ np.}\BuiltInTok{abs}\NormalTok{(X) }\OperatorTok{/}\NormalTok{ N}
\ControlFlowTok{if}\NormalTok{ N }\OperatorTok{\%} \DecValTok{2} \OperatorTok{==} \DecValTok{0}\NormalTok{:}
\NormalTok{    A[}\DecValTok{1}\NormalTok{:}\OperatorTok{{-}}\DecValTok{1}\NormalTok{] }\OperatorTok{*=} \DecValTok{2}
\ControlFlowTok{else}\NormalTok{:}
\NormalTok{    A[}\DecValTok{1}\NormalTok{:] }\OperatorTok{*=} \DecValTok{2}
\NormalTok{fase }\OperatorTok{=}\NormalTok{ np.angle(X)}
\end{Highlighting}
\end{Shaded}

\begin{alertblock}{Atención}

El código trata la paridad. La interpretación de amplitud supone ventana
rectangular y tonos alineados con bins.

\end{alertblock}

\note{\textbf{Libro guía:} John Semmlow, \emph{Circuits, Signals, and
Systems for Bioengineers}, 3.ª ed., 2018.

§3.4.2, \textbf{Matlab Implementation Of The Discrete Fourier Transform}
(página 148 del visor PDF); §3.4.3, \textbf{Details Of The Dft Spectrum}
(página 153 del visor PDF).

\textbf{Correspondencia:} Ejemplo, figura o implementación de
elaboración docente a partir de estos fundamentos. No es una
transcripción del código MATLAB del libro.

\textbf{Microcurrículo:} semana 6, sesión 2 (sesión global 12). Los
numerales del cronograma son identificadores curriculares; no son los
apartados del libro citados arriba.}
\end{frame}

\begin{frame}[fragile]{Prueba de amplitud para longitudes par e impar}
\phantomsection\label{d2-036}
\begin{Shaded}
\begin{Highlighting}[]
\ControlFlowTok{for}\NormalTok{ N\_test }\KeywordTok{in}\NormalTok{ (}\DecValTok{999}\NormalTok{, }\DecValTok{1000}\NormalTok{):}
\NormalTok{    n }\OperatorTok{=}\NormalTok{ np.arange(N\_test)}
\NormalTok{    xt }\OperatorTok{=} \DecValTok{2}\OperatorTok{*}\NormalTok{np.cos(}\DecValTok{2}\OperatorTok{*}\NormalTok{np.pi}\OperatorTok{*}\DecValTok{17}\OperatorTok{*}\NormalTok{n}\OperatorTok{/}\NormalTok{N\_test)}
\NormalTok{    At }\OperatorTok{=}\NormalTok{ np.}\BuiltInTok{abs}\NormalTok{(np.fft.rfft(xt))}\OperatorTok{/}\NormalTok{N\_test}
    \ControlFlowTok{if}\NormalTok{ N\_test }\OperatorTok{\%} \DecValTok{2} \OperatorTok{==} \DecValTok{0}\NormalTok{:}
\NormalTok{        At[}\DecValTok{1}\NormalTok{:}\OperatorTok{{-}}\DecValTok{1}\NormalTok{] }\OperatorTok{*=} \DecValTok{2}
    \ControlFlowTok{else}\NormalTok{:}
\NormalTok{        At[}\DecValTok{1}\NormalTok{:] }\OperatorTok{*=} \DecValTok{2}
    \ControlFlowTok{assert}\NormalTok{ np.isclose(At[}\DecValTok{17}\NormalTok{], }\FloatTok{2.0}\NormalTok{)}
\end{Highlighting}
\end{Shaded}

La prueba verifica recuperación de una amplitud conocida con longitud
par e impar. No evalúa fuga ni mezclas de fuentes.

\note{\textbf{Libro guía:} John Semmlow, \emph{Circuits, Signals, and
Systems for Bioengineers}, 3.ª ed., 2018.

§3.4.2, \textbf{Matlab Implementation Of The Discrete Fourier Transform}
(página 148 del visor PDF); §3.4.3, \textbf{Details Of The Dft Spectrum}
(página 153 del visor PDF).

\textbf{Correspondencia:} Ejemplo, figura o implementación de
elaboración docente a partir de estos fundamentos. No es una
transcripción del código MATLAB del libro.

\textbf{Microcurrículo:} semana 6, sesión 2 (sesión global 12). Los
numerales del cronograma son identificadores curriculares; no son los
apartados del libro citados arriba.}
\end{frame}

\begin{frame}[fragile]{Verificación de la sesión 12}
\phantomsection\label{d2-037}
Un registro tiene \(f_s=500\) Hz y \(N=2000\).

\begin{enumerate}
\tightlist
\item
  calcule su duración;
\item
  calcule \(\Delta f\);
\item
  identifique el rango de \texttt{rfft};
\item
  explique qué mejora y qué no mejora al usar 4096 puntos mediante zero
  padding;
\item
  indique por qué la fase de bins pequeños puede ser inestable.
\end{enumerate}

\begin{block}{Criterio de logro}

Debe distinguir duración, resolución, densidad gráfica y costo
computacional.

\end{block}

\note{\textbf{Libro guía:} John Semmlow, \emph{Circuits, Signals, and
Systems for Bioengineers}, 3.ª ed., 2018.

§3.4.1, \textbf{The Discrete Fourier Transform} (página 142 del visor
PDF); §3.4.2, \textbf{Matlab Implementation Of The Discrete Fourier
Transform} (página 148 del visor PDF); §3.4.3, \textbf{Details Of The
Dft Spectrum} (página 153 del visor PDF).

\textbf{Correspondencia:} El tema se desarrolla en las secciones
indicadas. La redacción es una síntesis docente.

\textbf{Microcurrículo:} semana 6, sesión 2 (sesión global 12). Los
numerales del cronograma son identificadores curriculares; no son los
apartados del libro citados arriba.}
\end{frame}

\begin{frame}{Cierre de la semana 6}
\phantomsection\label{d2-038}
\begin{itemize}
\tightlist
\item
  Las propiedades conectan operaciones temporales con cambios
  espectrales.
\item
  El retraso actúa sobre fase; el escalamiento temporal modifica el
  ancho.
\item
  La convolución temporal equivale a multiplicación frecuencial.
\item
  La DFT representa un bloque finito sobre bins discretos.
\item
  La FFT acelera el cálculo sin alterar la interpretación.
\end{itemize}

\begin{block}{Pregunta de salida}

¿Por qué dos registros con la misma frecuencia de muestreo pueden tener
resoluciones frecuenciales distintas?

\end{block}

\note{\textbf{Libro guía:} John Semmlow, \emph{Circuits, Signals, and
Systems for Bioengineers}, 3.ª ed., 2018.

§3.3.7, \textbf{The Continuous Fourier Transform} (página 138 del visor
PDF); §3.4, \textbf{Time--Frequency Transformation In The Discrete
Domain} (página 142 del visor PDF).

\textbf{Correspondencia:} Síntesis o encuadre que reúne varias
secciones; no corresponde a un único apartado del libro.

\textbf{Microcurrículo:} semana 6, sesión 2 (sesión global 12). Los
numerales del cronograma son identificadores curriculares; no son los
apartados del libro citados arriba.}
\end{frame}

\begin{frame}{Semana 7 · Resolución, fuga y adquisición}
\phantomsection\label{d2-039}
\note{\textbf{Libro guía:} John Semmlow, \emph{Circuits, Signals, and
Systems for Bioengineers}, 3.ª ed., 2018.

§3.4, \textbf{Time--Frequency Transformation In The Discrete Domain}
(página 142 del visor PDF); §4.2, \textbf{Data Acquisition And Storage}
(página 175 del visor PDF); §4.6, \textbf{Time--Frequency Analysis}
(página 205 del visor PDF); §5.4, \textbf{Using The Impulse Response To
Find A Systems Output To Any Input---Convolution} (página 219 del visor
PDF).

\textbf{Correspondencia:} Síntesis o encuadre que reúne varias
secciones; no corresponde a un único apartado del libro.

\textbf{Lectura:} use estas secciones como ruta de preparación y
consulta. La bibliografía aparece al final del bloque.}
\end{frame}

\begin{frame}{Sesión 13 · La resolución proviene de la duración}
\phantomsection\label{d2-040}
\[\Delta f=\frac{f_s}{N}=\frac{1}{T}.\]

Ejemplo: \(T=10\) s produce bins separados 0.1 Hz, con independencia de
la forma usada para calcular \(N=f_sT\).

\begin{block}{Concepto clave}

Observar durante más tiempo permite separar componentes más próximas,
siempre que la señal mantenga propiedades suficientemente estables.

\end{block}

\note{\textbf{Libro guía:} John Semmlow, \emph{Circuits, Signals, and
Systems for Bioengineers}, 3.ª ed., 2018.

§3.4.3, \textbf{Details Of The Dft Spectrum} (página 153 del visor PDF);
§4.2.3, \textbf{Data Truncation} (página 184 del visor PDF).

\textbf{Correspondencia:} El tema se desarrolla en las secciones
indicadas. La redacción es una síntesis docente.

\textbf{Microcurrículo:} semana 7, sesión 1 (sesión global 13). Los
numerales del cronograma son identificadores curriculares; no son los
apartados del libro citados arriba.}
\end{frame}

\begin{frame}{Resolución nominal no garantiza separación efectiva}
\phantomsection\label{d2-041}
Dos tonos separados por \(\Delta f\) pueden seguir sin distinguirse si:

\begin{itemize}
\tightlist
\item
  la ventana ensancha el lóbulo principal;
\item
  una componente es mucho mayor que la otra;
\item
  el ruido oculta la componente débil;
\item
  la frecuencia cambia durante el registro.
\end{itemize}

\begin{alertblock}{Atención}

El ancho del lóbulo principal y el nivel de lóbulos laterales importan
tanto como el espaciado entre bins.

\end{alertblock}

\note{\textbf{Libro guía:} John Semmlow, \emph{Circuits, Signals, and
Systems for Bioengineers}, 3.ª ed., 2018.

§4.2.4, \textbf{Data Truncation---Window Functions} (página 188 del
visor PDF).

\textbf{Correspondencia:} El tema se desarrolla en las secciones
indicadas. La redacción es una síntesis docente.

\textbf{Microcurrículo:} semana 7, sesión 1 (sesión global 13). Los
numerales del cronograma son identificadores curriculares; no son los
apartados del libro citados arriba.}
\end{frame}

\begin{frame}{Una frecuencia alineada con un bin produce un caso
especial}
\phantomsection\label{d2-042}
Si una cosenoide completa un número entero de ciclos en \(T\):

\[f_0=m\Delta f,\]

su energía coincide con bins de la DFT bajo una ventana rectangular.

\begin{block}{Nota}

Este caso facilita pruebas de escalamiento, pero rara vez describe
exactamente una bioseñal real.

\end{block}

\note{\textbf{Libro guía:} John Semmlow, \emph{Circuits, Signals, and
Systems for Bioengineers}, 3.ª ed., 2018.

§3.4.1, \textbf{The Discrete Fourier Transform} (página 142 del visor
PDF); §4.2.3, \textbf{Data Truncation} (página 184 del visor PDF).

\textbf{Correspondencia:} El tema se desarrolla en las secciones
indicadas. La redacción es una síntesis docente.

\textbf{Microcurrículo:} semana 7, sesión 1 (sesión global 13). Los
numerales del cronograma son identificadores curriculares; no son los
apartados del libro citados arriba.}
\end{frame}

\begin{frame}{La fuga distribuye energía entre bins}
\phantomsection\label{d2-043}
Una componente no alineada con los bins produce valores en múltiples
frecuencias.

La causa matemática es el producto temporal:

\[x_w[n]=x[n]w[n],\]

que corresponde a convolución espectral.

\begin{block}{Concepto clave}

La fuga no implica que la señal contenga todas esas frecuencias como
oscilaciones independientes.

\end{block}

\note{\textbf{Libro guía:} John Semmlow, \emph{Circuits, Signals, and
Systems for Bioengineers}, 3.ª ed., 2018.

§4.2.3, \textbf{Data Truncation} (página 184 del visor PDF).

\textbf{Correspondencia:} El tema se desarrolla en las secciones
indicadas. La redacción es una síntesis docente.

\textbf{Microcurrículo:} semana 7, sesión 1 (sesión global 13). Los
numerales del cronograma son identificadores curriculares; no son los
apartados del libro citados arriba.}
\end{frame}

\begin{frame}{Ventana rectangular: resolución y lóbulos laterales}
\phantomsection\label{d2-044}
La ventana rectangular equivale a recortar el registro sin suavizar
extremos.

\begin{itemize}
\tightlist
\item
  lóbulo principal estrecho;
\item
  lóbulos laterales relativamente altos;
\item
  buena separación para tonos comparables y alineados;
\item
  mayor contaminación alrededor de componentes intensas.
\end{itemize}

\begin{alertblock}{Atención}

Usar ``sin ventana'' significa usar una ventana rectangular.

\end{alertblock}

\note{\textbf{Libro guía:} John Semmlow, \emph{Circuits, Signals, and
Systems for Bioengineers}, 3.ª ed., 2018.

§4.2.4, \textbf{Data Truncation---Window Functions} (página 188 del
visor PDF).

\textbf{Correspondencia:} El tema se desarrolla en las secciones
indicadas. La redacción es una síntesis docente.

\textbf{Microcurrículo:} semana 7, sesión 1 (sesión global 13). Los
numerales del cronograma son identificadores curriculares; no son los
apartados del libro citados arriba.}
\end{frame}

\begin{frame}{Hann reduce lóbulos laterales y ensancha el principal}
\phantomsection\label{d2-045}
\[w[n]=\frac{1}{2}\left(1-\cos\frac{2\pi n}{N-1}\right).\]

La ventana lleva los extremos hacia cero y reduce discontinuidades en la
extensión periódica.

\begin{block}{Concepto clave}

La reducción de fuga se paga con menor capacidad para separar
componentes cercanas.

\end{block}

\note{\textbf{Libro guía:} John Semmlow, \emph{Circuits, Signals, and
Systems for Bioengineers}, 3.ª ed., 2018.

§4.2.4, \textbf{Data Truncation---Window Functions} (página 188 del
visor PDF).

\textbf{Correspondencia:} El tema se desarrolla en las secciones
indicadas. La redacción es una síntesis docente.

\textbf{Microcurrículo:} semana 7, sesión 1 (sesión global 13). Los
numerales del cronograma son identificadores curriculares; no son los
apartados del libro citados arriba.}
\end{frame}

\begin{frame}{Hamming y Blackman ofrecen otros compromisos}
\phantomsection\label{d2-046}
\begin{longtable}[]{@{}
  >{\raggedright\arraybackslash}p{(\linewidth - 4\tabcolsep) * \real{0.3333}}
  >{\raggedright\arraybackslash}p{(\linewidth - 4\tabcolsep) * \real{0.3333}}
  >{\raggedright\arraybackslash}p{(\linewidth - 4\tabcolsep) * \real{0.3333}}@{}}
\toprule\noalign{}
\begin{minipage}[b]{\linewidth}\raggedright
Ventana
\end{minipage} & \begin{minipage}[b]{\linewidth}\raggedright
Uso orientativo
\end{minipage} & \begin{minipage}[b]{\linewidth}\raggedright
Compromiso
\end{minipage} \\
\midrule\noalign{}
\endhead
Rectangular & tonos alineados y cercanos & lóbulos laterales altos \\
Hann & análisis general y Welch & balance habitual \\
Hamming & atenuación del primer lateral & extremos no nulos \\
Blackman & componente débil junto a fuerte & lóbulo principal más
ancho \\
\bottomrule\noalign{}
\end{longtable}

\begin{alertblock}{Atención}

La selección no debe basarse en una lista memorizada. Debe responder a
separación, rango dinámico y objetivo de amplitud.

\end{alertblock}

\note{\textbf{Libro guía:} John Semmlow, \emph{Circuits, Signals, and
Systems for Bioengineers}, 3.ª ed., 2018.

§4.2.4, \textbf{Data Truncation---Window Functions} (página 188 del
visor PDF).

\textbf{Correspondencia:} El tema se desarrolla en las secciones
indicadas. La redacción es una síntesis docente.

\textbf{Microcurrículo:} semana 7, sesión 1 (sesión global 13). Los
numerales del cronograma son identificadores curriculares; no son los
apartados del libro citados arriba.}
\end{frame}

\begin{frame}{La ganancia coherente corrige amplitud}
\phantomsection\label{d2-047}
Una ventana reduce la amplitud de un tono. Su ganancia coherente es

\[G_c=\frac{1}{N}\sum_{n=0}^{N-1}w[n].\]

Para estimar amplitud se corrige por \(G_c\) cuando el modelo de tono y
alineación lo justifican.

\begin{block}{Nota}

La corrección de amplitud no elimina fuga ni convierte una señal no
estacionaria en estacionaria.

\end{block}

\note{\textbf{Libro guía:} John Semmlow, \emph{Circuits, Signals, and
Systems for Bioengineers}, 3.ª ed., 2018.

§4.2.4, \textbf{Data Truncation---Window Functions} (página 188 del
visor PDF); §4.3, \textbf{Power Spectrum} (página 191 del visor PDF).

\textbf{Correspondencia:} Las secciones aportan el fundamento. La
diapositiva amplía la formulación o precisa condiciones que no se
presentan con idéntica notación en el libro. Ganancia coherente y ENBW
formalizan el escalamiento de ventanas. Consulte también la
documentación oficial de SciPy sobre escalamiento espectral.

\textbf{Microcurrículo:} semana 7, sesión 1 (sesión global 13). Los
numerales del cronograma son identificadores curriculares; no son los
apartados del libro citados arriba.}
\end{frame}

\begin{frame}{El ancho de banda equivalente afecta el ruido}
\phantomsection\label{d2-048}
Una forma discreta normalizada en bins es

\[\mathrm{ENBW}=N\frac{\sum_n w^2[n]}{\left(\sum_n w[n]\right)^2}.\]

Una ENBW mayor integra más potencia de ruido por bin espectral.

\begin{block}{Concepto clave}

La ventana afecta tanto la forma de una componente tonal como el piso de
ruido estimado.

\end{block}

\note{\textbf{Libro guía:} John Semmlow, \emph{Circuits, Signals, and
Systems for Bioengineers}, 3.ª ed., 2018.

§4.2.4, \textbf{Data Truncation---Window Functions} (página 188 del
visor PDF); §4.3, \textbf{Power Spectrum} (página 191 del visor PDF).

\textbf{Correspondencia:} Las secciones aportan el fundamento. La
diapositiva amplía la formulación o precisa condiciones que no se
presentan con idéntica notación en el libro. Ganancia coherente y ENBW
formalizan el escalamiento de ventanas. Consulte también la
documentación oficial de SciPy sobre escalamiento espectral.

\textbf{Microcurrículo:} semana 7, sesión 1 (sesión global 13). Los
numerales del cronograma son identificadores curriculares; no son los
apartados del libro citados arriba.}
\end{frame}

\begin{frame}{Zero padding densifica el eje}
\phantomsection\label{d2-049}
Agregar ceros hasta \(N_{FFT}>N\) produce separación gráfica

\[\Delta f_{grid}=\frac{f_s}{N_{FFT}},\]

pero la resolución física sigue determinada principalmente por
\(T=N/f_s\) y la ventana.

\begin{alertblock}{Atención}

Una curva más suave no contiene más observación ni permite separar
componentes que el registro original no resolvía.

\end{alertblock}

\note{\textbf{Libro guía:} John Semmlow, \emph{Circuits, Signals, and
Systems for Bioengineers}, 3.ª ed., 2018.

§3.4.3, \textbf{Details Of The Dft Spectrum} (página 153 del visor PDF);
§4.2.3, \textbf{Data Truncation} (página 184 del visor PDF).

\textbf{Correspondencia:} El tema se desarrolla en las secciones
indicadas. La redacción es una síntesis docente.

\textbf{Microcurrículo:} semana 7, sesión 1 (sesión global 13). Los
numerales del cronograma son identificadores curriculares; no son los
apartados del libro citados arriba.}
\end{frame}

\begin{frame}{Duraciones largas también tienen costos}
\phantomsection\label{d2-050}
\begin{itemize}
\tightlist
\item
  pueden mezclar estados fisiológicos;
\item
  reducen localización temporal;
\item
  aumentan probabilidad de artefactos;
\item
  exigen mayor estabilidad del reloj y del protocolo.
\end{itemize}

\begin{block}{Concepto clave}

La duración óptima surge del equilibrio entre resolución frecuencial y
estacionariedad local.

\end{block}

\note{\textbf{Libro guía:} John Semmlow, \emph{Circuits, Signals, and
Systems for Bioengineers}, 3.ª ed., 2018.

§4.2.3, \textbf{Data Truncation} (página 184 del visor PDF); §10.2,
\textbf{Stochastic Processes: Stationarity And Ergodicity} (página 450
del visor PDF).

\textbf{Correspondencia:} El tema se desarrolla en las secciones
indicadas. La redacción es una síntesis docente.

\textbf{Microcurrículo:} semana 7, sesión 1 (sesión global 13). Los
numerales del cronograma son identificadores curriculares; no son los
apartados del libro citados arriba.}
\end{frame}

\begin{frame}{Ejemplo biomédico: ritmo respiratorio}
\phantomsection\label{d2-051}
Para distinguir 0.20 Hz de 0.25 Hz se necesita una observación
suficientemente larga.

\begin{itemize}
\tightlist
\item
  con \(T=10\) s, \(\Delta f=0.1\) Hz;
\item
  con \(T=40\) s, \(\Delta f=0.025\) Hz.
\end{itemize}

\begin{alertblock}{Atención}

La mejor separación nominal de 40 s solo es útil si el ritmo no cambia
de manera importante durante esos 40 s.

\end{alertblock}

\note{\textbf{Libro guía:} John Semmlow, \emph{Circuits, Signals, and
Systems for Bioengineers}, 3.ª ed., 2018.

§4.2.3, \textbf{Data Truncation} (página 184 del visor PDF).

\textbf{Correspondencia:} Actividad o interpretación biomédica de
elaboración docente apoyada en las secciones indicadas.

\textbf{Microcurrículo:} semana 7, sesión 1 (sesión global 13). Los
numerales del cronograma son identificadores curriculares; no son los
apartados del libro citados arriba.}
\end{frame}

\begin{frame}[fragile]{Código: comparar ventanas}
\phantomsection\label{d2-052}
\begin{Shaded}
\begin{Highlighting}[]
\ImportTok{import}\NormalTok{ numpy }\ImportTok{as}\NormalTok{ np}
\ImportTok{from}\NormalTok{ scipy.signal }\ImportTok{import}\NormalTok{ get\_window}

\NormalTok{N }\OperatorTok{=} \BuiltInTok{len}\NormalTok{(x)}
\ControlFlowTok{for}\NormalTok{ nombre }\KeywordTok{in}\NormalTok{ [}\StringTok{"boxcar"}\NormalTok{, }\StringTok{"hann"}\NormalTok{, }\StringTok{"hamming"}\NormalTok{, }\StringTok{"blackman"}\NormalTok{]:}
\NormalTok{    w }\OperatorTok{=}\NormalTok{ get\_window(nombre, N, fftbins}\OperatorTok{=}\VariableTok{False}\NormalTok{)}
\NormalTok{    Xw }\OperatorTok{=}\NormalTok{ np.fft.rfft((x }\OperatorTok{{-}}\NormalTok{ np.mean(x)) }\OperatorTok{*}\NormalTok{ w)}
\NormalTok{    Gc }\OperatorTok{=}\NormalTok{ np.mean(w)}
\NormalTok{    A }\OperatorTok{=}\NormalTok{ np.}\BuiltInTok{abs}\NormalTok{(Xw) }\OperatorTok{/}\NormalTok{ (N }\OperatorTok{*}\NormalTok{ Gc)}
    \ControlFlowTok{if}\NormalTok{ N }\OperatorTok{\%} \DecValTok{2} \OperatorTok{==} \DecValTok{0}\NormalTok{:}
\NormalTok{        A[}\DecValTok{1}\NormalTok{:}\OperatorTok{{-}}\DecValTok{1}\NormalTok{] }\OperatorTok{*=} \DecValTok{2}
    \ControlFlowTok{else}\NormalTok{:}
\NormalTok{        A[}\DecValTok{1}\NormalTok{:] }\OperatorTok{*=} \DecValTok{2}
\end{Highlighting}
\end{Shaded}

\begin{block}{Nota}

La comparación debe mantener el mismo segmento y declarar la corrección
aplicada.

\end{block}

\note{\textbf{Libro guía:} John Semmlow, \emph{Circuits, Signals, and
Systems for Bioengineers}, 3.ª ed., 2018.

§4.2.4, \textbf{Data Truncation---Window Functions} (página 188 del
visor PDF).

\textbf{Correspondencia:} Ejemplo, figura o implementación de
elaboración docente a partir de estos fundamentos. No es una
transcripción del código MATLAB del libro.

\textbf{Microcurrículo:} semana 7, sesión 1 (sesión global 13). Los
numerales del cronograma son identificadores curriculares; no son los
apartados del libro citados arriba.}
\end{frame}

\begin{frame}{Verificación de la sesión 13}
\phantomsection\label{d2-053}
Un EEG se analiza con \(f_s=256\) Hz y ventanas de 2 s.

\begin{enumerate}
\tightlist
\item
  determine \(N\) y \(\Delta f\);
\item
  explique si zero padding a 2048 puntos cambia la resolución real;
\item
  compare el efecto esperado de rectangular y Hann;
\item
  proponga una duración para separar 9.8 Hz de 10.2 Hz;
\item
  discuta el supuesto fisiológico de esa duración.
\end{enumerate}

\begin{block}{Criterio de logro}

La respuesta debe separar espaciado de bins, ancho del lóbulo y
estabilidad temporal.

\end{block}

\note{\textbf{Libro guía:} John Semmlow, \emph{Circuits, Signals, and
Systems for Bioengineers}, 3.ª ed., 2018.

§4.2.3, \textbf{Data Truncation} (página 184 del visor PDF); §4.2.4,
\textbf{Data Truncation---Window Functions} (página 188 del visor PDF).

\textbf{Correspondencia:} El tema se desarrolla en las secciones
indicadas. La redacción es una síntesis docente.

\textbf{Microcurrículo:} semana 7, sesión 1 (sesión global 13). Los
numerales del cronograma son identificadores curriculares; no son los
apartados del libro citados arriba.}
\end{frame}

\begin{frame}{Sesión 14 · La cadena de adquisición condiciona el dato}
\phantomsection\label{d2-054}
\begin{enumerate}
\tightlist
\item
  El fenómeno fisiológico origina una magnitud observable.
\item
  El sensor y el acoplamiento convierten esa magnitud.
\item
  El acondicionamiento limita banda y adapta nivel.
\item
  El ADC discretiza tiempo y amplitud.
\item
  El registro conserva valores, unidades y metadatos.
\end{enumerate}

\begin{block}{Concepto clave}

Cada etapa limita la información disponible para el procesamiento
posterior.

\end{block}

\note{\textbf{Libro guía:} John Semmlow, \emph{Circuits, Signals, and
Systems for Bioengineers}, 3.ª ed., 2018.

§1.2.2, \textbf{Biotransducers And Common Physiological Measurements}
(página 12 del visor PDF); §4.2.1, \textbf{Data Sampling: The Sampling
Theorem} (página 175 del visor PDF); §4.2.2, \textbf{Amplitude
Slicing---Quantization} (página 182 del visor PDF).

\textbf{Correspondencia:} El tema se desarrolla en las secciones
indicadas. La redacción es una síntesis docente.

\textbf{Microcurrículo:} semana 7, sesión 2 (sesión global 14). Los
numerales del cronograma son identificadores curriculares; no son los
apartados del libro citados arriba.}
\end{frame}

\begin{frame}{El sensor define sensibilidad y ancho de banda}
\phantomsection\label{d2-055}
Modelo local:

\[v_s(t)=Sx(t)+b+n_s(t).\]

\begin{itemize}
\tightlist
\item
  \(S\): sensibilidad con unidades físicas.
\item
  \(b\): offset.
\item
  \(n_s(t)\): perturbación propia del sensor.
\end{itemize}

\begin{alertblock}{Atención}

La relación lineal suele ser válida solo dentro de un rango. Saturación,
histéresis y acoplamiento alteran el modelo.

\end{alertblock}

\note{\textbf{Libro guía:} John Semmlow, \emph{Circuits, Signals, and
Systems for Bioengineers}, 3.ª ed., 2018.

§1.2.2, \textbf{Biotransducers And Common Physiological Measurements}
(página 12 del visor PDF).

\textbf{Correspondencia:} Actividad o interpretación biomédica de
elaboración docente apoyada en las secciones indicadas.

\textbf{Microcurrículo:} semana 7, sesión 2 (sesión global 14). Los
numerales del cronograma son identificadores curriculares; no son los
apartados del libro citados arriba.}
\end{frame}

\begin{frame}{El acoplamiento forma parte de la medición}
\phantomsection\label{d2-056}
Ejemplos:

\begin{itemize}
\tightlist
\item
  impedancia electrodo-piel en ECG y EEG;
\item
  presión de contacto en PPG;
\item
  orientación y fijación en acelerometría;
\item
  posición de banda en respiración.
\end{itemize}

\begin{block}{Concepto clave}

Un sensor correcto con acoplamiento deficiente puede producir un
registro inválido.

\end{block}

\note{\textbf{Libro guía:} John Semmlow, \emph{Circuits, Signals, and
Systems for Bioengineers}, 3.ª ed., 2018.

§1.2.2, \textbf{Biotransducers And Common Physiological Measurements}
(página 12 del visor PDF).

\textbf{Correspondencia:} Actividad o interpretación biomédica de
elaboración docente apoyada en las secciones indicadas.

\textbf{Microcurrículo:} semana 7, sesión 2 (sesión global 14). Los
numerales del cronograma son identificadores curriculares; no son los
apartados del libro citados arriba.}
\end{frame}

\begin{frame}{El acondicionamiento protege y adapta la señal}
\phantomsection\label{d2-057}
Funciones frecuentes:

\begin{itemize}
\tightlist
\item
  amplificación;
\item
  rechazo de modo común;
\item
  limitación de banda;
\item
  ajuste de nivel y protección;
\item
  aislamiento cuando corresponde.
\end{itemize}

\begin{block}{Nota}

La ganancia debe usar el rango del convertidor sin recortar transitorios
fisiológicos ni artefactos plausibles.

\end{block}

\note{\textbf{Libro guía:} John Semmlow, \emph{Circuits, Signals, and
Systems for Bioengineers}, 3.ª ed., 2018.

§1.2.2, \textbf{Biotransducers And Common Physiological Measurements}
(página 12 del visor PDF).

\textbf{Correspondencia:} Actividad o interpretación biomédica de
elaboración docente apoyada en las secciones indicadas.

\textbf{Microcurrículo:} semana 7, sesión 2 (sesión global 14). Los
numerales del cronograma son identificadores curriculares; no son los
apartados del libro citados arriba.}
\end{frame}

\begin{frame}{El filtro antialias actúa antes del ADC}
\phantomsection\label{d2-058}
La etapa analógica atenúa componentes por encima de la banda que el
sistema digital puede representar.

\begin{alertblock}{Atención}

Un filtro digital posterior no puede distinguir una frecuencia real de
su alias cuando ambas ya producen las mismas muestras.

\end{alertblock}

\note{\textbf{Libro guía:} John Semmlow, \emph{Circuits, Signals, and
Systems for Bioengineers}, 3.ª ed., 2018.

§4.2.1, \textbf{Data Sampling: The Sampling Theorem} (página 175 del
visor PDF).

\textbf{Correspondencia:} El tema se desarrolla en las secciones
indicadas. La redacción es una síntesis docente.

\textbf{Microcurrículo:} semana 7, sesión 2 (sesión global 14). Los
numerales del cronograma son identificadores curriculares; no son los
apartados del libro citados arriba.}
\end{frame}

\begin{frame}{El ADC discretiza tiempo y amplitud}
\phantomsection\label{d2-059}
Dos decisiones distintas:

\begin{itemize}
\tightlist
\item
  muestreo: intervalos temporales \(T_s=1/f_s\);
\item
  cuantización: niveles de amplitud.
\end{itemize}

Para un rango \(V_{FS}\) y \(B\) bits:

\[\Delta_q\approx\frac{V_{FS}}{2^B}.\]

\begin{block}{Concepto clave}

Más bits reducen el escalón ideal de cuantización, pero no evitan ruido
analógico, mala calibración ni clipping.

\end{block}

\note{\textbf{Libro guía:} John Semmlow, \emph{Circuits, Signals, and
Systems for Bioengineers}, 3.ª ed., 2018.

§4.2.2, \textbf{Amplitude Slicing---Quantization} (página 182 del visor
PDF).

\textbf{Correspondencia:} El tema se desarrolla en las secciones
indicadas. La redacción es una síntesis docente.

\textbf{Microcurrículo:} semana 7, sesión 2 (sesión global 14). Los
numerales del cronograma son identificadores curriculares; no son los
apartados del libro citados arriba.}
\end{frame}

\begin{frame}{El rango dinámico exige margen}
\phantomsection\label{d2-060}
Si la señal supera el rango del ADC, aparece saturación. Si ocupa una
fracción muy pequeña, se desaprovecha resolución.

\begin{alertblock}{Atención}

El clipping es una distorsión no lineal. El filtrado no reconstruye de
manera confiable las amplitudes recortadas.

\end{alertblock}

\note{\textbf{Libro guía:} John Semmlow, \emph{Circuits, Signals, and
Systems for Bioengineers}, 3.ª ed., 2018.

§4.2.2, \textbf{Amplitude Slicing---Quantization} (página 182 del visor
PDF).

\textbf{Correspondencia:} El tema se desarrolla en las secciones
indicadas. La redacción es una síntesis docente.

\textbf{Microcurrículo:} semana 7, sesión 2 (sesión global 14). Los
numerales del cronograma son identificadores curriculares; no son los
apartados del libro citados arriba.}
\end{frame}

\begin{frame}{El reloj determina el eje temporal}
\phantomsection\label{d2-061}
Un registro requiere:

\begin{itemize}
\tightlist
\item
  frecuencia nominal de muestreo;
\item
  sellos temporales o instante inicial;
\item
  información sobre muestras perdidas;
\item
  tolerancia o deriva del reloj;
\item
  sincronización con otros equipos.
\end{itemize}

\begin{block}{Nota}

Pequeños errores de reloj pueden alterar fase, coherencia y alineación
de eventos en registros prolongados.

\end{block}

\note{\textbf{Libro guía:} John Semmlow, \emph{Circuits, Signals, and
Systems for Bioengineers}, 3.ª ed., 2018.

§1.2.3, \textbf{Signal Representation: Continuous (Analog) Signals
Versus Discrete (Digital) Signals} (página 14 del visor PDF); §4.2,
\textbf{Data Acquisition And Storage} (página 175 del visor PDF).

\textbf{Correspondencia:} Tema previsto en el cronograma, desarrollado
como ampliación docente. Las referencias del libro son antecedentes
conceptuales. Metadatos, contrato de datos y sincronización: extensión
de la cadena de adquisición para el trabajo reproducible.

\textbf{Microcurrículo:} semana 7, sesión 2 (sesión global 14). Los
numerales del cronograma son identificadores curriculares; no son los
apartados del libro citados arriba.}
\end{frame}

\begin{frame}{La sincronización sostiene análisis multicanal}
\phantomsection\label{d2-062}
Para comparar ECG, presión y respiración se necesita una referencia
temporal común.

Fuentes de desalineación:

\begin{itemize}
\tightlist
\item
  retardos distintos de sensores y filtros;
\item
  buffers de software;
\item
  marcas manuales;
\item
  relojes independientes.
\end{itemize}

\begin{block}{Concepto clave}

Un desfase instrumental puede confundirse con un retardo fisiológico.

\end{block}

\note{\textbf{Libro guía:} John Semmlow, \emph{Circuits, Signals, and
Systems for Bioengineers}, 3.ª ed., 2018.

§1.2.3, \textbf{Signal Representation: Continuous (Analog) Signals
Versus Discrete (Digital) Signals} (página 14 del visor PDF); §4.2,
\textbf{Data Acquisition And Storage} (página 175 del visor PDF).

\textbf{Correspondencia:} Tema previsto en el cronograma, desarrollado
como ampliación docente. Las referencias del libro son antecedentes
conceptuales. Metadatos, contrato de datos y sincronización: extensión
de la cadena de adquisición para el trabajo reproducible.

\textbf{Microcurrículo:} semana 7, sesión 2 (sesión global 14). Los
numerales del cronograma son identificadores curriculares; no son los
apartados del libro citados arriba.}
\end{frame}

\begin{frame}{Almacenar valores exige conservar contexto}
\phantomsection\label{d2-063}
Metadatos mínimos:

\begin{itemize}
\tightlist
\item
  unidades y escala;
\item
  frecuencia de muestreo;
\item
  nombres y ubicación de canales;
\item
  configuración de ganancia y filtros;
\item
  eventos del protocolo;
\item
  identificador seudonimizado del participante.
\end{itemize}

\begin{alertblock}{Atención}

Un archivo numérico sin metadatos puede ser ilegible desde el punto de
vista científico aunque abra correctamente.

\end{alertblock}

\note{\textbf{Libro guía:} John Semmlow, \emph{Circuits, Signals, and
Systems for Bioengineers}, 3.ª ed., 2018.

§1.2.3, \textbf{Signal Representation: Continuous (Analog) Signals
Versus Discrete (Digital) Signals} (página 14 del visor PDF); §4.2,
\textbf{Data Acquisition And Storage} (página 175 del visor PDF).

\textbf{Correspondencia:} Tema previsto en el cronograma, desarrollado
como ampliación docente. Las referencias del libro son antecedentes
conceptuales. Metadatos, contrato de datos y sincronización: extensión
de la cadena de adquisición para el trabajo reproducible.

\textbf{Microcurrículo:} semana 7, sesión 2 (sesión global 14). Los
numerales del cronograma son identificadores curriculares; no son los
apartados del libro citados arriba.}
\end{frame}

\begin{frame}{Formato tabular no siempre basta}
\phantomsection\label{d2-064}
\begin{longtable}[]{@{}ll@{}}
\toprule\noalign{}
Necesidad & Riesgo de un CSV aislado \\
\midrule\noalign{}
\endhead
tiempo irregular & se pierde si solo hay índice \\
múltiples unidades & columnas ambiguas \\
eventos & quedan en archivos separados \\
calibración & no se documenta \\
datos faltantes & pueden confundirse con cero \\
\bottomrule\noalign{}
\end{longtable}

\begin{block}{Comprobación}

El contrato de datos debe describir estructura, tipos, unidades, valores
faltantes y relación entre archivos.

\end{block}

\note{\textbf{Libro guía:} John Semmlow, \emph{Circuits, Signals, and
Systems for Bioengineers}, 3.ª ed., 2018.

§1.2.3, \textbf{Signal Representation: Continuous (Analog) Signals
Versus Discrete (Digital) Signals} (página 14 del visor PDF); §4.2,
\textbf{Data Acquisition And Storage} (página 175 del visor PDF).

\textbf{Correspondencia:} Tema previsto en el cronograma, desarrollado
como ampliación docente. Las referencias del libro son antecedentes
conceptuales. Metadatos, contrato de datos y sincronización: extensión
de la cadena de adquisición para el trabajo reproducible.

\textbf{Microcurrículo:} semana 7, sesión 2 (sesión global 14). Los
numerales del cronograma son identificadores curriculares; no son los
apartados del libro citados arriba.}
\end{frame}

\begin{frame}{Control de calidad antes del análisis}
\phantomsection\label{d2-065}
\begin{enumerate}
\tightlist
\item
  verificar longitud y tasa efectiva;
\item
  comprobar monotonicidad del tiempo;
\item
  identificar faltantes y saturación;
\item
  revisar unidades y calibración;
\item
  contrastar eventos con el protocolo;
\item
  inspeccionar canales auxiliares.
\end{enumerate}

\begin{block}{Concepto clave}

La inspección debe preceder al filtrado para evitar que el procesamiento
oculte problemas de adquisición.

\end{block}

\note{\textbf{Libro guía:} John Semmlow, \emph{Circuits, Signals, and
Systems for Bioengineers}, 3.ª ed., 2018.

§1.2.3, \textbf{Signal Representation: Continuous (Analog) Signals
Versus Discrete (Digital) Signals} (página 14 del visor PDF); §4.2,
\textbf{Data Acquisition And Storage} (página 175 del visor PDF).

\textbf{Correspondencia:} Tema previsto en el cronograma, desarrollado
como ampliación docente. Las referencias del libro son antecedentes
conceptuales. Metadatos, contrato de datos y sincronización: extensión
de la cadena de adquisición para el trabajo reproducible.

\textbf{Microcurrículo:} semana 7, sesión 2 (sesión global 14). Los
numerales del cronograma son identificadores curriculares; no son los
apartados del libro citados arriba.}
\end{frame}

\begin{frame}{Verificación de la sesión 14}
\phantomsection\label{d2-066}
Un ECG de 12 bits usa un rango de \(\pm 5\) mV y \(f_s=500\) Hz.

\begin{enumerate}
\tightlist
\item
  estime el escalón ideal de cuantización;
\item
  identifique el máximo tiempo entre muestras;
\item
  explique por qué necesita filtro antialias analógico;
\item
  enumere metadatos para interpretar derivaciones;
\item
  proponga una prueba de saturación y pérdida de muestras.
\end{enumerate}

\begin{block}{Criterio de logro}

La solución debe conectar especificaciones del sensor,
acondicionamiento, ADC y archivo final.

\end{block}

\note{\textbf{Libro guía:} John Semmlow, \emph{Circuits, Signals, and
Systems for Bioengineers}, 3.ª ed., 2018.

§1.2.2, \textbf{Biotransducers And Common Physiological Measurements}
(página 12 del visor PDF); §4.2.1, \textbf{Data Sampling: The Sampling
Theorem} (página 175 del visor PDF); §4.2.2, \textbf{Amplitude
Slicing---Quantization} (página 182 del visor PDF).

\textbf{Correspondencia:} Actividad o interpretación biomédica de
elaboración docente apoyada en las secciones indicadas.

\textbf{Microcurrículo:} semana 7, sesión 2 (sesión global 14). Los
numerales del cronograma son identificadores curriculares; no son los
apartados del libro citados arriba.}
\end{frame}

\begin{frame}{Cierre de la semana 7}
\phantomsection\label{d2-067}
\begin{itemize}
\tightlist
\item
  La duración determina el espaciado básico de frecuencias.
\item
  La ventana modifica resolución, fuga y potencia de ruido.
\item
  Zero padding interpola el eje sin mejorar la información.
\item
  La adquisición define qué contenido llega al computador.
\item
  Metadatos y sincronización forman parte de la validez.
\end{itemize}

\begin{block}{Pregunta de salida}

¿Qué problema espectral puede parecer un fenómeno fisiológico cuando dos
equipos no comparten el mismo reloj?

\end{block}

\note{\textbf{Libro guía:} John Semmlow, \emph{Circuits, Signals, and
Systems for Bioengineers}, 3.ª ed., 2018.

§4.2, \textbf{Data Acquisition And Storage} (página 175 del visor PDF).

\textbf{Correspondencia:} Síntesis o encuadre que reúne varias
secciones; no corresponde a un único apartado del libro.

\textbf{Microcurrículo:} semana 7, sesión 2 (sesión global 14). Los
numerales del cronograma son identificadores curriculares; no son los
apartados del libro citados arriba.}
\end{frame}

\begin{frame}{Semana 8 · Muestreo y estimación de potencia}
\phantomsection\label{d2-068}
\note{\textbf{Libro guía:} John Semmlow, \emph{Circuits, Signals, and
Systems for Bioengineers}, 3.ª ed., 2018.

§3.4, \textbf{Time--Frequency Transformation In The Discrete Domain}
(página 142 del visor PDF); §4.2, \textbf{Data Acquisition And Storage}
(página 175 del visor PDF); §4.6, \textbf{Time--Frequency Analysis}
(página 205 del visor PDF); §5.4, \textbf{Using The Impulse Response To
Find A Systems Output To Any Input---Convolution} (página 219 del visor
PDF).

\textbf{Correspondencia:} Síntesis o encuadre que reúne varias
secciones; no corresponde a un único apartado del libro.

\textbf{Lectura:} use estas secciones como ruta de preparación y
consulta. La bibliografía aparece al final del bloque.}
\end{frame}

\begin{frame}{Sesión 15 · Muestrear observa la señal en instantes
discretos}
\phantomsection\label{d2-069}
Un modelo ideal usa un tren de impulsos:

\[p(t)=\sum_{n=-\infty}^{\infty}\delta(t-nT_s).\]

La señal muestreada es

\[x_s(t)=x(t)p(t)=\sum_n x(nT_s)\delta(t-nT_s).\]

\begin{block}{Concepto clave}

El muestreo conserva valores en instantes separados por \(T_s\). Entre
muestras no existe observación directa.

\end{block}

\note{\textbf{Libro guía:} John Semmlow, \emph{Circuits, Signals, and
Systems for Bioengineers}, 3.ª ed., 2018.

§4.2.1, \textbf{Data Sampling: The Sampling Theorem} (página 175 del
visor PDF); §5.6.1, \textbf{Data Sampling Revisited} (página 238 del
visor PDF).

\textbf{Correspondencia:} El tema se desarrolla en las secciones
indicadas. La redacción es una síntesis docente.

\textbf{Microcurrículo:} semana 8, sesión 1 (sesión global 15). Los
numerales del cronograma son identificadores curriculares; no son los
apartados del libro citados arriba.}
\end{frame}

\begin{frame}{Muestrear replica el espectro}
\phantomsection\label{d2-070}
El espectro del tren de impulsos también es periódico:

\[P(f)=f_s\sum_{k=-\infty}^{\infty}\delta(f-kf_s).\]

Por tanto:

\[X_s(f)=f_s\sum_{k=-\infty}^{\infty}X(f-kf_s).\]

\begin{block}{Concepto clave}

El muestreo crea réplicas del espectro separadas por \(f_s\). El
aliasing aparece cuando esas réplicas se superponen.

\end{block}

\note{\textbf{Libro guía:} John Semmlow, \emph{Circuits, Signals, and
Systems for Bioengineers}, 3.ª ed., 2018.

§5.6.1, \textbf{Data Sampling Revisited} (página 238 del visor PDF).

\textbf{Correspondencia:} El tema se desarrolla en las secciones
indicadas. La redacción es una síntesis docente.

\textbf{Microcurrículo:} semana 8, sesión 1 (sesión global 15). Los
numerales del cronograma son identificadores curriculares; no son los
apartados del libro citados arriba.}
\end{frame}

\begin{frame}{Nyquist establece una condición ideal}
\phantomsection\label{d2-071}
Si \(X(f)=0\) para \(|f|>B\), la reconstrucción ideal requiere

\[f_s>2B.\]

La frecuencia \(f_s/2\) se denomina frecuencia de Nyquist.

\begin{alertblock}{Atención}

La condición supone banda estrictamente limitada y un filtro ideal. Las
señales y filtros físicos requieren margen de transición.

\end{alertblock}

\note{\textbf{Libro guía:} John Semmlow, \emph{Circuits, Signals, and
Systems for Bioengineers}, 3.ª ed., 2018.

§4.2.1, \textbf{Data Sampling: The Sampling Theorem} (página 175 del
visor PDF).

\textbf{Correspondencia:} El tema se desarrolla en las secciones
indicadas. La redacción es una síntesis docente.

\textbf{Microcurrículo:} semana 8, sesión 1 (sesión global 15). Los
numerales del cronograma son identificadores curriculares; no son los
apartados del libro citados arriba.}
\end{frame}

\begin{frame}{La frecuencia de Nyquist no es la tasa de Nyquist}
\phantomsection\label{d2-072}
\begin{itemize}
\tightlist
\item
  frecuencia de Nyquist: \(f_s/2\) para una tasa dada;
\item
  tasa de Nyquist: \(2B\), límite teórico asociado a una banda \(B\);
  para el diseño adoptamos \(f_s>2B\) y margen de transición.
\end{itemize}

\begin{block}{Concepto clave}

Intercambiar estos términos produce especificaciones ambiguas. Debe
indicarse si se habla de la señal o del sistema de adquisición.

\end{block}

\note{\textbf{Libro guía:} John Semmlow, \emph{Circuits, Signals, and
Systems for Bioengineers}, 3.ª ed., 2018.

§4.2.1, \textbf{Data Sampling: The Sampling Theorem} (página 175 del
visor PDF).

\textbf{Correspondencia:} El tema se desarrolla en las secciones
indicadas. La redacción es una síntesis docente.

\textbf{Microcurrículo:} semana 8, sesión 1 (sesión global 15). Los
numerales del cronograma son identificadores curriculares; no son los
apartados del libro citados arriba.}
\end{frame}

\begin{frame}{El alias de una frecuencia depende de \(f_s\)}
\phantomsection\label{d2-073}
Una frecuencia \(f_0\) se pliega dentro de \([0,f_s/2]\). Una forma
operativa es

\[f_a=\left|\left((f_0+f_s/2)\bmod f_s\right)-f_s/2\right|.\]

Ejemplo: con \(f_s=100\) Hz, una componente de 70 Hz aparece a 30 Hz.

\begin{alertblock}{Atención}

El alias conserva una apariencia sinusoidal válida. El registro aislado
no permite saber si provino de 30 Hz o de 70 Hz.

\end{alertblock}

\note{\textbf{Libro guía:} John Semmlow, \emph{Circuits, Signals, and
Systems for Bioengineers}, 3.ª ed., 2018.

§4.2.1, \textbf{Data Sampling: The Sampling Theorem} (página 175 del
visor PDF).

\textbf{Correspondencia:} El tema se desarrolla en las secciones
indicadas. La redacción es una síntesis docente.

\textbf{Microcurrículo:} semana 8, sesión 1 (sesión global 15). Los
numerales del cronograma son identificadores curriculares; no son los
apartados del libro citados arriba.}
\end{frame}

\begin{frame}{El filtro antialias necesita una banda de transición}
\phantomsection\label{d2-074}
Especificación conceptual:

\begin{itemize}
\tightlist
\item
  banda útil hasta \(f_p\);
\item
  banda de rechazo desde \(f_a\);
\item
  borde digital máximo \(f_s/2\);
\item
  atenuación suficiente antes de Nyquist.
\end{itemize}

\begin{block}{Concepto clave}

Elegir \(f_s\) implica reservar espacio para que el filtro analógico
pase de la banda útil a la banda de rechazo.

\end{block}

\note{\textbf{Libro guía:} John Semmlow, \emph{Circuits, Signals, and
Systems for Bioengineers}, 3.ª ed., 2018.

§4.2.1, \textbf{Data Sampling: The Sampling Theorem} (página 175 del
visor PDF).

\textbf{Correspondencia:} El tema se desarrolla en las secciones
indicadas. La redacción es una síntesis docente.

\textbf{Microcurrículo:} semana 8, sesión 1 (sesión global 15). Los
numerales del cronograma son identificadores curriculares; no son los
apartados del libro citados arriba.}
\end{frame}

\begin{frame}{Sobremuestrear simplifica parte del diseño}
\phantomsection\label{d2-075}
Una tasa mayor puede:

\begin{itemize}
\tightlist
\item
  ampliar la transición del filtro antialias;
\item
  mejorar la localización temporal;
\item
  reducir sensibilidad a ciertos errores analógicos.
\end{itemize}

También aumenta datos, consumo y costo de transmisión.

\begin{block}{Nota}

Sobremuestrear no aumenta por sí mismo la duración ni la resolución
frecuencial de un registro fijo en segundos.

\end{block}

\note{\textbf{Libro guía:} John Semmlow, \emph{Circuits, Signals, and
Systems for Bioengineers}, 3.ª ed., 2018.

§4.2.1, \textbf{Data Sampling: The Sampling Theorem} (página 175 del
visor PDF).

\textbf{Correspondencia:} El tema se desarrolla en las secciones
indicadas. La redacción es una síntesis docente.

\textbf{Microcurrículo:} semana 8, sesión 1 (sesión global 15). Los
numerales del cronograma son identificadores curriculares; no son los
apartados del libro citados arriba.}
\end{frame}

\begin{frame}{La resolución temporal básica es \(T_s\)}
\phantomsection\label{d2-076}
\[T_s=\frac{1}{f_s}.\]

Si \(f_s=500\) Hz, la separación nominal entre muestras es 2 ms.

\begin{alertblock}{Atención}

La precisión de un evento no queda determinada solo por \(T_s\).
Interpolación, ruido, ancho de banda y algoritmo también influyen.

\end{alertblock}

\note{\textbf{Libro guía:} John Semmlow, \emph{Circuits, Signals, and
Systems for Bioengineers}, 3.ª ed., 2018.

§4.2.1, \textbf{Data Sampling: The Sampling Theorem} (página 175 del
visor PDF).

\textbf{Correspondencia:} El tema se desarrolla en las secciones
indicadas. La redacción es una síntesis docente.

\textbf{Microcurrículo:} semana 8, sesión 1 (sesión global 15). Los
numerales del cronograma son identificadores curriculares; no son los
apartados del libro citados arriba.}
\end{frame}

\begin{frame}{Duración y muestreo controlan aspectos distintos}
\phantomsection\label{d2-077}
\begin{longtable}[]{@{}ll@{}}
\toprule\noalign{}
Parámetro & Relación principal \\
\midrule\noalign{}
\endhead
\(f_s\) & rango frecuencial representable y paso temporal \\
\(T\) & resolución frecuencial básica y contexto observado \\
\(N=f_sT\) & tamaño del registro y costo computacional \\
\bottomrule\noalign{}
\end{longtable}

\begin{block}{Concepto clave}

Una mayor cantidad de muestras puede provenir de mayor duración o mayor
tasa. Esas decisiones no son equivalentes.

\end{block}

\note{\textbf{Libro guía:} John Semmlow, \emph{Circuits, Signals, and
Systems for Bioengineers}, 3.ª ed., 2018.

§4.2.1, \textbf{Data Sampling: The Sampling Theorem} (página 175 del
visor PDF).

\textbf{Correspondencia:} El tema se desarrolla en las secciones
indicadas. La redacción es una síntesis docente.

\textbf{Microcurrículo:} semana 8, sesión 1 (sesión global 15). Los
numerales del cronograma son identificadores curriculares; no son los
apartados del libro citados arriba.}
\end{frame}

\begin{frame}{Muestreo irregular exige otro modelo}
\phantomsection\label{d2-078}
Si los instantes son \(t_n\) y no cumplen \(t_n=nT_s\), la DFT estándar
no representa exactamente el eje temporal real.

Opciones:

\begin{itemize}
\tightlist
\item
  corregir o interpolar con justificación;
\item
  usar métodos para muestreo no uniforme;
\item
  modelar explícitamente los sellos temporales.
\end{itemize}

\begin{alertblock}{Atención}

Asignar una frecuencia nominal a datos irregulares puede introducir
errores de fase y frecuencia.

\end{alertblock}

\note{\textbf{Libro guía:} John Semmlow, \emph{Circuits, Signals, and
Systems for Bioengineers}, 3.ª ed., 2018.

§4.2.1, \textbf{Data Sampling: The Sampling Theorem} (página 175 del
visor PDF).

\textbf{Correspondencia:} Las secciones aportan el fundamento. La
diapositiva amplía la formulación o precisa condiciones que no se
presentan con idéntica notación en el libro. El muestreo uniforme es el
modelo del teorema; datos irregulares requieren otro modelo o un
remuestreo justificado.

\textbf{Microcurrículo:} semana 8, sesión 1 (sesión global 15). Los
numerales del cronograma son identificadores curriculares; no son los
apartados del libro citados arriba.}
\end{frame}

\begin{frame}[fragile]{Código: observar aliasing}
\phantomsection\label{d2-079}
\begin{Shaded}
\begin{Highlighting}[]
\ImportTok{import}\NormalTok{ numpy }\ImportTok{as}\NormalTok{ np}

\NormalTok{fs }\OperatorTok{=} \FloatTok{100.0}
\NormalTok{t }\OperatorTok{=}\NormalTok{ np.arange(}\DecValTok{0}\NormalTok{, }\DecValTok{1}\NormalTok{, }\DecValTok{1}\OperatorTok{/}\NormalTok{fs)}
\NormalTok{x30 }\OperatorTok{=}\NormalTok{ np.cos(}\DecValTok{2}\OperatorTok{*}\NormalTok{np.pi}\OperatorTok{*}\DecValTok{30}\OperatorTok{*}\NormalTok{t)}
\NormalTok{x70 }\OperatorTok{=}\NormalTok{ np.cos(}\DecValTok{2}\OperatorTok{*}\NormalTok{np.pi}\OperatorTok{*}\DecValTok{70}\OperatorTok{*}\NormalTok{t)}

\NormalTok{error }\OperatorTok{=}\NormalTok{ np.}\BuiltInTok{max}\NormalTok{(np.}\BuiltInTok{abs}\NormalTok{(x30 }\OperatorTok{{-}}\NormalTok{ x70))}
\end{Highlighting}
\end{Shaded}

\begin{block}{Nota}

Con esos instantes, ambas secuencias coinciden salvo error numérico. El
ejemplo demuestra indistinguibilidad discreta, no igualdad de las
señales continuas.

\end{block}

\note{\textbf{Libro guía:} John Semmlow, \emph{Circuits, Signals, and
Systems for Bioengineers}, 3.ª ed., 2018.

§4.2.1, \textbf{Data Sampling: The Sampling Theorem} (página 175 del
visor PDF).

\textbf{Correspondencia:} Actividad o interpretación biomédica de
elaboración docente apoyada en las secciones indicadas.

\textbf{Microcurrículo:} semana 8, sesión 1 (sesión global 15). Los
numerales del cronograma son identificadores curriculares; no son los
apartados del libro citados arriba.}
\end{frame}

\begin{frame}{Selección de \(f_s\) para una bioseñal}
\phantomsection\label{d2-080}
Procedimiento:

\begin{enumerate}
\tightlist
\item
  definir la banda fisiológica necesaria para la tarea;
\item
  incluir transitorios que deban conservarse;
\item
  establecer margen para el filtro antialias;
\item
  revisar precisión temporal y sincronización;
\item
  evaluar almacenamiento y transmisión;
\item
  documentar la tasa efectiva.
\end{enumerate}

\begin{block}{Comprobación}

La tasa se justifica desde la pregunta y la cadena completa, no desde
una regla fija por modalidad.

\end{block}

\note{\textbf{Libro guía:} John Semmlow, \emph{Circuits, Signals, and
Systems for Bioengineers}, 3.ª ed., 2018.

§4.2.1, \textbf{Data Sampling: The Sampling Theorem} (página 175 del
visor PDF).

\textbf{Correspondencia:} Actividad o interpretación biomédica de
elaboración docente apoyada en las secciones indicadas.

\textbf{Microcurrículo:} semana 8, sesión 1 (sesión global 15). Los
numerales del cronograma son identificadores curriculares; no son los
apartados del libro citados arriba.}
\end{frame}

\begin{frame}{Verificación de la sesión 15}
\phantomsection\label{d2-081}
Una señal contiene información hasta 120 Hz y se muestrea a 250 Hz.

\begin{enumerate}
\tightlist
\item
  calcule la frecuencia de Nyquist;
\item
  evalúe qué atenuación, orden y tolerancia se necesitarían en una
  transición de 5 Hz;
\item
  determine el alias de una interferencia de 180 Hz;
\item
  calcule \(T_s\);
\item
  proponga una tasa razonada, sin asumir un valor universal.
\end{enumerate}

\begin{block}{Criterio de logro}

La propuesta debe incluir banda de transición, atenuación y precisión
temporal.

\end{block}

\note{\textbf{Libro guía:} John Semmlow, \emph{Circuits, Signals, and
Systems for Bioengineers}, 3.ª ed., 2018.

§4.2.1, \textbf{Data Sampling: The Sampling Theorem} (página 175 del
visor PDF).

\textbf{Correspondencia:} Actividad o interpretación biomédica de
elaboración docente apoyada en las secciones indicadas.

\textbf{Microcurrículo:} semana 8, sesión 1 (sesión global 15). Los
numerales del cronograma son identificadores curriculares; no son los
apartados del libro citados arriba.}
\end{frame}

\begin{frame}{Sesión 16 · Potencia espectral describe distribución
cuadrática}
\phantomsection\label{d2-082}
Para una señal de potencia, la PSD \(S_{xx}(f)\) distribuye potencia por
unidad de frecuencia.

\[P=\int_{-\infty}^{\infty}S_{xx}(f)\,df.\]

\begin{block}{Concepto clave}

La PSD permite integrar potencia en bandas, comparar ritmos y
cuantificar cómo se distribuye la variabilidad.

\end{block}

\note{\textbf{Libro guía:} John Semmlow, \emph{Circuits, Signals, and
Systems for Bioengineers}, 3.ª ed., 2018.

§4.3, \textbf{Power Spectrum} (página 191 del visor PDF).

\textbf{Correspondencia:} El tema se desarrolla en las secciones
indicadas. La redacción es una síntesis docente.

\textbf{Microcurrículo:} semana 8, sesión 2 (sesión global 16). Los
numerales del cronograma son identificadores curriculares; no son los
apartados del libro citados arriba.}
\end{frame}

\begin{frame}{Energía y potencia no son intercambiables}
\phantomsection\label{d2-083}
Energía discreta:

\[E_x=\sum_{n=-\infty}^{\infty}|x[n]|^2.\]

Potencia media:

\[P_x=\lim_{N\to\infty}\frac{1}{2N+1}\sum_{n=-N}^{N}|x[n]|^2.\]

\begin{alertblock}{Atención}

Un registro finito siempre tiene suma cuadrática finita, pero su modelo
fisiológico puede tratarse como realización de un proceso de potencia.

\end{alertblock}

\note{\textbf{Libro guía:} John Semmlow, \emph{Circuits, Signals, and
Systems for Bioengineers}, 3.ª ed., 2018.

§2.2.1, \textbf{Mean And Root Mean Squared Signal Measurements} (página
58 del visor PDF); §4.3, \textbf{Power Spectrum} (página 191 del visor
PDF).

\textbf{Correspondencia:} Las secciones aportan el fundamento. La
diapositiva amplía la formulación o precisa condiciones que no se
presentan con idéntica notación en el libro.

\textbf{Microcurrículo:} semana 8, sesión 2 (sesión global 16). Los
numerales del cronograma son identificadores curriculares; no son los
apartados del libro citados arriba.}
\end{frame}

\begin{frame}{Espectro de energía y PSD responden a modelos distintos}
\phantomsection\label{d2-084}
\begin{itemize}
\tightlist
\item
  \(|X(f)|^2\) describe densidad de energía para señales de energía.
\item
  \(S_{xx}(f)\) describe densidad de potencia para procesos o señales
  persistentes.
\item
  Un registro finito permite estimar estas cantidades, no observarlas
  sin incertidumbre.
\end{itemize}

\begin{block}{Concepto clave}

El periodograma es un estimador de PSD. No es la PSD verdadera del
proceso.

\end{block}

\note{\textbf{Libro guía:} John Semmlow, \emph{Circuits, Signals, and
Systems for Bioengineers}, 3.ª ed., 2018.

§2.2.1, \textbf{Mean And Root Mean Squared Signal Measurements} (página
58 del visor PDF); §4.3, \textbf{Power Spectrum} (página 191 del visor
PDF).

\textbf{Correspondencia:} Las secciones aportan el fundamento. La
diapositiva amplía la formulación o precisa condiciones que no se
presentan con idéntica notación en el libro.

\textbf{Microcurrículo:} semana 8, sesión 2 (sesión global 16). Los
numerales del cronograma son identificadores curriculares; no son los
apartados del libro citados arriba.}
\end{frame}

\begin{frame}{El periodograma parte de la DFT}
\phantomsection\label{d2-085}
Una forma bilateral para \(N\) muestras y tasa \(f_s\) es

\[\widehat S_{xx}[k]=\frac{1}{Nf_s}|X[k]|^2.\]

La constante cambia si se usa ventana, espectro unilateral o una
convención diferente.

\begin{alertblock}{Atención}

No se puede comparar PSD entre programas si no se verifican
escalamiento, ventana, detrend y lados del espectro.

\end{alertblock}

\note{\textbf{Libro guía:} John Semmlow, \emph{Circuits, Signals, and
Systems for Bioengineers}, 3.ª ed., 2018.

§3.4.1, \textbf{The Discrete Fourier Transform} (página 142 del visor
PDF); §4.3, \textbf{Power Spectrum} (página 191 del visor PDF).

\textbf{Correspondencia:} El tema se desarrolla en las secciones
indicadas. La redacción es una síntesis docente.

\textbf{Microcurrículo:} semana 8, sesión 2 (sesión global 16). Los
numerales del cronograma son identificadores curriculares; no son los
apartados del libro citados arriba.}
\end{frame}

\begin{frame}{Las unidades distinguen espectro y densidad}
\phantomsection\label{d2-086}
Si \(x\) se mide en mV:

\begin{itemize}
\tightlist
\item
  amplitud espectral: mV;
\item
  potencia por bin: mV\(^2\);
\item
  PSD: mV\(^2\)/Hz.
\end{itemize}

\begin{block}{Concepto clave}

La integración de PSD sobre Hz debe recuperar una potencia en unidades
cuadráticas de la señal.

\end{block}

\note{\textbf{Libro guía:} John Semmlow, \emph{Circuits, Signals, and
Systems for Bioengineers}, 3.ª ed., 2018.

§4.3, \textbf{Power Spectrum} (página 191 del visor PDF).

\textbf{Correspondencia:} Las secciones aportan el fundamento. La
diapositiva amplía la formulación o precisa condiciones que no se
presentan con idéntica notación en el libro.

\textbf{Microcurrículo:} semana 8, sesión 2 (sesión global 16). Los
numerales del cronograma son identificadores curriculares; no son los
apartados del libro citados arriba.}
\end{frame}

\begin{frame}{La PSD unilateral conserva potencia}
\phantomsection\label{d2-087}
Para señales reales se combinan contribuciones positivas y negativas.

\begin{itemize}
\tightlist
\item
  DC no se duplica.
\item
  Nyquist no se duplica cuando existe.
\item
  Los bins interiores se duplican.
\end{itemize}

\begin{alertblock}{Atención}

La regla de duplicación depende de si se trabaja con amplitud, potencia
o densidad. Debe aplicarse en la escala correspondiente.

\end{alertblock}

\note{\textbf{Libro guía:} John Semmlow, \emph{Circuits, Signals, and
Systems for Bioengineers}, 3.ª ed., 2018.

§4.3, \textbf{Power Spectrum} (página 191 del visor PDF).

\textbf{Correspondencia:} Las secciones aportan el fundamento. La
diapositiva amplía la formulación o precisa condiciones que no se
presentan con idéntica notación en el libro.

\textbf{Microcurrículo:} semana 8, sesión 2 (sesión global 16). Los
numerales del cronograma son identificadores curriculares; no son los
apartados del libro citados arriba.}
\end{frame}

\begin{frame}{Parseval ofrece una prueba de consistencia}
\phantomsection\label{d2-088}
Para la DFT adoptada:

\[\sum_{n=0}^{N-1}|x[n]|^2=\frac{1}{N}\sum_{k=0}^{N-1}|X[k]|^2.\]

\begin{block}{Comprobación}

Con ventana rectangular y sin quitar media,
\(\Delta f\sum_k\widehat S[k]=N^{-1}\sum_n|x[n]|^2\). Una ventana
distinta produce una potencia ponderada: la referencia temporal debe
usar la misma ventana y el mismo centrado.

\end{block}

\note{\textbf{Libro guía:} John Semmlow, \emph{Circuits, Signals, and
Systems for Bioengineers}, 3.ª ed., 2018.

§4.3, \textbf{Power Spectrum} (página 191 del visor PDF).

\textbf{Correspondencia:} Las secciones aportan el fundamento. La
diapositiva amplía la formulación o precisa condiciones que no se
presentan con idéntica notación en el libro.

\textbf{Microcurrículo:} semana 8, sesión 2 (sesión global 16). Los
numerales del cronograma son identificadores curriculares; no son los
apartados del libro citados arriba.}
\end{frame}

\begin{frame}[fragile]{Parseval con ventana: comparación correcta}
\phantomsection\label{d2-089}
\begin{Shaded}
\begin{Highlighting}[]
\ImportTok{from}\NormalTok{ scipy.signal }\ImportTok{import}\NormalTok{ get\_window, periodogram}
\NormalTok{x0 }\OperatorTok{=}\NormalTok{ np.asarray(x) }\OperatorTok{{-}}\NormalTok{ np.mean(x)}
\NormalTok{w }\OperatorTok{=}\NormalTok{ get\_window(}\StringTok{"hann"}\NormalTok{, }\BuiltInTok{len}\NormalTok{(x0), fftbins}\OperatorTok{=}\VariableTok{True}\NormalTok{)}
\NormalTok{f, S }\OperatorTok{=}\NormalTok{ periodogram(x0, fs}\OperatorTok{=}\NormalTok{fs, window}\OperatorTok{=}\NormalTok{w,}
\NormalTok{                   detrend}\OperatorTok{=}\VariableTok{False}\NormalTok{, scaling}\OperatorTok{=}\StringTok{"density"}\NormalTok{)}
\NormalTok{P\_f }\OperatorTok{=}\NormalTok{ (f[}\DecValTok{1}\NormalTok{]}\OperatorTok{{-}}\NormalTok{f[}\DecValTok{0}\NormalTok{])}\OperatorTok{*}\NormalTok{np.}\BuiltInTok{sum}\NormalTok{(S)}
\NormalTok{P\_t }\OperatorTok{=}\NormalTok{ np.}\BuiltInTok{sum}\NormalTok{((w}\OperatorTok{*}\NormalTok{x0)}\OperatorTok{**}\DecValTok{2}\NormalTok{)}\OperatorTok{/}\NormalTok{np.}\BuiltInTok{sum}\NormalTok{(w}\OperatorTok{**}\DecValTok{2}\NormalTok{)}
\ControlFlowTok{assert}\NormalTok{ np.isclose(P\_f, P\_t)}
\end{Highlighting}
\end{Shaded}

\begin{block}{Concepto clave}

La prueba usa el mismo segmento, ventana y centrado. La identidad
compara \(P_f\) con \(P_t\) ponderada, no con la media cuadrática sin
ventana.

\end{block}

\note{\textbf{Libro guía:} John Semmlow, \emph{Circuits, Signals, and
Systems for Bioengineers}, 3.ª ed., 2018.

§4.3, \textbf{Power Spectrum} (página 191 del visor PDF); §4.4,
\textbf{Spectral Averaging} (página 195 del visor PDF).

\textbf{Correspondencia:} Ejemplo, figura o implementación de
elaboración docente a partir de estos fundamentos. No es una
transcripción del código MATLAB del libro. Escalamiento contrastado con
la documentación oficial de scipy.signal.periodogram.

\textbf{Microcurrículo:} semana 8, sesión 2 (sesión global 16). Los
numerales del cronograma son identificadores curriculares; no son los
apartados del libro citados arriba.

\textbf{Implementación:}
\href{https://docs.scipy.org/doc/scipy/reference/generated/scipy.signal.periodogram.html}{SciPy:
periodogram y escalamiento espectral}.}
\end{frame}

\begin{frame}{Retirar la media cambia la PSD cerca de 0 Hz}
\phantomsection\label{d2-090}
La componente DC representa el promedio del segmento.

\[x_0[n]=x[n]-\bar{x}.\]

\begin{alertblock}{Atención}

Retirar media o tendencia debe declararse. Puede ser conveniente para
estudiar oscilaciones, pero elimina información de nivel.

\end{alertblock}

\note{\textbf{Libro guía:} John Semmlow, \emph{Circuits, Signals, and
Systems for Bioengineers}, 3.ª ed., 2018.

§2.2.1, \textbf{Mean And Root Mean Squared Signal Measurements} (página
58 del visor PDF); §4.3, \textbf{Power Spectrum} (página 191 del visor
PDF).

\textbf{Correspondencia:} El tema se desarrolla en las secciones
indicadas. La redacción es una síntesis docente.

\textbf{Microcurrículo:} semana 8, sesión 2 (sesión global 16). Los
numerales del cronograma son identificadores curriculares; no son los
apartados del libro citados arriba.}
\end{frame}

\begin{frame}{Autocorrelación y PSD: el vínculo estadístico}
\phantomsection\label{d2-091}
Para un proceso estacionario en sentido amplio:

\[R_{xx}(\tau)=\mathbb E\{x(t+\tau)x^*(t)\},\qquad
S_{xx}(f)=\mathcal F\{R_{xx}(\tau)\}.\]

Si la media es \(\mu\), \(C_{xx}(\tau)=R_{xx}(\tau)-|\mu|^2\). Su
transformada elimina la contribución DC \(|\mu|^2\delta(f)\).

\begin{block}{Nota}

La autocorrelación de un registro finito es un estimador. Tiempo de
observación, normalización y estacionariedad determinan cómo se
relaciona con la PSD estimada.

\end{block}

\note{\textbf{Libro guía:} John Semmlow, \emph{Circuits, Signals, and
Systems for Bioengineers}, 3.ª ed., 2018.

§2.4.5, \textbf{Autocorrelation} (página 104 del visor PDF); §2.4.6,
\textbf{Autocovariance And Cross-Covariance} (página 108 del visor PDF);
§4.3, \textbf{Power Spectrum} (página 191 del visor PDF).

\textbf{Correspondencia:} Las secciones aportan el fundamento. La
diapositiva amplía la formulación o precisa condiciones que no se
presentan con idéntica notación en el libro. Formalización estadística
del vínculo entre correlación y potencia; no confundir una estimación
finita con el parámetro del proceso.

\textbf{Microcurrículo:} semana 8, sesión 2 (sesión global 16). Los
numerales del cronograma son identificadores curriculares; no son los
apartados del libro citados arriba.}
\end{frame}

\begin{frame}{El periodograma tiene alta variabilidad}
\phantomsection\label{d2-092}
Aumentar \(N\) mejora el espaciado frecuencial, pero el periodograma
simple no se vuelve necesariamente estable como estimador puntual.

\begin{block}{Concepto clave}

La varianza del estimador motiva el promediado de segmentos y el método
de Welch.

\end{block}

\note{\textbf{Libro guía:} John Semmlow, \emph{Circuits, Signals, and
Systems for Bioengineers}, 3.ª ed., 2018.

§4.4, \textbf{Spectral Averaging} (página 195 del visor PDF).

\textbf{Correspondencia:} El tema se desarrolla en las secciones
indicadas. La redacción es una síntesis docente.

\textbf{Microcurrículo:} semana 8, sesión 2 (sesión global 16). Los
numerales del cronograma son identificadores curriculares; no son los
apartados del libro citados arriba.}
\end{frame}

\begin{frame}{Potencia en una banda}
\phantomsection\label{d2-093}
Para una banda \([f_1,f_2]\):

\[P_{banda}=\int_{f_1}^{f_2}S_{xx}(f)\,df.\]

En un eje discreto uniforme se aproxima mediante suma o integración
trapezoidal.

\begin{alertblock}{Atención}

Las bandas deben justificarse para la modalidad, población, estado y
método de adquisición.

\end{alertblock}

\note{\textbf{Libro guía:} John Semmlow, \emph{Circuits, Signals, and
Systems for Bioengineers}, 3.ª ed., 2018.

§4.3, \textbf{Power Spectrum} (página 191 del visor PDF); §4.5,
\textbf{Signal Bandwidth} (página 200 del visor PDF).

\textbf{Correspondencia:} El tema se desarrolla en las secciones
indicadas. La redacción es una síntesis docente.

\textbf{Microcurrículo:} semana 8, sesión 2 (sesión global 16). Los
numerales del cronograma son identificadores curriculares; no son los
apartados del libro citados arriba.}
\end{frame}

\begin{frame}{Potencia absoluta y relativa}
\phantomsection\label{d2-094}
\[P_{rel}=\frac{P_{banda}}{P_{referencia}}.\]

La potencia relativa reduce algunas diferencias de escala, pero acopla
el numerador al denominador.

\begin{block}{Nota}

Un aumento de potencia relativa puede deberse a aumento de la banda o
disminución de la potencia de referencia.

\end{block}

\note{\textbf{Libro guía:} John Semmlow, \emph{Circuits, Signals, and
Systems for Bioengineers}, 3.ª ed., 2018.

§4.3, \textbf{Power Spectrum} (página 191 del visor PDF).

\textbf{Correspondencia:} Las secciones aportan el fundamento. La
diapositiva amplía la formulación o precisa condiciones que no se
presentan con idéntica notación en el libro.

\textbf{Microcurrículo:} semana 8, sesión 2 (sesión global 16). Los
numerales del cronograma son identificadores curriculares; no son los
apartados del libro citados arriba.}
\end{frame}

\begin{frame}[fragile]{Código: periodograma con unidades explícitas}
\phantomsection\label{d2-095}
\begin{Shaded}
\begin{Highlighting}[]
\ImportTok{from}\NormalTok{ scipy.signal }\ImportTok{import}\NormalTok{ periodogram}

\NormalTok{f, Pxx }\OperatorTok{=}\NormalTok{ periodogram(}
\NormalTok{    x,}
\NormalTok{    fs}\OperatorTok{=}\NormalTok{fs,}
\NormalTok{    window}\OperatorTok{=}\StringTok{"hann"}\NormalTok{,}
\NormalTok{    detrend}\OperatorTok{=}\StringTok{"constant"}\NormalTok{,}
\NormalTok{    scaling}\OperatorTok{=}\StringTok{"density"}\NormalTok{,}
\NormalTok{    return\_onesided}\OperatorTok{=}\VariableTok{True}\NormalTok{,}
\NormalTok{)}
\NormalTok{potencia }\OperatorTok{=}\NormalTok{ (f[}\DecValTok{1}\NormalTok{] }\OperatorTok{{-}}\NormalTok{ f[}\DecValTok{0}\NormalTok{]) }\OperatorTok{*}\NormalTok{ Pxx.}\BuiltInTok{sum}\NormalTok{()}
\end{Highlighting}
\end{Shaded}

\begin{block}{Nota}

La suma de bins por su ancho incluye completamente DC y Nyquist. Con
Hann integra potencia del registro centrado y ponderado por la ventana,
dividida por su energía. No tiene por qué coincidir exactamente con la
potencia sin ponderar.

\end{block}

\note{\textbf{Libro guía:} John Semmlow, \emph{Circuits, Signals, and
Systems for Bioengineers}, 3.ª ed., 2018.

§4.3, \textbf{Power Spectrum} (página 191 del visor PDF).

\textbf{Correspondencia:} Actividad o interpretación biomédica de
elaboración docente apoyada en las secciones indicadas.

\textbf{Microcurrículo:} semana 8, sesión 2 (sesión global 16). Los
numerales del cronograma son identificadores curriculares; no son los
apartados del libro citados arriba.

\textbf{Implementación:}
\href{https://docs.scipy.org/doc/scipy/reference/generated/scipy.signal.periodogram.html}{SciPy:
periodogram y escalamiento espectral}.}
\end{frame}

\begin{frame}{Lectura biomédica prudente}
\phantomsection\label{d2-096}
Antes de interpretar un pico o una banda:

\begin{itemize}
\tightlist
\item
  verificar estado fisiológico y protocolo;
\item
  revisar artefactos y canales auxiliares;
\item
  confirmar duración y resolución;
\item
  comprobar preprocesamiento;
\item
  estimar incertidumbre o variabilidad.
\end{itemize}

\begin{block}{Concepto clave}

Una asociación entre banda y fenómeno requiere evidencia adicional. La
frecuencia por sí sola no identifica la fuente.

\end{block}

\note{\textbf{Libro guía:} John Semmlow, \emph{Circuits, Signals, and
Systems for Bioengineers}, 3.ª ed., 2018.

§4.3, \textbf{Power Spectrum} (página 191 del visor PDF).

\textbf{Correspondencia:} Actividad o interpretación biomédica de
elaboración docente apoyada en las secciones indicadas.

\textbf{Microcurrículo:} semana 8, sesión 2 (sesión global 16). Los
numerales del cronograma son identificadores curriculares; no son los
apartados del libro citados arriba.}
\end{frame}

\begin{frame}[fragile]{Verificación de la sesión 16}
\phantomsection\label{d2-097}
Una PSD de aceleración se expresa en \((m/s^2)^2/Hz\).

\begin{enumerate}
\tightlist
\item
  indique las unidades de su integral entre 0 y 20 Hz;
\item
  explique qué hace \texttt{detrend="constant"};
\item
  compare periodograma unilateral y bilateral;
\item
  proponga una prueba de Parseval;
\item
  explique por qué un pico no identifica automáticamente un mecanismo.
\end{enumerate}

\begin{block}{Criterio de logro}

La respuesta debe conservar unidades y diferenciar estimador, parámetro
y explicación fisiológica.

\end{block}

\note{\textbf{Libro guía:} John Semmlow, \emph{Circuits, Signals, and
Systems for Bioengineers}, 3.ª ed., 2018.

§4.3, \textbf{Power Spectrum} (página 191 del visor PDF).

\textbf{Correspondencia:} Actividad o interpretación biomédica de
elaboración docente apoyada en las secciones indicadas.

\textbf{Microcurrículo:} semana 8, sesión 2 (sesión global 16). Los
numerales del cronograma son identificadores curriculares; no son los
apartados del libro citados arriba.}
\end{frame}

\begin{frame}{Cierre de la semana 8}
\phantomsection\label{d2-098}
\begin{itemize}
\tightlist
\item
  El muestreo replica el espectro y puede producir aliasing.
\item
  Nyquist ofrece una condición ideal, no un margen de diseño suficiente.
\item
  \(f_s\) y \(T\) controlan propiedades distintas.
\item
  El periodograma estima densidad de potencia a partir de un bloque
  finito.
\item
  Las unidades y Parseval permiten verificar el escalamiento.
\end{itemize}

\begin{block}{Pregunta de salida}

¿Por qué duplicar la frecuencia de muestreo sin cambiar la duración no
divide a la mitad la incertidumbre del periodograma?

\end{block}

\note{\textbf{Libro guía:} John Semmlow, \emph{Circuits, Signals, and
Systems for Bioengineers}, 3.ª ed., 2018.

§4.2.1, \textbf{Data Sampling: The Sampling Theorem} (página 175 del
visor PDF); §4.3, \textbf{Power Spectrum} (página 191 del visor PDF).

\textbf{Correspondencia:} Síntesis o encuadre que reúne varias
secciones; no corresponde a un único apartado del libro.

\textbf{Microcurrículo:} semana 8, sesión 2 (sesión global 16). Los
numerales del cronograma son identificadores curriculares; no son los
apartados del libro citados arriba.}
\end{frame}

\begin{frame}{Semana 9 · Welch y análisis tiempo-frecuencia}
\phantomsection\label{d2-099}
\note{\textbf{Libro guía:} John Semmlow, \emph{Circuits, Signals, and
Systems for Bioengineers}, 3.ª ed., 2018.

§3.4, \textbf{Time--Frequency Transformation In The Discrete Domain}
(página 142 del visor PDF); §4.2, \textbf{Data Acquisition And Storage}
(página 175 del visor PDF); §4.6, \textbf{Time--Frequency Analysis}
(página 205 del visor PDF); §5.4, \textbf{Using The Impulse Response To
Find A Systems Output To Any Input---Convolution} (página 219 del visor
PDF).

\textbf{Correspondencia:} Síntesis o encuadre que reúne varias
secciones; no corresponde a un único apartado del libro.

\textbf{Lectura:} use estas secciones como ruta de preparación y
consulta. La bibliografía aparece al final del bloque.}
\end{frame}

\begin{frame}{Sesión 17 · Welch reduce variabilidad mediante promediado}
\phantomsection\label{d2-100}
El método divide el registro en segmentos, aplica una ventana y promedia
periodogramas modificados.

\[\widehat S_{Welch}(f)=\frac{1}{K}\sum_{r=1}^{K}\widehat S_r(f).\]

\begin{block}{Concepto clave}

El promediado reduce variabilidad del estimador a cambio de resolución
frecuencial y dependencia entre segmentos solapados.

\end{block}

\note{\textbf{Libro guía:} John Semmlow, \emph{Circuits, Signals, and
Systems for Bioengineers}, 3.ª ed., 2018.

§4.4, \textbf{Spectral Averaging} (página 195 del visor PDF).

\textbf{Correspondencia:} El tema se desarrolla en las secciones
indicadas. La redacción es una síntesis docente.

\textbf{Microcurrículo:} semana 9, sesión 1 (sesión global 17). Los
numerales del cronograma son identificadores curriculares; no son los
apartados del libro citados arriba.}
\end{frame}

\begin{frame}{El procedimiento contiene cuatro decisiones}
\phantomsection\label{d2-101}
\begin{enumerate}
\tightlist
\item
  longitud de segmento \(L\);
\item
  solapamiento \(O\);
\item
  ventana \(w[n]\);
\item
  número de FFT \(N_{FFT}\).
\end{enumerate}

\begin{alertblock}{Atención}

Dos PSD calculadas con parámetros distintos pueden diferir aunque
provengan del mismo registro.

\end{alertblock}

\note{\textbf{Libro guía:} John Semmlow, \emph{Circuits, Signals, and
Systems for Bioengineers}, 3.ª ed., 2018.

§4.4, \textbf{Spectral Averaging} (página 195 del visor PDF).

\textbf{Correspondencia:} El tema se desarrolla en las secciones
indicadas. La redacción es una síntesis docente.

\textbf{Microcurrículo:} semana 9, sesión 1 (sesión global 17). Los
numerales del cronograma son identificadores curriculares; no son los
apartados del libro citados arriba.}
\end{frame}

\begin{frame}{Segmentos cortos producen más promedios}
\phantomsection\label{d2-102}
Al reducir \(L\):

\begin{itemize}
\tightlist
\item
  aumenta el número de segmentos;
\item
  suele disminuir la variabilidad;
\item
  empeora la resolución frecuencial;
\item
  mejora la aproximación de estacionariedad local.
\end{itemize}

\begin{block}{Concepto clave}

La longitud debe relacionarse con la separación frecuencial necesaria y
la escala temporal del fenómeno.

\end{block}

\note{\textbf{Libro guía:} John Semmlow, \emph{Circuits, Signals, and
Systems for Bioengineers}, 3.ª ed., 2018.

§4.4, \textbf{Spectral Averaging} (página 195 del visor PDF).

\textbf{Correspondencia:} El tema se desarrolla en las secciones
indicadas. La redacción es una síntesis docente.

\textbf{Microcurrículo:} semana 9, sesión 1 (sesión global 17). Los
numerales del cronograma son identificadores curriculares; no son los
apartados del libro citados arriba.}
\end{frame}

\begin{frame}{El solapamiento reutiliza muestras}
\phantomsection\label{d2-103}
Con avance \(H=L-O\):

\[\text{solapamiento relativo}=\frac{O}{L}=1-\frac{H}{L}.\]

\begin{alertblock}{Atención}

Más solapamiento no crea observaciones independientes. Su beneficio
depende de la ventana y del costo computacional.

\end{alertblock}

\note{\textbf{Libro guía:} John Semmlow, \emph{Circuits, Signals, and
Systems for Bioengineers}, 3.ª ed., 2018.

§4.4, \textbf{Spectral Averaging} (página 195 del visor PDF).

\textbf{Correspondencia:} El tema se desarrolla en las secciones
indicadas. La redacción es una síntesis docente.

\textbf{Microcurrículo:} semana 9, sesión 1 (sesión global 17). Los
numerales del cronograma son identificadores curriculares; no son los
apartados del libro citados arriba.}
\end{frame}

\begin{frame}{Hann y 50 \% forman una elección frecuente}
\phantomsection\label{d2-104}
La combinación suele ofrecer cobertura uniforme razonable y eficiencia
computacional.

\begin{block}{Nota}

``Frecuente'' no significa universal. Eventos transitorios, registros
breves o necesidades de amplitud pueden exigir otra selección.

\end{block}

\note{\textbf{Libro guía:} John Semmlow, \emph{Circuits, Signals, and
Systems for Bioengineers}, 3.ª ed., 2018.

§4.4, \textbf{Spectral Averaging} (página 195 del visor PDF).

\textbf{Correspondencia:} El tema se desarrolla en las secciones
indicadas. La redacción es una síntesis docente.

\textbf{Microcurrículo:} semana 9, sesión 1 (sesión global 17). Los
numerales del cronograma son identificadores curriculares; no son los
apartados del libro citados arriba.}
\end{frame}

\begin{frame}{El escalamiento incluye la energía de la ventana}
\phantomsection\label{d2-105}
El periodograma modificado compensa, según la convención, la suma
cuadrática de \(w[n]\):

\[U=\frac{1}{L}\sum_{n=0}^{L-1}w^2[n].\]

\begin{block}{Concepto clave}

Cada periodograma modificado integra \(\sum|w x_0|^2/\sum w^2\). Welch
promedia esas potencias ponderadas. La equivalencia con la potencia del
proceso requiere los supuestos estadísticos del estimador.

\end{block}

\note{\textbf{Libro guía:} John Semmlow, \emph{Circuits, Signals, and
Systems for Bioengineers}, 3.ª ed., 2018.

§4.4, \textbf{Spectral Averaging} (página 195 del visor PDF).

\textbf{Correspondencia:} El tema se desarrolla en las secciones
indicadas. La redacción es una síntesis docente.

\textbf{Microcurrículo:} semana 9, sesión 1 (sesión global 17). Los
numerales del cronograma son identificadores curriculares; no son los
apartados del libro citados arriba.}
\end{frame}

\begin{frame}{\(N_{FFT}\) cambia la malla, no la evidencia}
\phantomsection\label{d2-106}
Si \(N_{FFT}>L\), se aplica zero padding dentro de cada segmento.

\begin{alertblock}{Atención}

La densidad de puntos puede facilitar lectura o integración, pero la
resolución física sigue limitada por \(L/f_s\) y la ventana.

\end{alertblock}

\note{\textbf{Libro guía:} John Semmlow, \emph{Circuits, Signals, and
Systems for Bioengineers}, 3.ª ed., 2018.

§4.4, \textbf{Spectral Averaging} (página 195 del visor PDF).

\textbf{Correspondencia:} El tema se desarrolla en las secciones
indicadas. La redacción es una síntesis docente.

\textbf{Microcurrículo:} semana 9, sesión 1 (sesión global 17). Los
numerales del cronograma son identificadores curriculares; no son los
apartados del libro citados arriba.}
\end{frame}

\begin{frame}{Promediar requiere segmentos comparables}
\phantomsection\label{d2-107}
Welch supone que los segmentos describen propiedades espectrales
suficientemente semejantes.

Problemas:

\begin{itemize}
\tightlist
\item
  cambio de tarea;
\item
  transición fisiológica;
\item
  artefactos concentrados;
\item
  deriva prolongada.
\end{itemize}

\begin{alertblock}{Atención}

Promediar estados diferentes produce una PSD precisa para una mezcla que
puede carecer de significado fisiológico.

\end{alertblock}

\note{\textbf{Libro guía:} John Semmlow, \emph{Circuits, Signals, and
Systems for Bioengineers}, 3.ª ed., 2018.

§4.4, \textbf{Spectral Averaging} (página 195 del visor PDF).

\textbf{Correspondencia:} El tema se desarrolla en las secciones
indicadas. La redacción es una síntesis docente.

\textbf{Microcurrículo:} semana 9, sesión 1 (sesión global 17). Los
numerales del cronograma son identificadores curriculares; no son los
apartados del libro citados arriba.}
\end{frame}

\begin{frame}{Ancho de banda puede definirse de varias formas}
\phantomsection\label{d2-108}
El ancho de banda es una \textbf{extensión frecuencial}, expresada en
Hz. Debe definirse el criterio:

\begin{itemize}
\tightlist
\item
  diferencia entre bordes de paso especificados;
\item
  separación entre puntos de media potencia (\(-3\) dB), cuando aplica;
\item
  ancho del intervalo que contiene una fracción fijada de potencia;
\item
  ancho equivalente de ruido, según el problema.
\end{itemize}

\begin{block}{Concepto clave}

La frecuencia media y la mediana espectral describen localización del
espectro. No son medidas de su ancho. Dos PSD pueden tener igual centro
y distinta dispersión.

\end{block}

\note{\textbf{Libro guía:} John Semmlow, \emph{Circuits, Signals, and
Systems for Bioengineers}, 3.ª ed., 2018.

§4.5, \textbf{Signal Bandwidth} (página 200 del visor PDF).

\textbf{Correspondencia:} El tema se desarrolla en las secciones
indicadas. La redacción es una síntesis docente.

\textbf{Microcurrículo:} semana 9, sesión 1 (sesión global 17). Los
numerales del cronograma son identificadores curriculares; no son los
apartados del libro citados arriba.}
\end{frame}

\begin{frame}{Frecuencia mediana espectral}
\phantomsection\label{d2-109}
La frecuencia \(f_m\) satisface

\[\int_{0}^{f_m}S_{xx}(f)df
=\frac{1}{2}\int_{0}^{f_{max}}S_{xx}(f)df.\]

\begin{block}{Nota}

En EMG puede describir redistribución espectral, pero depende de fatiga,
fuerza, geometría del registro y preprocesamiento.

\end{block}

\note{\textbf{Libro guía:} John Semmlow, \emph{Circuits, Signals, and
Systems for Bioengineers}, 3.ª ed., 2018.

§4.3, \textbf{Power Spectrum} (página 191 del visor PDF); §4.5,
\textbf{Signal Bandwidth} (página 200 del visor PDF).

\textbf{Correspondencia:} Las secciones aportan el fundamento. La
diapositiva amplía la formulación o precisa condiciones que no se
presentan con idéntica notación en el libro. Frecuencia media/mediana
son descriptores de localización y no sinónimos de ancho de banda.

\textbf{Microcurrículo:} semana 9, sesión 1 (sesión global 17). Los
numerales del cronograma son identificadores curriculares; no son los
apartados del libro citados arriba.}
\end{frame}

\begin{frame}{Frecuencia media espectral}
\phantomsection\label{d2-110}
\[f_{mean}=\frac{\int fS_{xx}(f)df}{\int S_{xx}(f)df}.\]

\begin{alertblock}{Atención}

La frecuencia media es sensible a contaminación de alta frecuencia y a
los límites escogidos para la integración.

\end{alertblock}

\note{\textbf{Libro guía:} John Semmlow, \emph{Circuits, Signals, and
Systems for Bioengineers}, 3.ª ed., 2018.

§4.3, \textbf{Power Spectrum} (página 191 del visor PDF); §4.5,
\textbf{Signal Bandwidth} (página 200 del visor PDF).

\textbf{Correspondencia:} Las secciones aportan el fundamento. La
diapositiva amplía la formulación o precisa condiciones que no se
presentan con idéntica notación en el libro. Frecuencia media/mediana
son descriptores de localización y no sinónimos de ancho de banda.

\textbf{Microcurrículo:} semana 9, sesión 1 (sesión global 17). Los
numerales del cronograma son identificadores curriculares; no son los
apartados del libro citados arriba.}
\end{frame}

\begin{frame}[fragile]{Potencia de banda mediante integración numérica}
\phantomsection\label{d2-111}
\begin{Shaded}
\begin{Highlighting}[]
\ImportTok{from}\NormalTok{ scipy.signal }\ImportTok{import}\NormalTok{ welch}
\ImportTok{import}\NormalTok{ numpy }\ImportTok{as}\NormalTok{ np}

\NormalTok{f, Pxx }\OperatorTok{=}\NormalTok{ welch(x, fs}\OperatorTok{=}\NormalTok{fs, window}\OperatorTok{=}\StringTok{"hann"}\NormalTok{,}
\NormalTok{               nperseg}\OperatorTok{=}\NormalTok{L, noverlap}\OperatorTok{=}\NormalTok{L}\OperatorTok{//}\DecValTok{2}\NormalTok{,}
\NormalTok{               detrend}\OperatorTok{=}\StringTok{"constant"}\NormalTok{, scaling}\OperatorTok{=}\StringTok{"density"}\NormalTok{)}
\NormalTok{mask }\OperatorTok{=}\NormalTok{ (f }\OperatorTok{\textgreater{}=}\NormalTok{ f1) }\OperatorTok{\&}\NormalTok{ (f }\OperatorTok{\textless{}=}\NormalTok{ f2)}
\NormalTok{P\_band }\OperatorTok{=}\NormalTok{ np.trapezoid(Pxx[mask], f[mask])}
\end{Highlighting}
\end{Shaded}

\begin{block}{Nota}

La integración debe incluir suficientes puntos y límites coherentes con
la resolución disponible.

\end{block}

\note{\textbf{Libro guía:} John Semmlow, \emph{Circuits, Signals, and
Systems for Bioengineers}, 3.ª ed., 2018.

§4.3, \textbf{Power Spectrum} (página 191 del visor PDF); §4.4,
\textbf{Spectral Averaging} (página 195 del visor PDF).

\textbf{Correspondencia:} Ejemplo, figura o implementación de
elaboración docente a partir de estos fundamentos. No es una
transcripción del código MATLAB del libro.

\textbf{Microcurrículo:} semana 9, sesión 1 (sesión global 17). Los
numerales del cronograma son identificadores curriculares; no son los
apartados del libro citados arriba.

\textbf{Implementación:}
\href{https://docs.scipy.org/doc/scipy/reference/generated/scipy.signal.periodogram.html}{SciPy:
periodogram y escalamiento espectral}.}
\end{frame}

\begin{frame}{Selección razonada de parámetros}
\phantomsection\label{d2-112}
\begin{enumerate}
\tightlist
\item
  fijar la mínima separación frecuencial de interés;
\item
  elegir una duración compatible con esa separación;
\item
  verificar estacionariedad local;
\item
  seleccionar ventana según fuga y resolución;
\item
  usar solapamiento justificado;
\item
  validar potencia y estabilidad.
\end{enumerate}

\begin{block}{Comprobación}

Los parámetros forman parte del método y deben quedar registrados junto
con el resultado.

\end{block}

\note{\textbf{Libro guía:} John Semmlow, \emph{Circuits, Signals, and
Systems for Bioengineers}, 3.ª ed., 2018.

§4.4, \textbf{Spectral Averaging} (página 195 del visor PDF); §4.5,
\textbf{Signal Bandwidth} (página 200 del visor PDF).

\textbf{Correspondencia:} Actividad o interpretación biomédica de
elaboración docente apoyada en las secciones indicadas.

\textbf{Microcurrículo:} semana 9, sesión 1 (sesión global 17). Los
numerales del cronograma son identificadores curriculares; no son los
apartados del libro citados arriba.}
\end{frame}

\begin{frame}{Comparación entre condiciones}
\phantomsection\label{d2-113}
Para comparar reposo y tarea:

\begin{itemize}
\tightlist
\item
  igualar unidades y preprocesamiento;
\item
  usar parámetros espectrales iguales;
\item
  controlar duración y calidad;
\item
  comparar por participante cuando corresponda;
\item
  reportar variabilidad, no solo promedios.
\end{itemize}

\begin{alertblock}{Atención}

Normalizar cada condición por un denominador distinto puede cambiar la
interpretación de las diferencias.

\end{alertblock}

\note{\textbf{Libro guía:} John Semmlow, \emph{Circuits, Signals, and
Systems for Bioengineers}, 3.ª ed., 2018.

§4.4, \textbf{Spectral Averaging} (página 195 del visor PDF); §4.5,
\textbf{Signal Bandwidth} (página 200 del visor PDF).

\textbf{Correspondencia:} Actividad o interpretación biomédica de
elaboración docente apoyada en las secciones indicadas.

\textbf{Microcurrículo:} semana 9, sesión 1 (sesión global 17). Los
numerales del cronograma son identificadores curriculares; no son los
apartados del libro citados arriba.}
\end{frame}

\begin{frame}{Verificación de la sesión 17}
\phantomsection\label{d2-114}
Un registro de 60 s a 1000 Hz usa Welch con segmentos de 2 s y 50 \% de
solapamiento.

\begin{enumerate}
\tightlist
\item
  determine \(L\) y el avance \(H\);
\item
  calcule la separación nominal de bins;
\item
  explique qué cambia con segmentos de 4 s;
\item
  identifique el supuesto al promediar;
\item
  proponga una prueba de conservación de potencia.
\end{enumerate}

\begin{block}{Criterio de logro}

La respuesta debe discutir resolución, varianza, dependencia y
estacionariedad.

\end{block}

\note{\textbf{Libro guía:} John Semmlow, \emph{Circuits, Signals, and
Systems for Bioengineers}, 3.ª ed., 2018.

§4.4, \textbf{Spectral Averaging} (página 195 del visor PDF); §4.5,
\textbf{Signal Bandwidth} (página 200 del visor PDF).

\textbf{Correspondencia:} Actividad o interpretación biomédica de
elaboración docente apoyada en las secciones indicadas.

\textbf{Microcurrículo:} semana 9, sesión 1 (sesión global 17). Los
numerales del cronograma son identificadores curriculares; no son los
apartados del libro citados arriba.}
\end{frame}

\begin{frame}{Sesión 18 · Señales no estacionarias}
\phantomsection\label{d2-115}
Ejemplos:

\begin{itemize}
\tightlist
\item
  transición de reposo a ejercicio;
\item
  ráfagas EEG;
\item
  activación muscular durante marcha;
\item
  cambios de ritmo cardíaco.
\end{itemize}

\begin{block}{Concepto clave}

Un espectro global mezcla tiempos. El análisis tiempo-frecuencia busca
localizar cómo cambia el contenido espectral.

\end{block}

\note{\textbf{Libro guía:} John Semmlow, \emph{Circuits, Signals, and
Systems for Bioengineers}, 3.ª ed., 2018.

§4.6, \textbf{Time--Frequency Analysis} (página 205 del visor PDF).

\textbf{Correspondencia:} El tema se desarrolla en las secciones
indicadas. La redacción es una síntesis docente.

\textbf{Microcurrículo:} semana 9, sesión 2 (sesión global 18). Los
numerales del cronograma son identificadores curriculares; no son los
apartados del libro citados arriba.}
\end{frame}

\begin{frame}{La estacionariedad es un supuesto de modelamiento}
\phantomsection\label{d2-116}
Para un proceso aleatorio con segundo momento finito, estacionariedad en
sentido amplio exige:

\[\mathbb{E}\{x(t)\}=\mu,\]

\[R_x(t_1,t_2)=R_x(t_1-t_2).\]

\begin{alertblock}{Atención}

Las bioseñales rara vez cumplen estacionariedad global. Se trabaja con
segmentos donde las propiedades cambian poco para el objetivo.

\end{alertblock}

\note{\textbf{Libro guía:} John Semmlow, \emph{Circuits, Signals, and
Systems for Bioengineers}, 3.ª ed., 2018.

§10.2, \textbf{Stochastic Processes: Stationarity And Ergodicity}
(página 450 del visor PDF).

\textbf{Correspondencia:} El tema se desarrolla en las secciones
indicadas. La redacción es una síntesis docente.

\textbf{Microcurrículo:} semana 9, sesión 2 (sesión global 18). Los
numerales del cronograma son identificadores curriculares; no son los
apartados del libro citados arriba.}
\end{frame}

\begin{frame}{La STFT aplica Fourier localmente}
\phantomsection\label{d2-117}
\[X(\tau,f)=\int_{-\infty}^{\infty}x(t)w(t-\tau)e^{-j2\pi ft}dt.\]

\begin{itemize}
\tightlist
\item
  \(\tau\): posición temporal de la ventana.
\item
  \(f\): frecuencia analizada.
\item
  \(w\): ventana de análisis.
\end{itemize}

\begin{block}{Concepto clave}

La STFT produce una colección de espectros locales con una ventana de
forma y longitud fijas.

\end{block}

\note{\textbf{Libro guía:} John Semmlow, \emph{Circuits, Signals, and
Systems for Bioengineers}, 3.ª ed., 2018.

§4.6, \textbf{Time--Frequency Analysis} (página 205 del visor PDF).

\textbf{Correspondencia:} El tema se desarrolla en las secciones
indicadas. La redacción es una síntesis docente.

\textbf{Microcurrículo:} semana 9, sesión 2 (sesión global 18). Los
numerales del cronograma son identificadores curriculares; no son los
apartados del libro citados arriba.}
\end{frame}

\begin{frame}{El espectrograma muestra magnitud o potencia}
\phantomsection\label{d2-118}
Una definición habitual:

\[\mathcal{S}(\tau,f)=|X(\tau,f)|^2.\]

\begin{alertblock}{Atención}

Color, escala logarítmica y normalización cambian la apariencia. Una
imagen atractiva no garantiza unidades correctas.

\end{alertblock}

\note{\textbf{Libro guía:} John Semmlow, \emph{Circuits, Signals, and
Systems for Bioengineers}, 3.ª ed., 2018.

§4.6, \textbf{Time--Frequency Analysis} (página 205 del visor PDF).

\textbf{Correspondencia:} El tema se desarrolla en las secciones
indicadas. La redacción es una síntesis docente.

\textbf{Microcurrículo:} semana 9, sesión 2 (sesión global 18). Los
numerales del cronograma son identificadores curriculares; no son los
apartados del libro citados arriba.}
\end{frame}

\begin{frame}{La ventana fija impone un compromiso}
\phantomsection\label{d2-119}
Ventana corta:

\begin{itemize}
\tightlist
\item
  mejor localización temporal;
\item
  peor separación frecuencial.
\end{itemize}

Ventana larga:

\begin{itemize}
\tightlist
\item
  mejor separación frecuencial;
\item
  mayor mezcla temporal.
\end{itemize}

\begin{block}{Concepto clave}

La STFT no ofrece máxima resolución en ambos dominios. La pregunta
define el compromiso aceptable.

\end{block}

\note{\textbf{Libro guía:} John Semmlow, \emph{Circuits, Signals, and
Systems for Bioengineers}, 3.ª ed., 2018.

§4.6, \textbf{Time--Frequency Analysis} (página 205 del visor PDF).

\textbf{Correspondencia:} El tema se desarrolla en las secciones
indicadas. La redacción es una síntesis docente.

\textbf{Microcurrículo:} semana 9, sesión 2 (sesión global 18). Los
numerales del cronograma son identificadores curriculares; no son los
apartados del libro citados arriba.}
\end{frame}

\begin{frame}{La incertidumbre expresa ese límite}
\phantomsection\label{d2-120}
Para una ventana de energía finita, normalice \(|w(t)|^2\) y
\(|W(f)|^2\) como distribuciones de energía. Si sus desviaciones
estándar son \(\sigma_t\) y \(\sigma_f\):

\[\sigma_t\sigma_f\ge\frac{1}{4\pi}.\]

\begin{block}{Concepto clave}

Esta constante corresponde a Fourier con \(e^{-j2\pi ft}\) y a
dispersión RMS de energía. Con frecuencia angular se obtiene
\(\sigma_t\sigma_\omega\ge1/2\).

\end{block}

La gaussiana ideal alcanza la igualdad. Duración rectangular, ancho
entre ceros y separación de bins son medidas diferentes.

\note{\textbf{Libro guía:} John Semmlow, \emph{Circuits, Signals, and
Systems for Bioengineers}, 3.ª ed., 2018.

§4.6, \textbf{Time--Frequency Analysis} (página 205 del visor PDF).

\textbf{Correspondencia:} Las secciones aportan el fundamento. La
diapositiva amplía la formulación o precisa condiciones que no se
presentan con idéntica notación en el libro. Límite de dispersión
tiempo--frecuencia formalizado para la convención en Hz y desviaciones
de energía.

\textbf{Microcurrículo:} semana 9, sesión 2 (sesión global 18). Los
numerales del cronograma son identificadores curriculares; no son los
apartados del libro citados arriba.}
\end{frame}

\begin{frame}{El avance controla densidad temporal}
\phantomsection\label{d2-121}
Con ventana \(L\) y avance \(H\):

\[t_r=\frac{rH+\text{referencia de ventana}}{f_s}.\]

\begin{alertblock}{Atención}

Mayor solapamiento produce más columnas en el espectrograma, pero no
aumenta la resolución temporal física en la misma proporción.

\end{alertblock}

\note{\textbf{Libro guía:} John Semmlow, \emph{Circuits, Signals, and
Systems for Bioengineers}, 3.ª ed., 2018.

§4.6, \textbf{Time--Frequency Analysis} (página 205 del visor PDF).

\textbf{Correspondencia:} Actividad o interpretación biomédica de
elaboración docente apoyada en las secciones indicadas.

\textbf{Microcurrículo:} semana 9, sesión 2 (sesión global 18). Los
numerales del cronograma son identificadores curriculares; no son los
apartados del libro citados arriba.}
\end{frame}

\begin{frame}{Las condiciones de borde necesitan atención}
\phantomsection\label{d2-122}
Cerca del inicio y final, la ventana puede:

\begin{itemize}
\tightlist
\item
  contener menos datos reales;
\item
  usar relleno;
\item
  introducir transitorios de borde;
\item
  cambiar la interpretación temporal.
\end{itemize}

\begin{block}{Concepto clave}

Debe documentarse si se usa padding, centrado de ventanas y extensión
del registro.

\end{block}

\note{\textbf{Libro guía:} John Semmlow, \emph{Circuits, Signals, and
Systems for Bioengineers}, 3.ª ed., 2018.

§4.6, \textbf{Time--Frequency Analysis} (página 205 del visor PDF).

\textbf{Correspondencia:} Actividad o interpretación biomédica de
elaboración docente apoyada en las secciones indicadas.

\textbf{Microcurrículo:} semana 9, sesión 2 (sesión global 18). Los
numerales del cronograma son identificadores curriculares; no son los
apartados del libro citados arriba.}
\end{frame}

\begin{frame}{La escala logarítmica revela componentes débiles}
\phantomsection\label{d2-123}
Una representación frecuente es

\[P_{dB}(\tau,f)=10\log_{10}\left(\frac{P(\tau,f)}{P_{ref}}\right).\]

\begin{alertblock}{Atención}

Sin \(P_{ref}\) y unidades, ``dB'' queda incompleto. Ceros o valores muy
pequeños requieren un piso numérico declarado.

\end{alertblock}

\note{\textbf{Libro guía:} John Semmlow, \emph{Circuits, Signals, and
Systems for Bioengineers}, 3.ª ed., 2018.

§4.6, \textbf{Time--Frequency Analysis} (página 205 del visor PDF).

\textbf{Correspondencia:} Actividad o interpretación biomédica de
elaboración docente apoyada en las secciones indicadas.

\textbf{Microcurrículo:} semana 9, sesión 2 (sesión global 18). Los
numerales del cronograma son identificadores curriculares; no son los
apartados del libro citados arriba.}
\end{frame}

\begin{frame}{STFT de una activación EMG}
\phantomsection\label{d2-124}
Una ventana corta puede localizar inicio y final de activación. Una
ventana más larga estabiliza la distribución espectral dentro de la
contracción.

\begin{block}{Nota}

La envolvente responde a variación de amplitud lenta. El espectrograma
describe además cómo se distribuye energía dentro de la banda.

\end{block}

\note{\textbf{Libro guía:} John Semmlow, \emph{Circuits, Signals, and
Systems for Bioengineers}, 3.ª ed., 2018.

§4.6, \textbf{Time--Frequency Analysis} (página 205 del visor PDF).

\textbf{Correspondencia:} Actividad o interpretación biomédica de
elaboración docente apoyada en las secciones indicadas.

\textbf{Microcurrículo:} semana 9, sesión 2 (sesión global 18). Los
numerales del cronograma son identificadores curriculares; no son los
apartados del libro citados arriba.}
\end{frame}

\begin{frame}{STFT de una ráfaga EEG}
\phantomsection\label{d2-125}
Preguntas relevantes:

\begin{itemize}
\tightlist
\item
  ¿cuándo comienza la ráfaga?
\item
  ¿qué banda domina?
\item
  ¿cambia su frecuencia central?
\item
  ¿aparece simultáneamente en varios canales?
\end{itemize}

\begin{alertblock}{Atención}

Una banda visible no identifica por sí sola una fuente neural.
Artefactos o referencia común pueden generar patrones similares.

\end{alertblock}

\note{\textbf{Libro guía:} John Semmlow, \emph{Circuits, Signals, and
Systems for Bioengineers}, 3.ª ed., 2018.

§4.6, \textbf{Time--Frequency Analysis} (página 205 del visor PDF).

\textbf{Correspondencia:} Actividad o interpretación biomédica de
elaboración docente apoyada en las secciones indicadas.

\textbf{Microcurrículo:} semana 9, sesión 2 (sesión global 18). Los
numerales del cronograma son identificadores curriculares; no son los
apartados del libro citados arriba.}
\end{frame}

\begin{frame}[fragile]{Código: STFT con densidad espectral}
\phantomsection\label{d2-126}
\begin{Shaded}
\begin{Highlighting}[]
\ImportTok{import}\NormalTok{ numpy }\ImportTok{as}\NormalTok{ np}
\ImportTok{from}\NormalTok{ scipy.signal }\ImportTok{import}\NormalTok{ stft}
\NormalTok{f, t, Z }\OperatorTok{=}\NormalTok{ stft(}
\NormalTok{    x, fs}\OperatorTok{=}\NormalTok{fs, window}\OperatorTok{=}\StringTok{"hann"}\NormalTok{,}
\NormalTok{    nperseg}\OperatorTok{=}\NormalTok{L, noverlap}\OperatorTok{=}\NormalTok{L}\OperatorTok{{-}}\NormalTok{H,}
\NormalTok{    nfft}\OperatorTok{=}\NormalTok{Nfft, detrend}\OperatorTok{=}\StringTok{"constant"}\NormalTok{,}
\NormalTok{    boundary}\OperatorTok{=}\VariableTok{None}\NormalTok{, padded}\OperatorTok{=}\VariableTok{False}\NormalTok{, scaling}\OperatorTok{=}\StringTok{"psd"}\NormalTok{,}
\NormalTok{)}
\NormalTok{P }\OperatorTok{=}\NormalTok{ np.}\BuiltInTok{abs}\NormalTok{(Z)}\OperatorTok{**}\DecValTok{2}
\ControlFlowTok{if}\NormalTok{ Nfft }\OperatorTok{\%} \DecValTok{2} \OperatorTok{==} \DecValTok{0}\NormalTok{:}
\NormalTok{    P[}\DecValTok{1}\NormalTok{:}\OperatorTok{{-}}\DecValTok{1}\NormalTok{, :] }\OperatorTok{*=} \DecValTok{2}
\ControlFlowTok{else}\NormalTok{:}
\NormalTok{    P[}\DecValTok{1}\NormalTok{:, :] }\OperatorTok{*=} \DecValTok{2}
\end{Highlighting}
\end{Shaded}

\begin{block}{Nota}

Para entrada real: densidad unilateral en \([x]^2\)/Hz.

\end{block}

\note{\textbf{Libro guía:} John Semmlow, \emph{Circuits, Signals, and
Systems for Bioengineers}, 3.ª ed., 2018.

§4.6, \textbf{Time--Frequency Analysis} (página 205 del visor PDF).

\textbf{Correspondencia:} Ejemplo, figura o implementación de
elaboración docente a partir de estos fundamentos. No es una
transcripción del código MATLAB del libro. Con scaling=``psd'', stft
devuelve aquí el semiespectro complejo. Su módulo cuadrado requiere
duplicar los bins interiores para combinar la potencia de las
frecuencias positivas y negativas.

\textbf{Microcurrículo:} semana 9, sesión 2 (sesión global 18). Los
numerales del cronograma son identificadores curriculares; no son los
apartados del libro citados arriba.

\textbf{Implementación:}
\href{https://docs.scipy.org/doc/scipy/reference/generated/scipy.signal.stft.html}{SciPy:
stft}. Función mantenida como legacy; se conserva para continuidad del
curso. Se explicita scaling y conversión de potencia unilateral.}
\end{frame}

\begin{frame}{Validar un espectrograma}
\phantomsection\label{d2-127}
\begin{enumerate}
\tightlist
\item
  contrastar eventos con el dominio temporal;
\item
  verificar frecuencias con señales sintéticas;
\item
  revisar efectos de ventana y borde;
\item
  repetir con otra longitud razonable;
\item
  comparar potencia integrada con estimaciones temporales;
\item
  inspeccionar canales de artefacto.
\end{enumerate}

\begin{block}{Comprobación}

Una interpretación robusta debe persistir bajo cambios moderados de
parámetros compatibles con la pregunta.

\end{block}

\note{\textbf{Libro guía:} John Semmlow, \emph{Circuits, Signals, and
Systems for Bioengineers}, 3.ª ed., 2018.

§4.6, \textbf{Time--Frequency Analysis} (página 205 del visor PDF).

\textbf{Correspondencia:} Actividad o interpretación biomédica de
elaboración docente apoyada en las secciones indicadas.

\textbf{Microcurrículo:} semana 9, sesión 2 (sesión global 18). Los
numerales del cronograma son identificadores curriculares; no son los
apartados del libro citados arriba.}
\end{frame}

\begin{frame}{Verificación de la sesión 18}
\phantomsection\label{d2-128}
Una ráfaga de 20 Hz dura 300 ms en un EEG a 250 Hz.

\begin{enumerate}
\tightlist
\item
  compare ventanas de 128 ms y 1 s;
\item
  indique cuál localiza mejor el inicio;
\item
  indique cuál separa mejor 20 Hz de 22 Hz;
\item
  discuta el efecto del solapamiento;
\item
  proponga canales o marcas para validar origen.
\end{enumerate}

\begin{block}{Criterio de logro}

La respuesta debe conectar duración del evento, lóbulo temporal,
resolución y evidencia fisiológica.

\end{block}

\note{\textbf{Libro guía:} John Semmlow, \emph{Circuits, Signals, and
Systems for Bioengineers}, 3.ª ed., 2018.

§4.6, \textbf{Time--Frequency Analysis} (página 205 del visor PDF).

\textbf{Correspondencia:} Actividad o interpretación biomédica de
elaboración docente apoyada en las secciones indicadas.

\textbf{Microcurrículo:} semana 9, sesión 2 (sesión global 18). Los
numerales del cronograma son identificadores curriculares; no son los
apartados del libro citados arriba.}
\end{frame}

\begin{frame}{Cierre de la semana 9}
\phantomsection\label{d2-129}
\begin{itemize}
\tightlist
\item
  Welch promedia periodogramas modificados para reducir variabilidad.
\item
  La longitud de segmento controla resolución y aproximación local.
\item
  Ancho de banda necesita una definición reproducible.
\item
  La STFT localiza cambios espectrales con resolución limitada.
\item
  Ventana, solapamiento, padding y escala forman parte del resultado.
\end{itemize}

\begin{block}{Pregunta de salida}

¿Por qué un espectrograma con más columnas puede no contener mejor
resolución temporal?

\end{block}

\note{\textbf{Libro guía:} John Semmlow, \emph{Circuits, Signals, and
Systems for Bioengineers}, 3.ª ed., 2018.

§4.4, \textbf{Spectral Averaging} (página 195 del visor PDF); §4.5,
\textbf{Signal Bandwidth} (página 200 del visor PDF); §4.6,
\textbf{Time--Frequency Analysis} (página 205 del visor PDF).

\textbf{Correspondencia:} Síntesis o encuadre que reúne varias
secciones; no corresponde a un único apartado del libro.

\textbf{Microcurrículo:} semana 9, sesión 2 (sesión global 18). Los
numerales del cronograma son identificadores curriculares; no son los
apartados del libro citados arriba.}
\end{frame}

\begin{frame}{Semana 10 · Ventanas tiempo-frecuencia y sistemas LTI}
\phantomsection\label{d2-130}
\note{\textbf{Libro guía:} John Semmlow, \emph{Circuits, Signals, and
Systems for Bioengineers}, 3.ª ed., 2018.

§3.4, \textbf{Time--Frequency Transformation In The Discrete Domain}
(página 142 del visor PDF); §4.2, \textbf{Data Acquisition And Storage}
(página 175 del visor PDF); §4.6, \textbf{Time--Frequency Analysis}
(página 205 del visor PDF); §5.4, \textbf{Using The Impulse Response To
Find A Systems Output To Any Input---Convolution} (página 219 del visor
PDF).

\textbf{Correspondencia:} Síntesis o encuadre que reúne varias
secciones; no corresponde a un único apartado del libro.

\textbf{Lectura:} use estas secciones como ruta de preparación y
consulta. La bibliografía aparece al final del bloque.}
\end{frame}

\begin{frame}{Sesión 19 · La ventana debe responder a la pregunta}
\phantomsection\label{d2-131}
Tres preguntas dominan la selección:

\begin{enumerate}
\tightlist
\item
  ¿qué duración tiene el evento?
\item
  ¿qué separación frecuencial se necesita?
\item
  ¿qué diferencia de amplitud existe entre componentes?
\end{enumerate}

\begin{block}{Concepto clave}

La forma y longitud de ventana se eligen conjuntamente. Ninguna ventana
domina en todos los criterios.

\end{block}

\note{\textbf{Libro guía:} John Semmlow, \emph{Circuits, Signals, and
Systems for Bioengineers}, 3.ª ed., 2018.

§4.2.4, \textbf{Data Truncation---Window Functions} (página 188 del
visor PDF); §4.6, \textbf{Time--Frequency Analysis} (página 205 del
visor PDF).

\textbf{Correspondencia:} El tema se desarrolla en las secciones
indicadas. La redacción es una síntesis docente.

\textbf{Microcurrículo:} semana 10, sesión 1 (sesión global 19). Los
numerales del cronograma son identificadores curriculares; no son los
apartados del libro citados arriba.}
\end{frame}

\begin{frame}{La longitud fija una escala de observación}
\phantomsection\label{d2-132}
Para \(L\) muestras:

\[T_w=\frac{L}{f_s},\qquad \Delta f_{nom}=\frac{f_s}{L}.\]

\begin{alertblock}{Atención}

La resolución efectiva depende además de la forma de la ventana y del
criterio usado para distinguir componentes.

\end{alertblock}

\note{\textbf{Libro guía:} John Semmlow, \emph{Circuits, Signals, and
Systems for Bioengineers}, 3.ª ed., 2018.

§4.2.4, \textbf{Data Truncation---Window Functions} (página 188 del
visor PDF); §4.6, \textbf{Time--Frequency Analysis} (página 205 del
visor PDF).

\textbf{Correspondencia:} El tema se desarrolla en las secciones
indicadas. La redacción es una síntesis docente.

\textbf{Microcurrículo:} semana 10, sesión 1 (sesión global 19). Los
numerales del cronograma son identificadores curriculares; no son los
apartados del libro citados arriba.}
\end{frame}

\begin{frame}{Rectangular prioriza un lóbulo principal estrecho}
\phantomsection\label{d2-133}
Ventajas:

\begin{itemize}
\tightlist
\item
  máxima conservación de un bloque alineado;
\item
  buena separación nominal de tonos comparables.
\end{itemize}

Limitaciones:

\begin{itemize}
\tightlist
\item
  lóbulos laterales altos;
\item
  sensibilidad a discontinuidades de borde.
\end{itemize}

\begin{block}{Nota}

Resulta útil como referencia para comprender por qué el recorte directo
produce fuga.

\end{block}

\note{\textbf{Libro guía:} John Semmlow, \emph{Circuits, Signals, and
Systems for Bioengineers}, 3.ª ed., 2018.

§4.2.4, \textbf{Data Truncation---Window Functions} (página 188 del
visor PDF); §4.6, \textbf{Time--Frequency Analysis} (página 205 del
visor PDF).

\textbf{Correspondencia:} El tema se desarrolla en las secciones
indicadas. La redacción es una síntesis docente.

\textbf{Microcurrículo:} semana 10, sesión 1 (sesión global 19). Los
numerales del cronograma son identificadores curriculares; no son los
apartados del libro citados arriba.}
\end{frame}

\begin{frame}{Hann ofrece un compromiso general}
\phantomsection\label{d2-134}
Ventajas:

\begin{itemize}
\tightlist
\item
  extremos nulos;
\item
  menor fuga lateral que rectangular;
\item
  uso frecuente en Welch y STFT.
\end{itemize}

Limitación principal: lóbulo central más ancho.

\begin{block}{Concepto clave}

Hann suele ser un punto de partida defendible, no una selección
automática.

\end{block}

\note{\textbf{Libro guía:} John Semmlow, \emph{Circuits, Signals, and
Systems for Bioengineers}, 3.ª ed., 2018.

§4.2.4, \textbf{Data Truncation---Window Functions} (página 188 del
visor PDF); §4.6, \textbf{Time--Frequency Analysis} (página 205 del
visor PDF).

\textbf{Correspondencia:} El tema se desarrolla en las secciones
indicadas. La redacción es una síntesis docente.

\textbf{Microcurrículo:} semana 10, sesión 1 (sesión global 19). Los
numerales del cronograma son identificadores curriculares; no son los
apartados del libro citados arriba.}
\end{frame}

\begin{frame}{Hamming mantiene extremos pequeños, no nulos}
\phantomsection\label{d2-135}
Para la forma simétrica:

\[w[n]=0.54-0.46\cos\left(\frac{2\pi n}{N-1}\right),\qquad
w[0]=w[N-1]=0.08.\]

Hamming atenúa especialmente el primer lóbulo lateral. Sus extremos no
nulos pueden dejar un salto en la señal ventaneada si los extremos del
registro difieren.

\begin{block}{Nota}

Extremos no nulos de la ventana no implican por sí solos discontinuidad
en su extensión periódica: ambos pueden ser iguales. Distinga la ventana
de la señal multiplicada por ella.

\end{block}

\note{\textbf{Libro guía:} John Semmlow, \emph{Circuits, Signals, and
Systems for Bioengineers}, 3.ª ed., 2018.

§4.2.4, \textbf{Data Truncation---Window Functions} (página 188 del
visor PDF); §4.6, \textbf{Time--Frequency Analysis} (página 205 del
visor PDF).

\textbf{Correspondencia:} El tema se desarrolla en las secciones
indicadas. La redacción es una síntesis docente.

\textbf{Microcurrículo:} semana 10, sesión 1 (sesión global 19). Los
numerales del cronograma son identificadores curriculares; no son los
apartados del libro citados arriba.}
\end{frame}

\begin{frame}{Blackman prioriza rango dinámico}
\phantomsection\label{d2-136}
Sus lóbulos laterales bajos ayudan a observar componentes débiles cerca
de componentes intensas.

Costo: lóbulo principal más ancho y menor separación de frecuencias
próximas.

\begin{alertblock}{Atención}

Reducir fuga lateral no mejora simultáneamente todas las formas de
resolución.

\end{alertblock}

\note{\textbf{Libro guía:} John Semmlow, \emph{Circuits, Signals, and
Systems for Bioengineers}, 3.ª ed., 2018.

§4.2.4, \textbf{Data Truncation---Window Functions} (página 188 del
visor PDF); §4.6, \textbf{Time--Frequency Analysis} (página 205 del
visor PDF).

\textbf{Correspondencia:} El tema se desarrolla en las secciones
indicadas. La redacción es una síntesis docente.

\textbf{Microcurrículo:} semana 10, sesión 1 (sesión global 19). Los
numerales del cronograma son identificadores curriculares; no son los
apartados del libro citados arriba.}
\end{frame}

\begin{frame}{Una ventana gaussiana localiza de forma equilibrada}
\phantomsection\label{d2-137}
La familia gaussiana permite ajustar el compromiso mediante su
dispersión y posee buenas propiedades de concentración
tiempo-frecuencia.

\begin{block}{Nota}

La elección del parámetro debe quedar documentada. ``Gaussiana'' no
define por sí sola una ventana única.

\end{block}

\note{\textbf{Libro guía:} John Semmlow, \emph{Circuits, Signals, and
Systems for Bioengineers}, 3.ª ed., 2018.

§4.2.4, \textbf{Data Truncation---Window Functions} (página 188 del
visor PDF); §4.6, \textbf{Time--Frequency Analysis} (página 205 del
visor PDF).

\textbf{Correspondencia:} El tema se desarrolla en las secciones
indicadas. La redacción es una síntesis docente.

\textbf{Microcurrículo:} semana 10, sesión 1 (sesión global 19). Los
numerales del cronograma son identificadores curriculares; no son los
apartados del libro citados arriba.}
\end{frame}

\begin{frame}{Ventana y evento deben tener escalas compatibles}
\phantomsection\label{d2-138}
Ejemplos:

\begin{itemize}
\tightlist
\item
  QRS: decenas a pocos cientos de milisegundos;
\item
  activación EMG: depende de la tarea y músculo;
\item
  oscilaciones respiratorias: varios segundos;
\item
  ritmos EEG sostenidos: duración definida por hipótesis.
\end{itemize}

\begin{block}{Concepto clave}

Una ventana mucho mayor que el evento diluye su localización. Una
ventana demasiado corta pierde detalle frecuencial.

\end{block}

\note{\textbf{Libro guía:} John Semmlow, \emph{Circuits, Signals, and
Systems for Bioengineers}, 3.ª ed., 2018.

§4.6, \textbf{Time--Frequency Analysis} (página 205 del visor PDF).

\textbf{Correspondencia:} Actividad o interpretación biomédica de
elaboración docente apoyada en las secciones indicadas.

\textbf{Microcurrículo:} semana 10, sesión 1 (sesión global 19). Los
numerales del cronograma son identificadores curriculares; no son los
apartados del libro citados arriba.}
\end{frame}

\begin{frame}{La frecuencia de interés determina ciclos observados}
\phantomsection\label{d2-139}
Número aproximado de ciclos dentro de la ventana:

\[N_c=f_0T_w.\]

\begin{alertblock}{Atención}

Estimar una frecuencia a partir de menos de un ciclo completo es
inestable y depende fuertemente del modelo.

\end{alertblock}

\note{\textbf{Libro guía:} John Semmlow, \emph{Circuits, Signals, and
Systems for Bioengineers}, 3.ª ed., 2018.

§4.6, \textbf{Time--Frequency Analysis} (página 205 del visor PDF).

\textbf{Correspondencia:} Actividad o interpretación biomédica de
elaboración docente apoyada en las secciones indicadas.

\textbf{Microcurrículo:} semana 10, sesión 1 (sesión global 19). Los
numerales del cronograma son identificadores curriculares; no son los
apartados del libro citados arriba.}
\end{frame}

\begin{frame}{Solapamiento mejora seguimiento visual}
\phantomsection\label{d2-140}
Reducir el avance \(H\) produce estimaciones más densas en el tiempo.

\begin{alertblock}{Atención}

Las columnas vecinas comparten datos y son dependientes. No deben
contarse como observaciones independientes en análisis estadístico.

\end{alertblock}

\note{\textbf{Libro guía:} John Semmlow, \emph{Circuits, Signals, and
Systems for Bioengineers}, 3.ª ed., 2018.

§4.6, \textbf{Time--Frequency Analysis} (página 205 del visor PDF).

\textbf{Correspondencia:} Actividad o interpretación biomédica de
elaboración docente apoyada en las secciones indicadas.

\textbf{Microcurrículo:} semana 10, sesión 1 (sesión global 19). Los
numerales del cronograma son identificadores curriculares; no son los
apartados del libro citados arriba.}
\end{frame}

\begin{frame}{Zero padding facilita localizar máximos}
\phantomsection\label{d2-141}
Una malla frecuencial densa puede ayudar a interpolar el máximo de un
lóbulo.

\begin{block}{Concepto clave}

La precisión de la estimación sigue limitada por SNR, duración, ventana
y estabilidad de la frecuencia.

\end{block}

\note{\textbf{Libro guía:} John Semmlow, \emph{Circuits, Signals, and
Systems for Bioengineers}, 3.ª ed., 2018.

§4.6, \textbf{Time--Frequency Analysis} (página 205 del visor PDF).

\textbf{Correspondencia:} Actividad o interpretación biomédica de
elaboración docente apoyada en las secciones indicadas.

\textbf{Microcurrículo:} semana 10, sesión 1 (sesión global 19). Los
numerales del cronograma son identificadores curriculares; no son los
apartados del libro citados arriba.}
\end{frame}

\begin{frame}{Normalizar entre ventanas requiere coherencia}
\phantomsection\label{d2-142}
Opciones distintas responden preguntas distintas:

\begin{itemize}
\tightlist
\item
  potencia absoluta;
\item
  potencia relativa por ventana;
\item
  cambio respecto de una línea base;
\item
  puntuación estandarizada por frecuencia.
\end{itemize}

\begin{alertblock}{Atención}

Normalizar cada columna por su propia potencia elimina cambios globales
de amplitud y puede ocultar activaciones reales.

\end{alertblock}

\note{\textbf{Libro guía:} John Semmlow, \emph{Circuits, Signals, and
Systems for Bioengineers}, 3.ª ed., 2018.

§4.6, \textbf{Time--Frequency Analysis} (página 205 del visor PDF).

\textbf{Correspondencia:} Actividad o interpretación biomédica de
elaboración docente apoyada en las secciones indicadas.

\textbf{Microcurrículo:} semana 10, sesión 1 (sesión global 19). Los
numerales del cronograma son identificadores curriculares; no son los
apartados del libro citados arriba.}
\end{frame}

\begin{frame}{Aplicación: activación muscular durante marcha}
\phantomsection\label{d2-143}
Una STFT puede describir cuándo aumenta energía EMG y si su distribución
cambia dentro del ciclo.

Diseño:

\begin{itemize}
\tightlist
\item
  alinear ciclos con eventos de marcha;
\item
  conservar escala o usar referencia fisiológica;
\item
  separar análisis temporal y espectral;
\item
  validar artefactos de movimiento.
\end{itemize}

\begin{block}{Concepto clave}

El espectrograma debe interpretarse junto con el ciclo de marcha y no
como una imagen aislada.

\end{block}

\note{\textbf{Libro guía:} John Semmlow, \emph{Circuits, Signals, and
Systems for Bioengineers}, 3.ª ed., 2018.

§4.6, \textbf{Time--Frequency Analysis} (página 205 del visor PDF).

\textbf{Correspondencia:} Actividad o interpretación biomédica de
elaboración docente apoyada en las secciones indicadas.

\textbf{Microcurrículo:} semana 10, sesión 1 (sesión global 19). Los
numerales del cronograma son identificadores curriculares; no son los
apartados del libro citados arriba.}
\end{frame}

\begin{frame}{Aplicación: cambio de ritmo respiratorio}
\phantomsection\label{d2-144}
Ventanas largas separan frecuencias respiratorias próximas. Ventanas
cortas localizan transiciones de ritmo.

\begin{block}{Nota}

Una estrategia útil compara varias longitudes plausibles y reporta qué
conclusión permanece estable.

\end{block}

\note{\textbf{Libro guía:} John Semmlow, \emph{Circuits, Signals, and
Systems for Bioengineers}, 3.ª ed., 2018.

§4.6, \textbf{Time--Frequency Analysis} (página 205 del visor PDF).

\textbf{Correspondencia:} Actividad o interpretación biomédica de
elaboración docente apoyada en las secciones indicadas.

\textbf{Microcurrículo:} semana 10, sesión 1 (sesión global 19). Los
numerales del cronograma son identificadores curriculares; no son los
apartados del libro citados arriba.}
\end{frame}

\begin{frame}{Flujo de selección documentada}
\phantomsection\label{d2-145}
\begin{enumerate}
\tightlist
\item
  formular evento o banda de interés;
\item
  estimar su escala temporal;
\item
  fijar separación frecuencial necesaria;
\item
  elegir longitud y forma;
\item
  definir avance y padding;
\item
  validar con señales conocidas;
\item
  realizar análisis de sensibilidad.
\end{enumerate}

\begin{block}{Comprobación}

La selección final debe poder reproducirse a partir de parámetros
guardados, no de ajustes visuales manuales.

\end{block}

\note{\textbf{Libro guía:} John Semmlow, \emph{Circuits, Signals, and
Systems for Bioengineers}, 3.ª ed., 2018.

§4.6, \textbf{Time--Frequency Analysis} (página 205 del visor PDF).

\textbf{Correspondencia:} Actividad o interpretación biomédica de
elaboración docente apoyada en las secciones indicadas.

\textbf{Microcurrículo:} semana 10, sesión 1 (sesión global 19). Los
numerales del cronograma son identificadores curriculares; no son los
apartados del libro citados arriba.}
\end{frame}

\begin{frame}{Verificación de la sesión 19}
\phantomsection\label{d2-146}
Se desea localizar una activación de 250 ms y distinguir componentes
separadas 2 Hz.

\begin{enumerate}
\tightlist
\item
  analice la tensión entre ambos objetivos;
\item
  calcule la duración mínima nominal para 2 Hz;
\item
  proponga una longitud inicial;
\item
  seleccione una ventana y justifique;
\item
  describa un análisis de sensibilidad.
\end{enumerate}

\begin{block}{Criterio de logro}

La solución debe reconocer que una única STFT puede no satisfacer
simultáneamente ambos objetivos.

\end{block}

\note{\textbf{Libro guía:} John Semmlow, \emph{Circuits, Signals, and
Systems for Bioengineers}, 3.ª ed., 2018.

§4.2.4, \textbf{Data Truncation---Window Functions} (página 188 del
visor PDF); §4.6, \textbf{Time--Frequency Analysis} (página 205 del
visor PDF).

\textbf{Correspondencia:} El tema se desarrolla en las secciones
indicadas. La redacción es una síntesis docente.

\textbf{Microcurrículo:} semana 10, sesión 1 (sesión global 19). Los
numerales del cronograma son identificadores curriculares; no son los
apartados del libro citados arriba.}
\end{frame}

\begin{frame}{Sesión 20 · Un sistema transforma una señal}
\phantomsection\label{d2-147}
\[y=\mathcal{T}\{x\}.\]

El análisis pregunta qué propiedades de \(\mathcal{T}\) permiten
predecir la salida para entradas diferentes.

\begin{block}{Concepto clave}

Las propiedades pertenecen al modelo del sistema dentro de condiciones
definidas. Un sistema fisiológico puede ser aproximadamente lineal solo
en un rango.

\end{block}

\note{\textbf{Libro guía:} John Semmlow, \emph{Circuits, Signals, and
Systems for Bioengineers}, 3.ª ed., 2018.

§5.2, \textbf{Linear Systems Analysis---An Overview} (página 212 del
visor PDF).

\textbf{Correspondencia:} El tema se desarrolla en las secciones
indicadas. La redacción es una síntesis docente.

\textbf{Microcurrículo:} semana 10, sesión 2 (sesión global 20). Los
numerales del cronograma son identificadores curriculares; no son los
apartados del libro citados arriba.}
\end{frame}

\begin{frame}{Linealidad combina aditividad y homogeneidad}
\phantomsection\label{d2-148}
\[\mathcal{T}\{a x_1+b x_2\}
=a\mathcal{T}\{x_1\}+b\mathcal{T}\{x_2\}.\]

\begin{alertblock}{Atención}

Verificar una sola entrada no demuestra linealidad. La propiedad debe
sostenerse para las combinaciones consideradas.

\end{alertblock}

\note{\textbf{Libro guía:} John Semmlow, \emph{Circuits, Signals, and
Systems for Bioengineers}, 3.ª ed., 2018.

§5.2.1, \textbf{Superposition And Linearity} (página 213 del visor PDF).

\textbf{Correspondencia:} El tema se desarrolla en las secciones
indicadas. La redacción es una síntesis docente.

\textbf{Microcurrículo:} semana 10, sesión 2 (sesión global 20). Los
numerales del cronograma son identificadores curriculares; no son los
apartados del libro citados arriba.}
\end{frame}

\begin{frame}{Saturación viola linealidad}
\phantomsection\label{d2-149}
Un amplificador ideal puede cumplir \(y=Gx\) dentro de su rango. Al
alcanzar los rieles, la salida se recorta.

\begin{block}{Concepto clave}

La linealidad depende del rango operativo. Un modelo local puede ser
útil aunque el dispositivo completo sea no lineal.

\end{block}

\note{\textbf{Libro guía:} John Semmlow, \emph{Circuits, Signals, and
Systems for Bioengineers}, 3.ª ed., 2018.

§1.4.4, \textbf{Linear Versus Nonlinear Signals And Systems} (página 44
del visor PDF); §5.2.1, \textbf{Superposition And Linearity} (página 213
del visor PDF).

\textbf{Correspondencia:} El tema se desarrolla en las secciones
indicadas. La redacción es una síntesis docente.

\textbf{Microcurrículo:} semana 10, sesión 2 (sesión global 20). Los
numerales del cronograma son identificadores curriculares; no son los
apartados del libro citados arriba.}
\end{frame}

\begin{frame}{Invariancia temporal conserva la regla}
\phantomsection\label{d2-150}
Si

\[y(t)=\mathcal{T}\{x(t)\},\]

un sistema invariante cumple

\[\mathcal{T}\{x(t-t_0)\}=y(t-t_0).\]

\begin{block}{Nota}

El desplazamiento de la entrada debe producir el mismo desplazamiento de
la salida, sin cambiar forma ni ganancia.

\end{block}

\note{\textbf{Libro guía:} John Semmlow, \emph{Circuits, Signals, and
Systems for Bioengineers}, 3.ª ed., 2018.

§1.4.4, \textbf{Linear Versus Nonlinear Signals And Systems} (página 44
del visor PDF); §5.2, \textbf{Linear Systems Analysis---An Overview}
(página 212 del visor PDF).

\textbf{Correspondencia:} Las secciones aportan el fundamento. La
diapositiva amplía la formulación o precisa condiciones que no se
presentan con idéntica notación en el libro. La invariancia temporal se
comprueba separadamente de la linealidad.

\textbf{Microcurrículo:} semana 10, sesión 2 (sesión global 20). Los
numerales del cronograma son identificadores curriculares; no son los
apartados del libro citados arriba.}
\end{frame}

\begin{frame}{Adaptación fisiológica rompe invariancia}
\phantomsection\label{d2-151}
La respuesta a un mismo estímulo puede cambiar por fatiga, aprendizaje,
medicación o estado autonómico.

\begin{alertblock}{Atención}

Aplicar herramientas LTI a un sistema biológico exige justificar
intervalo, amplitud y estado donde la aproximación resulta válida.

\end{alertblock}

\note{\textbf{Libro guía:} John Semmlow, \emph{Circuits, Signals, and
Systems for Bioengineers}, 3.ª ed., 2018.

§1.4.4, \textbf{Linear Versus Nonlinear Signals And Systems} (página 44
del visor PDF); §5.2, \textbf{Linear Systems Analysis---An Overview}
(página 212 del visor PDF).

\textbf{Correspondencia:} Las secciones aportan el fundamento. La
diapositiva amplía la formulación o precisa condiciones que no se
presentan con idéntica notación en el libro. La invariancia temporal se
comprueba separadamente de la linealidad.

\textbf{Microcurrículo:} semana 10, sesión 2 (sesión global 20). Los
numerales del cronograma son identificadores curriculares; no son los
apartados del libro citados arriba.}
\end{frame}

\begin{frame}{Memoria y causalidad son propiedades distintas}
\phantomsection\label{d2-152}
\begin{itemize}
\tightlist
\item
  sin memoria: \(y(t)\) depende solo de \(x(t)\);
\item
  con memoria: depende de otros instantes;
\item
  causal: no usa valores futuros de la entrada.
\end{itemize}

\begin{block}{Concepto clave}

Un promedio móvil causal tiene memoria. Un filtrado de fase cero fuera
de línea usa muestras futuras y no es causal.

\end{block}

\note{\textbf{Libro guía:} John Semmlow, \emph{Circuits, Signals, and
Systems for Bioengineers}, 3.ª ed., 2018.

§1.4.3, \textbf{Causal Versus Noncausal Signals And Systems} (página 43
del visor PDF); §8.4.3, \textbf{Causal And Noncausal Filters} (página
365 del visor PDF).

\textbf{Correspondencia:} El tema se desarrolla en las secciones
indicadas. La redacción es una síntesis docente.

\textbf{Microcurrículo:} semana 10, sesión 2 (sesión global 20). Los
numerales del cronograma son identificadores curriculares; no son los
apartados del libro citados arriba.}
\end{frame}

\begin{frame}{Estabilidad BIBO limita la salida}
\phantomsection\label{d2-153}
Si \(|x(t)|\leq M_x<\infty\), estabilidad BIBO exige
\(|y(t)|\leq M_y<\infty\).

Para un LTI continuo cuya respuesta \(h\) es una función ordinaria, el
criterio BIBO es

\[\int_{-\infty}^{\infty}|h(t)|dt<\infty.\]

\begin{alertblock}{Atención}

Un término directo \(K\delta(t)\) es estable y se trata como una medida
de masa finita, no con esta integral ordinaria. Un diferenciador ideal
no es BIBO estable.

\end{alertblock}

\note{\textbf{Libro guía:} John Semmlow, \emph{Circuits, Signals, and
Systems for Bioengineers}, 3.ª ed., 2018.

§5.4, \textbf{Using The Impulse Response To Find A Systems Output To Any
Input---Convolution} (página 219 del visor PDF); §7.5, \textbf{The
Laplace Domain, The Frequency Domain, And The Time Domain} (página 330
del visor PDF).

\textbf{Correspondencia:} Las secciones aportan el fundamento. La
diapositiva amplía la formulación o precisa condiciones que no se
presentan con idéntica notación en el libro. Se explicita el criterio
BIBO y el caso de término directo distribucional.

\textbf{Microcurrículo:} semana 10, sesión 2 (sesión global 20). Los
numerales del cronograma son identificadores curriculares; no son los
apartados del libro citados arriba.}
\end{frame}

\begin{frame}{El impulso ideal concentra área unitaria}
\phantomsection\label{d2-154}
\[\delta(t)=0\quad\text{para }t\neq0,\]

\[\int_{-\infty}^{\infty}\delta(t)dt=1.\]

Propiedad de selección:

\[\int x(\tau)\delta(t-\tau)d\tau=x(t).\]

\begin{block}{Nota}

El impulso de Dirac es una distribución matemática. Los experimentos
físicos usan pulsos breves con amplitud y energía limitadas.

\end{block}

\note{\textbf{Libro guía:} John Semmlow, \emph{Circuits, Signals, and
Systems for Bioengineers}, 3.ª ed., 2018.

§5.3, \textbf{A Slice In Time: The Impulse Signal} (página 214 del visor
PDF); §5.3.1, \textbf{Real-World Impulse Signals} (página 215 del visor
PDF).

\textbf{Correspondencia:} El tema se desarrolla en las secciones
indicadas. La redacción es una síntesis docente.

\textbf{Microcurrículo:} semana 10, sesión 2 (sesión global 20). Los
numerales del cronograma son identificadores curriculares; no son los
apartados del libro citados arriba.}
\end{frame}

\begin{frame}{En discreto, el impulso es una secuencia ordinaria}
\phantomsection\label{d2-155}
\[\delta[n]=\begin{cases}1,&n=0,\\0,&n\neq0.\end{cases}\]

Toda secuencia puede escribirse como

\[x[n]=\sum_{k=-\infty}^{\infty}x[k]\delta[n-k].\]

\begin{block}{Concepto clave}

La descomposición expresa la entrada como suma de impulsos desplazados y
ponderados.

\end{block}

\note{\textbf{Libro guía:} John Semmlow, \emph{Circuits, Signals, and
Systems for Bioengineers}, 3.ª ed., 2018.

§5.3, \textbf{A Slice In Time: The Impulse Signal} (página 214 del visor
PDF); §5.4, \textbf{Using The Impulse Response To Find A Systems Output
To Any Input---Convolution} (página 219 del visor PDF).

\textbf{Correspondencia:} El tema se desarrolla en las secciones
indicadas. La redacción es una síntesis docente.

\textbf{Microcurrículo:} semana 10, sesión 2 (sesión global 20). Los
numerales del cronograma son identificadores curriculares; no son los
apartados del libro citados arriba.}
\end{frame}

\begin{frame}{La respuesta al impulso caracteriza un sistema LTI}
\phantomsection\label{d2-156}
\[h(t)=\mathcal{T}\{\delta(t)\}.\]

Por linealidad e invariancia:

\[y(t)=\int_{-\infty}^{\infty}x(\tau)h(t-\tau)d\tau.\]

\begin{block}{Concepto clave}

La convolución representa la respuesta de estado cero. Si el sistema
almacena energía inicialmente, debe añadirse la respuesta de entrada
cero.

\end{block}

\note{\textbf{Libro guía:} John Semmlow, \emph{Circuits, Signals, and
Systems for Bioengineers}, 3.ª ed., 2018.

§5.4, \textbf{Using The Impulse Response To Find A Systems Output To Any
Input---Convolution} (página 219 del visor PDF).

\textbf{Correspondencia:} El tema se desarrolla en las secciones
indicadas. La redacción es una síntesis docente.

\textbf{Microcurrículo:} semana 10, sesión 2 (sesión global 20). Los
numerales del cronograma son identificadores curriculares; no son los
apartados del libro citados arriba.}
\end{frame}

\begin{frame}{La salida es una convolución}
\phantomsection\label{d2-157}
\[y(t)=x(t)*h(t).\]

En tiempo discreto:

\[y[n]=\sum_{k=-\infty}^{\infty}x[k]h[n-k].\]

\begin{block}{Nota}

La semana siguiente desarrolla interpretación gráfica, soporte y cálculo
de la convolución.

\end{block}

\note{\textbf{Libro guía:} John Semmlow, \emph{Circuits, Signals, and
Systems for Bioengineers}, 3.ª ed., 2018.

§5.4, \textbf{Using The Impulse Response To Find A Systems Output To Any
Input---Convolution} (página 219 del visor PDF).

\textbf{Correspondencia:} El tema se desarrolla en las secciones
indicadas. La redacción es una síntesis docente.

\textbf{Microcurrículo:} semana 10, sesión 2 (sesión global 20). Los
numerales del cronograma son identificadores curriculares; no son los
apartados del libro citados arriba.}
\end{frame}

\begin{frame}{Causalidad de un LTI se lee en \(h\)}
\phantomsection\label{d2-158}
Un sistema LTI causal cumple

\[h(t)=0\quad\text{para }t<0,\]

o en discreto \(h[n]=0\) para \(n<0\).

\begin{alertblock}{Atención}

Una respuesta al impulso simétrica alrededor de cero no puede
implementarse causalmente sin añadir retardo.

\end{alertblock}

\note{\textbf{Libro guía:} John Semmlow, \emph{Circuits, Signals, and
Systems for Bioengineers}, 3.ª ed., 2018.

§1.4.3, \textbf{Causal Versus Noncausal Signals And Systems} (página 43
del visor PDF); §5.4, \textbf{Using The Impulse Response To Find A
Systems Output To Any Input---Convolution} (página 219 del visor PDF).

\textbf{Correspondencia:} El tema se desarrolla en las secciones
indicadas. La redacción es una síntesis docente.

\textbf{Microcurrículo:} semana 10, sesión 2 (sesión global 20). Los
numerales del cronograma son identificadores curriculares; no son los
apartados del libro citados arriba.}
\end{frame}

\begin{frame}{Las exponenciales complejas son autofunciones LTI}
\phantomsection\label{d2-159}
Si \(x(t)=e^{j2\pi ft}\), entonces

\[y(t)=H(f)e^{j2\pi ft}.\]

\begin{block}{Concepto clave}

La relación exige que \(H(f)\) exista en esa frecuencia y que la
convolución converja. Para excitación iniciada en tiempo finito,
distinga transitorio y régimen permanente.

\end{block}

\note{\textbf{Libro guía:} John Semmlow, \emph{Circuits, Signals, and
Systems for Bioengineers}, 3.ª ed., 2018.

§6.3, \textbf{The Response Of System Elements To Sinusoidal Inputs:
Phasor Analysis} (página 251 del visor PDF).

\textbf{Correspondencia:} El tema se desarrolla en las secciones
indicadas. La redacción es una síntesis docente.

\textbf{Microcurrículo:} semana 10, sesión 2 (sesión global 20). Los
numerales del cronograma son identificadores curriculares; no son los
apartados del libro citados arriba.}
\end{frame}

\begin{frame}{\(H(f)\) es la transformada de \(h(t)\)}
\phantomsection\label{d2-160}
\[H(f)=\int_{-\infty}^{\infty}h(t)e^{-j2\pi ft}dt.\]

La relación

\[Y(f)=X(f)H(f)\]

coincide con el teorema de convolución estudiado en la semana 6.

\begin{block}{Comprobación}

Tiempo y frecuencia describen el mismo sistema mediante \(h(t)\) y
\(H(f)\).

\end{block}

\note{\textbf{Libro guía:} John Semmlow, \emph{Circuits, Signals, and
Systems for Bioengineers}, 3.ª ed., 2018.

§5.6, \textbf{Convolution In The Frequency Domain} (página 237 del visor
PDF); §6.7, \textbf{The Transfer Function And The Fourier Transform}
(página 287 del visor PDF).

\textbf{Correspondencia:} El tema se desarrolla en las secciones
indicadas. La redacción es una síntesis docente.

\textbf{Microcurrículo:} semana 10, sesión 2 (sesión global 20). Los
numerales del cronograma son identificadores curriculares; no son los
apartados del libro citados arriba.}
\end{frame}

\begin{frame}{Ejemplos de sistemas de procesamiento}
\phantomsection\label{d2-161}
\begin{longtable}[]{@{}lll@{}}
\toprule\noalign{}
Sistema & Respuesta temporal esperada & Acción frecuencial \\
\midrule\noalign{}
\endhead
promedio móvil & suaviza cambios rápidos & atenúa frecuencias altas \\
diferenciador & resalta pendientes & amplifica frecuencias altas \\
retardo & desplaza la señal & añade fase lineal \\
ganancia & escala amplitud & multiplica por constante \\
\bottomrule\noalign{}
\end{longtable}

\begin{alertblock}{Atención}

La acción útil depende de preservar los eventos y medidas relevantes
para la tarea biomédica.

\end{alertblock}

\note{\textbf{Libro guía:} John Semmlow, \emph{Circuits, Signals, and
Systems for Bioengineers}, 3.ª ed., 2018.

§5.5, \textbf{Applied Convolution---Basic Filters} (página 231 del visor
PDF); §6.5, \textbf{The Spectrum Of System Elements: The Bode Plot}
(página 263 del visor PDF).

\textbf{Correspondencia:} El tema se desarrolla en las secciones
indicadas. La redacción es una síntesis docente.

\textbf{Microcurrículo:} semana 10, sesión 2 (sesión global 20). Los
numerales del cronograma son identificadores curriculares; no son los
apartados del libro citados arriba.}
\end{frame}

\begin{frame}[fragile]{Código: probar linealidad numéricamente}
\phantomsection\label{d2-162}
\begin{Shaded}
\begin{Highlighting}[]
\ImportTok{import}\NormalTok{ numpy }\ImportTok{as}\NormalTok{ np}

\NormalTok{y\_combo }\OperatorTok{=}\NormalTok{ T(a}\OperatorTok{*}\NormalTok{x1 }\OperatorTok{+}\NormalTok{ b}\OperatorTok{*}\NormalTok{x2)}
\NormalTok{y\_super }\OperatorTok{=}\NormalTok{ a}\OperatorTok{*}\NormalTok{T(x1) }\OperatorTok{+}\NormalTok{ b}\OperatorTok{*}\NormalTok{T(x2)}
\NormalTok{error }\OperatorTok{=}\NormalTok{ np.linalg.norm(y\_combo }\OperatorTok{{-}}\NormalTok{ y\_super)}
\NormalTok{referencia }\OperatorTok{=}\NormalTok{ np.linalg.norm(y\_combo)}
\NormalTok{error\_rel }\OperatorTok{=}\NormalTok{ error }\OperatorTok{/} \BuiltInTok{max}\NormalTok{(referencia, }\FloatTok{1e{-}12}\NormalTok{)}
\end{Highlighting}
\end{Shaded}

\begin{alertblock}{Atención}

Un error pequeño para algunos vectores apoya la hipótesis, pero no
demuestra linealidad para todas las entradas.

\end{alertblock}

\note{\textbf{Libro guía:} John Semmlow, \emph{Circuits, Signals, and
Systems for Bioengineers}, 3.ª ed., 2018.

§5.2.1, \textbf{Superposition And Linearity} (página 213 del visor PDF).

\textbf{Correspondencia:} Ejemplo, figura o implementación de
elaboración docente a partir de estos fundamentos. No es una
transcripción del código MATLAB del libro.

\textbf{Microcurrículo:} semana 10, sesión 2 (sesión global 20). Los
numerales del cronograma son identificadores curriculares; no son los
apartados del libro citados arriba.}
\end{frame}

\begin{frame}{Verificación de la sesión 20}
\phantomsection\label{d2-163}
Para \(y[n]=0.5x[n]+0.5x[n-1]\):

\begin{enumerate}
\tightlist
\item
  verifique linealidad;
\item
  determine memoria y causalidad;
\item
  obtenga \(h[n]\);
\item
  evalúe estabilidad BIBO;
\item
  interprete su efecto sobre un ECG;
\item
  anticipe su respuesta frecuencial.
\end{enumerate}

\begin{block}{Criterio de logro}

La respuesta debe conectar la ecuación, las propiedades, \(h[n]\) y la
interpretación biomédica.

\end{block}

\note{\textbf{Libro guía:} John Semmlow, \emph{Circuits, Signals, and
Systems for Bioengineers}, 3.ª ed., 2018.

§5.4, \textbf{Using The Impulse Response To Find A Systems Output To Any
Input---Convolution} (página 219 del visor PDF); §5.5, \textbf{Applied
Convolution---Basic Filters} (página 231 del visor PDF).

\textbf{Correspondencia:} Actividad o interpretación biomédica de
elaboración docente apoyada en las secciones indicadas.

\textbf{Microcurrículo:} semana 10, sesión 2 (sesión global 20). Los
numerales del cronograma son identificadores curriculares; no son los
apartados del libro citados arriba.}
\end{frame}

\begin{frame}{Cierre de la semana 10}
\phantomsection\label{d2-164}
\begin{itemize}
\tightlist
\item
  La selección de ventana define una escala tiempo-frecuencia.
\item
  Ninguna configuración maximiza ambas resoluciones.
\item
  Linealidad e invariancia temporal permiten usar superposición.
\item
  La respuesta al impulso caracteriza un sistema LTI.
\item
  \(h(t)\) y \(H(f)\) son descripciones complementarias.
\end{itemize}

\begin{block}{Pregunta de salida}

¿Qué evidencia experimental necesitaría para tratar un sistema
fisiológico como LTI durante un intervalo específico?

\end{block}

\note{\textbf{Libro guía:} John Semmlow, \emph{Circuits, Signals, and
Systems for Bioengineers}, 3.ª ed., 2018.

§4.6, \textbf{Time--Frequency Analysis} (página 205 del visor PDF);
§5.4, \textbf{Using The Impulse Response To Find A Systems Output To Any
Input---Convolution} (página 219 del visor PDF).

\textbf{Correspondencia:} Síntesis o encuadre que reúne varias
secciones; no corresponde a un único apartado del libro.

\textbf{Microcurrículo:} semana 10, sesión 2 (sesión global 20). Los
numerales del cronograma son identificadores curriculares; no son los
apartados del libro citados arriba.}
\end{frame}

\begin{frame}{Síntesis y estudio autónomo}
\phantomsection\label{d2-165}
\note{\textbf{Libro guía:} John Semmlow, \emph{Circuits, Signals, and
Systems for Bioengineers}, 3.ª ed., 2018.

§3.4, \textbf{Time--Frequency Transformation In The Discrete Domain}
(página 142 del visor PDF); §4.2, \textbf{Data Acquisition And Storage}
(página 175 del visor PDF); §4.6, \textbf{Time--Frequency Analysis}
(página 205 del visor PDF); §5.4, \textbf{Using The Impulse Response To
Find A Systems Output To Any Input---Convolution} (página 219 del visor
PDF).

\textbf{Correspondencia:} Síntesis o encuadre que reúne varias
secciones; no corresponde a un único apartado del libro.

\textbf{Lectura:} use estas secciones como ruta de preparación y
consulta. La bibliografía aparece al final del bloque.}
\end{frame}

\begin{frame}{La cadena completa de decisiones}
\phantomsection\label{d2-166}
\[\begin{aligned}
\text{pregunta}&\rightarrow\text{adquisición}\rightarrow\text{muestreo}\rightarrow\text{segmento}\\
&\rightarrow\text{ventana}\rightarrow\text{estimador}\rightarrow\text{validación}\rightarrow\text{inferencia}
\end{aligned}\]

\begin{block}{Concepto clave}

Cada resultado espectral debe conservar trazabilidad hacia el registro,
los parámetros y el estado fisiológico observado.

\end{block}

\note{\textbf{Libro guía:} John Semmlow, \emph{Circuits, Signals, and
Systems for Bioengineers}, 3.ª ed., 2018.

§3.4, \textbf{Time--Frequency Transformation In The Discrete Domain}
(página 142 del visor PDF); §4.2, \textbf{Data Acquisition And Storage}
(página 175 del visor PDF); §4.6, \textbf{Time--Frequency Analysis}
(página 205 del visor PDF); §5.4, \textbf{Using The Impulse Response To
Find A Systems Output To Any Input---Convolution} (página 219 del visor
PDF).

\textbf{Correspondencia:} Síntesis o encuadre que reúne varias
secciones; no corresponde a un único apartado del libro.

\textbf{Lectura:} use estas secciones como ruta de preparación y
consulta. La bibliografía aparece al final del bloque.}
\end{frame}

\begin{frame}{Errores que deben poder diagnosticarse}
\phantomsection\label{d2-167}
\begin{enumerate}
\tightlist
\item
  sumar magnitudes cuando corresponde sumar espectros complejos;
\item
  interpretar zero padding como mayor resolución real;
\item
  atribuir fuga a nuevas oscilaciones fisiológicas;
\item
  usar Nyquist sin margen para el filtro analógico;
\item
  confundir potencia por bin con PSD;
\item
  promediar estados incompatibles con Welch;
\item
  leer resolución desde el número de columnas del espectrograma;
\item
  asumir que un sistema biológico es LTI sin verificar rango y estado.
\end{enumerate}

\note{\textbf{Libro guía:} John Semmlow, \emph{Circuits, Signals, and
Systems for Bioengineers}, 3.ª ed., 2018.

§3.4, \textbf{Time--Frequency Transformation In The Discrete Domain}
(página 142 del visor PDF); §4.2, \textbf{Data Acquisition And Storage}
(página 175 del visor PDF); §4.6, \textbf{Time--Frequency Analysis}
(página 205 del visor PDF); §5.4, \textbf{Using The Impulse Response To
Find A Systems Output To Any Input---Convolution} (página 219 del visor
PDF).

\textbf{Correspondencia:} Síntesis o encuadre que reúne varias
secciones; no corresponde a un único apartado del libro.

\textbf{Lectura:} use estas secciones como ruta de preparación y
consulta. La bibliografía aparece al final del bloque.}
\end{frame}

\begin{frame}{Autoevaluación acumulativa}
\phantomsection\label{d2-168}
Puede explicar, sin consultar notas:

\begin{itemize}
\tightlist
\item
  cómo un retraso modifica la fase;
\item
  por qué \(\Delta f=1/T\);
\item
  qué cambia al aplicar Hann;
\item
  cómo se produce aliasing;
\item
  qué unidades tiene una PSD;
\item
  qué intercambia Welch entre resolución y varianza;
\item
  cómo se selecciona una ventana STFT;
\item
  por qué \(h(t)\) caracteriza un sistema LTI.
\end{itemize}

\begin{block}{Nota}

Una explicación completa incluye ecuación, unidades, supuesto, efecto
observable y límite de interpretación.

\end{block}

\note{\textbf{Libro guía:} John Semmlow, \emph{Circuits, Signals, and
Systems for Bioengineers}, 3.ª ed., 2018.

§3.4, \textbf{Time--Frequency Transformation In The Discrete Domain}
(página 142 del visor PDF); §4.2, \textbf{Data Acquisition And Storage}
(página 175 del visor PDF); §4.6, \textbf{Time--Frequency Analysis}
(página 205 del visor PDF); §5.4, \textbf{Using The Impulse Response To
Find A Systems Output To Any Input---Convolution} (página 219 del visor
PDF).

\textbf{Correspondencia:} Síntesis o encuadre que reúne varias
secciones; no corresponde a un único apartado del libro.

\textbf{Lectura:} use estas secciones como ruta de preparación y
consulta. La bibliografía aparece al final del bloque.}
\end{frame}

\begin{frame}{Actividad integradora sugerida}
\phantomsection\label{d2-169}
Con un registro biomédico real o simulado:

\begin{enumerate}
\tightlist
\item
  audite muestreo, unidades y duración;
\item
  compare periodograma y Welch;
\item
  analice dos ventanas;
\item
  verifique potencia mediante integración;
\item
  construya una STFT con dos longitudes;
\item
  interprete cambios y artefactos;
\item
  documente parámetros en Python;
\item
  formule una conclusión limitada por la evidencia.
\end{enumerate}

\begin{block}{Producto esperado}

Un cuaderno reproducible que permita reconstruir cada figura y
justificar cada decisión.

\end{block}

\note{\textbf{Libro guía:} John Semmlow, \emph{Circuits, Signals, and
Systems for Bioengineers}, 3.ª ed., 2018.

§3.4, \textbf{Time--Frequency Transformation In The Discrete Domain}
(página 142 del visor PDF); §4.2, \textbf{Data Acquisition And Storage}
(página 175 del visor PDF); §4.6, \textbf{Time--Frequency Analysis}
(página 205 del visor PDF); §5.4, \textbf{Using The Impulse Response To
Find A Systems Output To Any Input---Convolution} (página 219 del visor
PDF).

\textbf{Correspondencia:} Síntesis o encuadre que reúne varias
secciones; no corresponde a un único apartado del libro.

\textbf{Lectura:} use estas secciones como ruta de preparación y
consulta. La bibliografía aparece al final del bloque.}
\end{frame}

\begin{frame}{Referencias y alcance}
\phantomsection\label{d2-170}
La programación de semanas 6--10: Fourier discreto, adquisición,
estimación espectral y sistemas LTI se contrasta con el microcurrículo
suministrado. El texto principal es \textbf{Semmlow, 3.ª edición (2018)}
(\citeproc{ref-semmlow2018}{J. Semmlow 2018}).

Las notas distinguen correspondencia directa, adaptación computacional y
ampliación. Los números del cronograma son identificadores curriculares
y no se usan como secciones del libro.

Complementos: Semmlow y Griffel (2014), para wavelets y clasificación
(\citeproc{ref-semmlowgriffel2014}{J. L. Semmlow y Griffel 2014}), y
Esakkirajan, Veerakumar y Subudhi (2024), para DSP con Python
(\citeproc{ref-esakkirajan2024}{Esakkirajan, Veerakumar, y Subudhi
2024}).

\note{\textbf{Libro guía:} John Semmlow, \emph{Circuits, Signals, and
Systems for Bioengineers}, 3.ª ed., 2018.

§3.4, \textbf{Time--Frequency Transformation In The Discrete Domain}
(página 142 del visor PDF); §4.2, \textbf{Data Acquisition And Storage}
(página 175 del visor PDF); §4.6, \textbf{Time--Frequency Analysis}
(página 205 del visor PDF); §5.4, \textbf{Using The Impulse Response To
Find A Systems Output To Any Input---Convolution} (página 219 del visor
PDF).

\textbf{Correspondencia:} Síntesis o encuadre que reúne varias
secciones; no corresponde a un único apartado del libro.

\textbf{Lectura:} use estas secciones como ruta de preparación y
consulta. La bibliografía aparece al final del bloque.}
\end{frame}

\begin{frame}{Cierre: criterio de dominio}
\phantomsection\label{d2-171}
El estudiante debe sostener esta afirmación con evidencia:

\begin{quote}
Un espectro no es una propiedad aislada del organismo. Es el resultado
de observar, segmentar y estimar una señal bajo decisiones explícitas.
\end{quote}

\begin{block}{Concepto clave}

El dominio del tema se demuestra cuando una conclusión puede rastrearse
hasta la adquisición, las unidades, los parámetros y los supuestos.

\end{block}

\note{\textbf{Libro guía:} John Semmlow, \emph{Circuits, Signals, and
Systems for Bioengineers}, 3.ª ed., 2018.

§3.4, \textbf{Time--Frequency Transformation In The Discrete Domain}
(página 142 del visor PDF); §4.2, \textbf{Data Acquisition And Storage}
(página 175 del visor PDF); §4.6, \textbf{Time--Frequency Analysis}
(página 205 del visor PDF); §5.4, \textbf{Using The Impulse Response To
Find A Systems Output To Any Input---Convolution} (página 219 del visor
PDF).

\textbf{Correspondencia:} Síntesis o encuadre que reúne varias
secciones; no corresponde a un único apartado del libro.

\textbf{Lectura:} use estas secciones como ruta de preparación y
consulta. La bibliografía aparece al final del bloque.}
\end{frame}

\begin{frame}[allowframebreaks]{Referencias}
\phantomsection\label{d2-172}
\phantomsection\label{refs}
\begin{CSLReferences}{1}{0}
\bibitem[\citeproctext]{ref-esakkirajan2024}
Esakkirajan, S., T. Veerakumar, y Badri N. Subudhi. 2024. \emph{Digital
Signal Processing: Illustration Using Python}. Springer.
\url{https://doi.org/10.1007/978-981-99-6752-0}.

\bibitem[\citeproctext]{ref-semmlow2018}
Semmlow, John. 2018. \emph{Circuits, Signals, and Systems for
Bioengineers: A MATLAB-Based Introduction}. 3.ª ed. Academic Press.

\bibitem[\citeproctext]{ref-semmlowgriffel2014}
Semmlow, John L., y Benjamin Griffel. 2014. \emph{Biosignal and Medical
Image Processing}. 3.ª ed. CRC Press.

\end{CSLReferences}

\note{\textbf{Libro guía:} John Semmlow, \emph{Circuits, Signals, and
Systems for Bioengineers}, 3.ª ed., 2018.

§3.4, \textbf{Time--Frequency Transformation In The Discrete Domain}
(página 142 del visor PDF); §4.2, \textbf{Data Acquisition And Storage}
(página 175 del visor PDF); §4.6, \textbf{Time--Frequency Analysis}
(página 205 del visor PDF); §5.4, \textbf{Using The Impulse Response To
Find A Systems Output To Any Input---Convolution} (página 219 del visor
PDF).

\textbf{Correspondencia:} Síntesis o encuadre que reúne varias
secciones; no corresponde a un único apartado del libro.

\textbf{Lectura:} use estas secciones como ruta de preparación y
consulta. La bibliografía aparece al final del bloque.}
\end{frame}




\end{document}
